\documentclass[a4paper,11pt,oneside,fleqn]{book}


\usepackage[T1]{fontenc}
\usepackage[utf8]{inputenc}
\usepackage{graphicx}
\usepackage{xcolor}
\usepackage[colorlinks = true, linkcolor = blue]{hyperref}

%algorithm
\usepackage{algpseudocode}

\usepackage{amsmath,amssymb,amsthm,textcomp}
\usepackage{multicol}
\usepackage{tikz}
\usepackage[normalem]{ulem}
\usepackage{bbm} % using for the characteristic function
\usepackage{ulem}
\usepackage{cancel}
\usepackage{esint} %integral mean

\usepackage{geometry}
\geometry{left=10mm,right=10mm,%
bindingoffset=0mm, top=10mm,bottom=10mm}



\renewcommand{\thechapter}{\Roman{chapter}}
\renewcommand*\thesection{\arabic{section}}


%\newtheorem{theorem}{Theorem}
\newtheorem{theorem}{Theorem}[section]
\newtheorem{algorithm}{Algorithm}[section]
\renewcommand{\thetheorem}{\arabic{section}.\arabic{theorem}}
\newtheorem{corollary}{Corollary}[theorem]
\newtheorem{lemma}[theorem]{Lemma}
\newtheorem{proposition}[theorem]{Proposition}

\theoremstyle{definition}
\newtheorem{definition}{Definition}[section]
\renewcommand{\thedefinition}{\arabic{section}.\arabic{definition}}

\newtheorem*{example}{Example}
\newtheorem{assumption}{Assumption}[section]
\newtheorem{exercise}{Exercise}[chapter]
\theoremstyle{remark}
\newtheorem*{remark}{Remark}
\newtheorem{review}{Review}

%contradiction
\usepackage{marvosym}

\linespread{1.3}
\newcommand{\linia}{\rule{\linewidth}{0.5pt}}


% my own titles
\makeatletter
\renewcommand{\maketitle}{
\begin{center}
\vspace{2ex}
{\huge \textsc{\@title}}
\vspace{1ex}
\\
\linia\\
\@author \hfill \@date
\vspace{4ex}
\end{center}
}
\makeatother
%%%

% custom footers and headers
\usepackage{fancyhdr}
\pagestyle{fancy}
\lhead{}
\chead{}
\rhead{}
\lfoot{}
\cfoot{}
\rfoot{\thepage}
\renewcommand{\headrulewidth}{0pt}
\renewcommand{\footrulewidth}{0pt}
%

% code listing settings
\usepackage{listings}
\lstset{
    language=Python,
    basicstyle=\ttfamily\small,
    aboveskip={1.0\baselineskip},
    belowskip={1.0\baselineskip},
    columns=fixed,
    extendedchars=true,
    breaklines=true,
    tabsize=4,
    prebreak=\raisebox{0ex}[0ex][0ex]{\ensuremath{\hookleftarrow}},
    frame=lines,
    showtabs=false,
    showspaces=false,
    showstringspaces=false,
    keywordstyle=\color[rgb]{0.627,0.126,0.941},
    commentstyle=\color[rgb]{0.133,0.545,0.133},
    stringstyle=\color[rgb]{01,0,0},
    numbers=left,
    numberstyle=\small,
    stepnumber=1,
    numbersep=10pt,
    captionpos=t,
    escapeinside={\%*}{*)}
}
\DeclareMathAlphabet{\mathcal}{OMS}{cmsy}{m}{n}

%symbol
\AtBeginDocument{
  \let\ran\relax
  \DeclareMathOperator{\dom}{dom}
}
\DeclareMathOperator{\var}{Var}

\AtBeginDocument{
  \let\div\relax
  \DeclareMathOperator{\div}{div}
}

\AtBeginDocument{
  \let\ran\relax
  \DeclareMathOperator{\ran}{ran}
}
\DeclareMathOperator*{\esssup}{ess\, sup}
\DeclareMathOperator*{\spa}{span}
\DeclareMathOperator*{\dist}{dist}
\DeclareMathOperator*{\support}{supp}


\begin{document}
\title{Analysis}
\maketitle
\tableofcontents



\chapter{Measure Theory}
\section{Measure and measurable sets}

\subsection{Set theory basics}
For the development of measure theory, Borel has introduced the following defintion:
\begin{definition}[Limes superior and limes inferior]
	Given $(A_n)_{n \geq 1}$ a sequence of subsets in $X$. We define:
	\begin{align*}
		\limsup_n A_n &:= \{ x \in X: x \in A_n \textbf{ for infinitely many } A_n \}\\
		\liminf_n A_n &:= \{ x \in X: \exists n_0(x) \in \mathbb{N}\forall n \geq n_0(x): x \in A_n \} \textbf{ x is in all but finitely many $A_n$}
	\end{align*}
	equivalently we see:
	\begin{align*}
		\limsup_n A_n &= \cap_{n = 1}^\infty \cup_{k = n}^\infty A_n\\ 
		\liminf_n A_n &= \cup_{n = 1}^\infty \cap_{k = n}^\infty A_n
	\end{align*}
	We define the convergence of sets as :
	\begin{align*}
		\limsup_n A_n = \liminf_n A_n := \lim_n A_n
	\end{align*}
\end{definition}
Note that the definition of $\liminf$ coincides with the definition of limes we are familiar with. However for a sequence of sets it always exists. The limes of a sequence of sets could be interpreted as: all elements that appear infinitely many times is in $A_n$ for $n$ large enough.


\subsection{($\sigma)$- Ring and $\sigma$-Algebra}
\begin{theorem}
	Given a power set of $X$ $\mathfrak{B}(X)$, $(\mathfrak{B}(X), \Delta, \cap)$ is a commutative ring with the Null $\emptyset$ and the identity $X$.
\end{theorem}
\begin{proof}
	There is an isomorphism between the power set $\mathfrak{B}(X)$ and the ring of characteristic function(for each $A \subset X$ we can define a homomorphism $\chi_A: X \rightarrow \{\bar{1}, \bar{0} \}$ to the ring $\{\bar{0}, \bar{1}\}$). Obviously it's a isomorphism.
\end{proof}

The algebraic structure works for the finite case. For developing the integration theory, finite combination is not enough, therefore we introduce the concept of $\sigma$- ring and $\sigma$- field.

\begin{definition}[$\sigma$-Ring and Field]
Given a set $X$
	\begin{itemize}
		\item Ring and field: A subring of the ring defined above is a \textbf{ring }over $X$. Such ring is an \textbf{Field} if it contains the $X$.  
		\item $\sigma$-Ring  and $\sigma$-field: A ring is a \textbf{$\sigma$-ring} if it's closed under \textbf{countable union}
	\end{itemize}
\end{definition}

\begin{theorem}
	$\mathfrak{R} \subset \mathfrak B(X)$ is a ring iff one of the following is true:
	\begin{enumerate}
		\item $\emptyset \in \mathfrak{R}, \forall A, B \in \mathfrak{R}: A \Delta B \in \mathfrak{R}, A\bigcap B \in \mathfrak{R}$
		\item $\dots\Delta\dots \bigcup \dots$
		\item $\dots\backslash\dots \bigcup \dots$
	\end{enumerate}
\end{theorem}


\subsection{Generator of the Borel-sets}
In the following theorem, we introduced an useful technique called \textit{Principle of Good sets}. It's often applied when we want to prove that a generator produces the same set under 2 different sets of operations.
\begin{theorem}
	$f: X \rightarrow Y$ a measurable map, $\mathfrak{G} \subset \mathfrak{B}(X)$ a generator of $\mathfrak{B}$. we have:
	\begin{align*}
		\sigma(f^{-1}(\mathfrak{G}) ) = f^{-1}(\sigma (\mathfrak{G}))
	\end{align*}
\end{theorem} 
\begin{proof}
\end{proof}

\subsection{Semiring}
In practice it's often easy to generate the $\sigma$-ring or $\sigma$-field with generators of simple structure. A class of this type is what we'll introduce: the \textbf{Semi-ring}. This concept was first introduced by Johann von Neumann.
\begin{definition}[Semi-ring]
 A set of sets $\mathfrak{H}$ is called semi-ring if it satisfies
 \begin{enumerate}
 	\item $\emptyset \in \mathfrak{H}$
 	\item $\forall A, B \in \mathfrak{H}: A \cap B \in \mathfrak{H}$
 	\item $\forall A, B \in \mathfrak{H}$ there exists disjoint $C_1, C_2, \dots C_n: A\slash B = \bigcup_{k = 1}^n C_k$
 \end{enumerate}
\end{definition}
$\quad$\\
By \textbf{H.Hahn} we know semi-ring can generate a ring:
\begin{theorem}
	Is $\mathfrak{H} \subset \mathfrak{B}(X)$ a semiring, then the disjoint union of the element in $\mathfrak{H}$ generates a ring:
	\begin{align*}
		\mathfrak{R} := \{ \bigcup_{I < \infty} A_i,\quad \textbf{$A_i\in \mathfrak{H}$ disjoint } \} 
	\end{align*}
\end{theorem}
\begin{proof}
\end{proof}


\subsection{Monotone class and Dynkin System}
In the last 2 subsections, we have introduced the ring-structure of the set of sets. Also, given a semi-ring, we can generate a ring with the operation of disjoint union. Since the main object of study is the $\sigma$-field, we need to find the operation, under which a ring becomes a $\sigma$-field. In this spirit, we'll introduce the concept of \textbf{Dynkin system}, the main result is the following theorem:

	"Given an intersection-stable generator $\mathfrak{G} \subset \mathfrak{B}(X)$, the Dynkin-system generated by $\mathfrak{G}$ equals to the $\sigma$-field generated by $\mathfrak{G}$."
	 
\subsubsection*{The monotone class}
\begin{definition}[Monotone Class]
	A system is a monotone class if it is closed under countable union, i.e. $\mathfrak{M} \subset \mathfrak{B}(X)$
	\begin{align*}
		\forall (A_i)_{i \in I} \in \mathfrak{M} \textbf{ monotone}: \lim_{i\rightarrow \infty} A_i \in \mathfrak{M} \}
	\end{align*}
\end{definition}
The importance of the monotone class lies in the fact that it's the necessary and sufficient condition for the underlying ring to generate a $\sigma$-ring.

\begin{theorem}
	The $\sigma$-ring generated by the $\mathfrak{R} \subset \mathfrak{B}(X)$ equals to the monotone class generated by $\mathfrak{R}$.
\end{theorem}
\begin{proof}
	The claim is $\mathfrak{M}(\mathfrak{R}) = \sigma(\mathfrak{R})$\\
	"$\subset$": trivial \\
	"$\supset$": again we apply \textit{Principle of the good sets}. Given $\mathfrak{M}(\mathfrak{R}):= \mathfrak{M}$. We define the sets in $X$ that are "good" under the ring operation(as $\mathfrak{D}$), and show that $\mathfrak{M} \subset \mathfrak{D} $.\par 
	Given $A \in \mathfrak{M}$, we define:
	\begin{align*}
		\mathfrak{D}(A):= \{ B \subset X: A \backslash B \in \mathfrak{M}, B \backslash A \in \mathfrak{M}, B \bigcup A \in \mathfrak{M} \}
	\end{align*}
	
It's easy to verify that $\mathfrak{D}(A)$ is a monotone class, since the ring operation upon $A$ would enclose all the sets in $\mathfrak{R}$, and $\mathfrak{M}$ is the smallest monotone class generated by $\mathfrak{R}$, it follows $\mathfrak{R}  \subset \mathfrak{M} \subset \mathfrak{D}(A)$.\par
Which means for any $A \in \mathfrak{M}$, to get a set in $\mathfrak{M}$ after performing the ring operation with a set $B$ in $X$, we can select $B$ as any set in $\mathfrak{M}$. This means $\mathfrak{M}$ is closed under the ring operation. 
\end{proof}
\begin{remark}[principle of the good sets]
\end{remark}


\begin{definition}[Dynkin system($\lambda$ System)]
$\mathfrak{D}\subset \mathfrak{B}(X)$ is a Dynkin system if
\begin{enumerate}
	\item $X \in \mathfrak{D}$
	\item $\forall A \in \mathfrak{D}: A^c \in \mathfrak{D}$
	\item $\forall (A_i)_{i \geq 1} \in \mathfrak{D} \textbf{ disjoint }: \bigcup_{i = 1}^\infty A_i \in \mathfrak{D}$
\end{enumerate} 
An equivalent definition is given by replacing 2. with "$\backslash$" and 3. with \textbf{monotone}
\end{definition}
\begin{remark}
The first definition corresponds the definition of measure. The $\sigma$-additivity of the measure is perfectly mirrored by the properties of the Dynkin-system.\par 
The Dykin-system differs from a $\sigma$-field in that it's closed under countable union of \textbf{disjoint} or \textbf{monotone} sets, where the $\sigma$-field of arbitrary sets.\par 
As mentioned at the beginning, the following most important theorem of this subchapter provides a tool to prove the uniqueness of the measure:
\end{remark}

\begin{lemma}
	A Dykin-system is a $\sigma$-field if it's intersection-stable.
\end{lemma}
\begin{proof}
	From sequence of arbitrary sets $A_n$ we may construct a sequence of monotone sets by:
	\begin{align*}
		B_1 &= A_1\\
		B_n: &= A_n \backslash B_{n-1}
	\end{align*}
	Then we have
	\begin{align*}
		\bigcup A_n = \bigcup B_n = \bigcup(A_n \cap B_{n-1}^C)
	\end{align*}
	If the Dykin-system is stable under intersection then we know it's stable under arbitrary union
\end{proof}

\begin{theorem}
	Given a intersection-stable generator $\mathfrak{G}$, the Dykin-system and $\sigma$-field generated by the generator is the same.
\end{theorem}







\subsection{Content, premeasure, and measure}
In the last section we dealt with the problem of the domain of a measure. From now on we shift our topic to the measure function. We follow the induction process of semi-ring, $\sigma$-ring and $\sigma$-algebra to introduce the measure. Axiomatically, a measure would satisfies positiveness and additivity.
\begin{definition}[Content, premeasure and measure]
Given $X$ a set, we define:
	\begin{itemize}
		\item \textbf{Content}: defined on a semiring: $\mathcal{Z} \rightarrow \mathbb{R}^+$ satisfies:
		\begin{enumerate}
			\item $\mu(\emptyset) = 0$. 
			\item $\mu > 0$
			\item \textbf{Finite additivity:} $A_i$ disjoint: $\mu(\bigcup_{i = 1}^n A_i) = \sum_{i = 1}^n \mu(A_i)$
		\end{enumerate}
		\item \textbf{premeasure}: function with properties 1), 2) whereas in 3) we would have $\sigma$-additivity.
		\item \textbf{measure}: defined on a $\sigma$-field.
	\end{itemize}
\end{definition}


\subsection{Extension from premeasure to measure}
Now we have a premeasure defined on a semi-ring. The tools at hand are 1) every set in the $\sigma$-field generated by $\mathfrak{H}$ could be approximated by the union of sets in $\mathfrak{H}$, 2) The $\sigma$-additivity, i.e.:
\begin{align*}
	\forall A \in \sigma(\mathfrak{R}): A = \bigcup_{n =1}^\infty A_n, \quad  A_n \in \mathfrak{R} \textbf{ disjoint }
\end{align*}
therefore
\begin{align*}
	\mu(A) = \sum_{n = 1}^\infty \mu(A_n)
\end{align*}
The problem is, the choice of $A_n$ is not unique. For an arbitrary $A \in \sigma(\mathfrak{R})$, we need to find a well-defined measure function.\par 
There are 2 ways to prove it's well-defined.
\begin{enumerate}
	\item \textbf{Young}
	\item \textbf{Caratheodory}
\end{enumerate} 

\begin{definition}[outer-measure]
	A outer measure $\eta: \mathfrak{B}(X) \rightarrow \bar{\mathbb{R}}$ is a map with the following properties:
	\begin{enumerate}
		\item $\eta(\emptyset) = 0$
		\item For $A \subset B$: $\eta(A) \leq \eta(B)$
		\item For a sequence $(A_n)_{n \geq 1}$: $\eta(\bigcup_n A_n) \leq \sum_{n = 1}^\infty \eta(A_n)$
	\end{enumerate}
\end{definition}

\begin{definition}
	Given a outer measure $\eta$, a set $A \subset X$ is $\eta$-measurable if:
	\begin{align*}
		\forall Q \subset X: \eta(Q) = \eta(Q \cap A) + \eta(Q \cap A^c)
	\end{align*}
\end{definition}

The convenience of Caratheodory is partially because the outer-measure defines a canonical $\sigma$-field. We can further prove the existence of an unique measure on the $\sigma$-field.
\begin{theorem}[Caratheodory]
	Given a outer measure $\eta: \mathfrak{B}(X) \rightarrow \mathbb{R}$:
	\begin{align*}
		\mathfrak{A}:= \{ A \in \mathfrak{B}(X): A \quad \eta-measurable \}
	\end{align*}
	defines a $\sigma$ algebra and $\eta|_{\mathfrak{A}}$ is a measure.
\end{theorem}

\begin{theorem}[The 1. extension theorem]
	Given $\mu: \mathfrak{R} \rightarrow \widehat{\mathbb{R}}$ a content and for $A$ define:
	\begin{align*}
		\eta(A):= \inf\{\sum_{n = 1}^\infty \mu(A_n): A_n \in \mathfrak{R}, \cup A_n \subset A \}
	\end{align*}
	Then $\eta$ is an outer measure.
\end{theorem}

\subsubsection*{Uniqueness of the Extension}
one of the central result is the following theorem, which is often used to prove the uniqueness of a random variable.
\begin{theorem}[Uniqueness of the extension]
	Given measure $\mu, \nu$ on the measurable space $(X, \mathfrak{A})$. $\mathfrak{E}$ a \textbf{intersection-stable }generator that satisfies:
	\begin{enumerate}
		\item $\forall E \in \mathfrak{E}: \mu(E) = \nu(E)$
		\item $\exists (E_n)_{n\geq 1} < \infty$ in $\mathfrak{E}$: $\bigcup_{n = 1}^\infty E_n = X$($\sigma$-finite)
	\end{enumerate}
\end{theorem}

\subsubsection*{Completeness of the Extension}

\begin{theorem}[continuity of measure]
	Suppose $\mu : \mathfrak{R} \rightarrow \mathbb{R}$ is a content defined on the ring $\mathfrak{R}$. Then the followings are equivalent:
	\begin{itemize}
		\item $\mu$ is a premeasure
	\end{itemize}
\end{theorem}

\begin{theorem}[Caratheodory Extension theorem]
\end{theorem}

Generator: Why do we always use the generator?\\
What is the $\sigma$-field defined by a random variable? -It is the domain.\\

Function class: continuous function with compact support\\
bounded linear operator\\
$L^1$ function\\
$L^p$ function\\

If we can find a continuous function that is not measurable(using Cantor function), then it is not in $L^1$, then we have a continuous function that is not integrable.

\begin{example}[Find a Lebegue measurable set that is not Borel measurable]
\end{example}

- Using Cantor function we can construct a function, where $\mu (g(C)) = 1$, thus $g^{-1}$ is a continuous but not measurable function. \\
The difference between Lebegue measurable function and Borel measurable function(The definition of the Lebegue measure. Why is the counter-example constructed using Cantor function is Lebegue measurable but not $\mathcal{L}-\mathcal{L}$ measurable.\\
Extension:
1. Extension of a field\\
2. Extension of the pre-measure\\
3. Extension from an outer measure to a measure.\\
3. Extension from Borel measurable function to Lebegue measurable\\ function.
3. Caratheodory\\



\subsection{The Lebesgue measure}
\begin{theorem}[Approximation theorem]
\end{theorem}



\section{Measurable and Integrable Functions}
Why is the Riemann integral not working for functions that are discontinuous on a set with positive measure.

\begin{theorem}[Young's Construction]
	For a non-negative function $f$ on $X$, f is measurable iff there exists a sequence of step functions $u_n$ with $u_n \uparrow f$. 
\end{theorem}
\begin{proof}
\end{proof}


\subsection*{numerical measurable functions}
\subsubsection*{Approximation by the simple functions}


\subsubsection*{countably generated measure space} 
First we could like to introduce the following definition:\\
1. \textbf{measurable isomorphism}:  \\
2. \textbf{$\mathfrak{G}$ separated $X$}: $\forall x, y \in X, x\neq y \exists A \in \mathfrak{G}: \chi_A(x) \neq \chi_A(y)$    \\
3. \textbf{$X$ is separated}: The $\sigma$-field $\mathfrak{A}$ separate $X$.\\
A space $(X, \mathfrak{A})$ is \textbf{countably generated}, if $\mathfrak{A}$ could be generated by a countable generator $\mathfrak{G}$: $\mathfrak{A} = \sigma{(\mathfrak{G}})$

\begin{lemma}
	$\mathfrak{G} \subset \mathfrak{A}$ separate $X$ iff $(X, \sigma(\mathfrak{G}))$ is separated
\end{lemma}

\begin{theorem}
	$(X, \mathfrak{A})$ is \textbf{countably generated} iff there exists a bijection $f: (X, \mathfrak{A}) \rightarrow ([0, 1], \mathfrak{B}([0, 1]))$.
\end{theorem}
\begin{proof}
	"$\Leftarrow$" is clear. For "$\Rightarrow$", see the function 
	\begin{align*}
		f:= \sum_{n = 1}^\infty 3^{-n} \chi_{A_n}
	\end{align*}
\end{proof}


\subsection{integrable function with compact support}
The functional spaces:
\begin{enumerate}
	\item $\mathcal{T}$: the step functions(simple functions). Have the form $f = \sum_{i = 1}^n \alpha \mu(A_i)$. 
\end{enumerate}
\begin{lemma}[the monotone convergence theorem]
	
\end{lemma}

\begin{theorem}
	$C_c(\mathbb{R}^n)$ is dense in $L^1(\mathbb{R}^n)$.
\end{theorem}
\begin{remark}
The technique of measure theory induction is applied here. Notice the result could be extended to $L^p$ space. See the chapter of separability.	
\end{remark}

\begin{corollary}
	\begin{align*}
		\forall \epsilon > 0 \exists \delta > 0 \forall E \in \mathfrak{A} \textbf{ with } \mu(E) < \delta:
		\int_E f d\mu < \epsilon
	\end{align*}
\end{corollary}


\subsection{integration of non-negative functions}

\subsection{The theorem of convergence}



\begin{theorem}[The dominated convergence theorem]
	Given measurable function $f_n, f : X \rightarrow \mathbb{K}$ with $f_n \overset{a.s.}{\rightarrow} f$. And $|f| \leq g$ a.s.  for some integrable function $g$. Then all $f_n$ and $f$ are integrable with:
	\begin{align*}
		\lim_{n \rightarrow \infty} \int f_n d\mu = \int f d\mu
	\end{align*}
	and 
	\begin{align*}
		\lim_{n \rightarrow \infty} \int |f_n - f| d\mu = 0
	\end{align*}
\end{theorem}
\begin{remark}
It's not enough to have $\mathbb{E}f_n < \infty$ for all $n$	
\end{remark}

\begin{example}[absolute continuous function]
	$f:[a, b] \rightarrow \mathbb{R}$ is differentiable and $f'$ is bounded. Then $f'$ is Lebesgue integrable over $[a, b]$ and 
	\begin{align*}
		\int_a^b f'd\lambda = f(a) - f(b)
	\end{align*}
\end{example}

\begin{theorem}[Fatu's lemma]
	For a sequence of measurable functions $f_n$ that non-negative, we have
	\begin{align*}
		\int \liminf f_n d\mu \le \liminf \int f_n d\mu 
	\end{align*}
\end{theorem}
\begin{proof}
\begin{align*}
	\int \lim_n \inf_{k \le n} f_k d\mu = \lim_n \int \inf_{k \le n} f_k d\mu \le \liminf \int f_n d\mu 
\end{align*}
\end{proof}

\begin{corollary}
	
\end{corollary}
\begin{proof}
\begin{align*}
	\int \lim_n \sup_{k \le k} f_k  = \lim_n \int \sup_{k \le n}f_k 
\end{align*}
\end{proof}



\newpage
\section{Product Measures}
The motivation of product measure is to define a measure in the cartesian product space which could be given by $\mu(A \times B) = \mu_X(A) \mu_Y(B)$. We want to explore the difference between 


\subsection{the product topology and the product measure}
\begin{theorem}
	A map $f: (Z, \mathfrak{G}) \rightarrow (X, \mathfrak{I})$ (where $\mathfrak{I}$ is the product $\sigma$-Algebra on the product space $X := \prod X_i$) is measurable iff it's measurable in every dimension.
\end{theorem}
\begin{proof}
\end{proof}
\quad\\
In the spirit of the last theorem we define the measurable function $j_K: \prod_{i \in I} X_i \rightarrow \prod_{i \in I} X_i$ :
\begin{equation*}
j_K(x) := \left\{
\begin{aligned}
	pr_i \qquad &i \in K\\
	a_i \qquad  &i \in I \backslash K \quad \textbf{for some $a_i \in X_i$}
\end{aligned}
\right.	
\end{equation*}
For any measurable set $M = \prod M_i \in \mathfrak{I}$ and $m_i \in M_i$ with respect to $a_i$, we can define the \textbf{cut} of $M$ as
\begin{align*}
	j_{I \backslash i} ^{-1}(\prod_{I \backslash i} \times a_i)
\end{align*}
By the last theorem we've established that every cut is measurable, therefore the following corollary
\begin{corollary}
	Every cut of a measurable set in the product $\sigma$-algebra is measurable.
\end{corollary}


\subsubsection{The measurability of the diagonal}
\begin{theorem}
	For the measurable space $(X, \mathfrak{A})$, the following claims are equivalent.
	\begin{enumerate}
		\item $X$ is separated by a countable $\mathfrak{G} \subset \mathfrak{A}$.
		\item There exists a bijection $f: (X, \mathfrak{A}) \rightarrow ([0, 1], \mathfrak{B}([0, 1]))$
		\item $\Delta_X \in \mathfrak{A} \otimes \mathfrak{A}$ 
		\item $\forall x \in X: \{x\} \in \mathfrak{A}$ and $(X, \mathfrak{A})$ is countably generated.
	\end{enumerate}
\end{theorem}
\begin{proof}
1)$\Rightarrow$2) trivial. \\
	3)$\Rightarrow$4): Notice: $\mathfrak{A} = \sigma(\mathfrak{G}) \Leftrightarrow \exists (A_n)_{n \geq 1} \in \mathfrak{G}: X = \bigcup_{n = 1}^\infty A_n$.(*)\\ 
	$\Delta_X \in \mathfrak{A} \bigotimes \mathfrak{A}$\\
	 $\Rightarrow \Delta_X \in \sigma(\mathfrak{G})$ with $\mathfrak{G} \subset \mathfrak{A} \otimes \mathfrak{A}$, countable.\\
	$\Rightarrow$ By (*) $\mathfrak{A} \otimes \mathfrak{A} = \sigma (\mathfrak{G})$\\
	$\Rightarrow$ $\mathfrak{A}$ is countably generated.
\end{proof}


\begin{definition}[Initial $\sigma$-Algebra and Product $\sigma$-Algebra]
	We have:
	\begin{itemize}
		\item\textbf{The initial $\sigma$-Algebra: } The smallest $\sigma$-Algebra on A, such that the family of function $\{f_i : A \rightarrow  (B, \mathfrak{B})| i \in I\}$ is measurable:
			\begin{equation*}
				\mathfrak{A} = \mathfrak{I}(f_i: i \in I) := \sigma(\cup f_i^{-1}(B_i) )
			\end{equation*} 
		\item \textbf{The product $\sigma$-Algebra}:
			The smallest $\sigma$-Algebra on A, such that the family of projection is measurable:
			\begin{equation*}
				\mathfrak{I}(pr_i: i \in I) := \bigotimes_{i \in I} \mathfrak{A}_i , \qquad pr_i: \prod X_j \rightarrow X_i, \quad  pr_i(x) = x_i
			\end{equation*}
		\item \textbf{Product Topology:} The smallest topology, such that the family of projection is continuous, a short deduction gives:
			\begin{equation*}
				(\prod X_j, \mathcal{T}) = \{ \prod_{\alpha \in I <\infty}  U_\alpha \times \prod_{\beta \neq \alpha} U_\beta \}
			\end{equation*}
	\end{itemize}
\end{definition}

\begin{remark}
Recall a \textbf{topology} is defined as a collection of subsets of: 1. contain $X$ and $\emptyset$, 2.any union, 3. pairwise intersection. 
 $B$ is a \textbf{subbase} if any subset in the topology is the union of the finite intersections of the subsets in the subbase. \par 
Given a space $X$ and a collection of subsets $B$ with $\cup_{U \in B} U = X$, the smallest topology that contains $B$ is the topological space generated by the subbase $B$.	
\end{remark}

\begin{theorem}
	$(X, \mathcal{T})$ is the \textbf{countable} product topology of the \textbf{separable} topological space $(X_i, \mathcal{T}_i$, then we have:
	\begin{equation*}
		\mathcal{B}(X) = \otimes_i \mathcal{B}(X_i)
	\end{equation*}
\end{theorem}

\begin{theorem}[Existence of Product measure]
$(X_j, \mathcal{A}_j, \mu_j)$ $(j = 1, \dots, n)$ \textbf{$\sigma$-finite measure space}, then there exists a \textbf{unique} measure:
\begin{equation*}
	\otimes_j \mu_j : \otimes \mathcal{A}_j \rightarrow \mathbb{R}^+, 
	\quad \otimes_j^n \mu_j(A_1 \times \dots \times A_n) = \prod_j^n \mu_j(A_j)
\end{equation*}
\end{theorem}


\subsection{Fubini theorem and the transformation rule}


\begin{theorem}[Fubini theorem]
	Given 2 \textbf{$\sigma$-finite} measure $\mu, \nu$ and define the canonical product measure $\mu \otimes \nu$ on the measurable space $(X \times B, \mathfrak{A} \otimes \mathfrak{B}$), given $f \in (A \times B, \mathfrak{A} \otimes \mathfrak{B}, \mu \otimes \nu)$ if one of the following criteria holds:
	\begin{enumerate}
		\item $f(x, \cdot)$ and $f(\cdot, y)$ non-negative.
		\item $f(\cdot, y)$ is integrable w.r.t $x$ for every $y$, and $f(x, \cdot)$ for every $x$.
		\item For $f$ measurable w.r.t the product $\sigma$-Algebra and either of the Integral is finite.
	\end{enumerate}
	Then we can change the order of integration.
\end{theorem}

\begin{theorem}[The Transformation Formula]
	Given $X, Y \subset \mathbb{R}^p$ open and $t: X \rightarrow Y$ a $C^1$-diffeomorphismus 
	\begin{enumerate}
		\item For all $A \in \mathcal{B}^p_X$, 
			\begin{equation*}
				\beta^p(t(A)) = \int_A |det Dt| d\beta^p
			\end{equation*}
	\end{enumerate}
\end{theorem}



\newpage
\section{Convergence Concept in the measure theory}
Throughout we assume the existence of $L^p$ space. The proof see Riezs-Fischer.


\subsection{Young, Holder and Minkowski}
\begin{proposition}[continuity of the convex functions on an Interval] Given a interval $I$, the convex function is continuous in the inner interval
\end{proposition}
\begin{proof}
	Page 241 of the Elstrodt.
\end{proof}

\begin{lemma}[Jensen's Inequality]
\end{lemma}
\begin{proof}
	The affine approximation and convex approximation.
\end{proof}

\begin{lemma}[Young's Ineqaulity]
if $1 < p, q < \infty$, $\frac{1}{p} + \frac{1}{q} = 1$, then:
\begin{equation*}
	ab \leq \frac{1}{p}a^p + \frac{1}{q}b^q
\end{equation*}
\end{lemma}

\begin{example}[An application of Jensen: Arithmetical and geometrical mean]
For $\sum_{i = 1}^n \alpha_i = 1$, $x_k > 0$, we have
	\begin{align*}
		\prod x_i^{\alpha_i} \leq \sum \alpha_i x_i 
	\end{align*}
\end{example}
\begin{proof}
	We see $\{ X := \{x_1, \dots, x_n\}, 2^X, \mu(x_i) = \alpha_i  \}$ defines a space with counting measure, For all measurable $Y: X \rightarrow\mathbb{R}$,  we therefore have by Jensen:
	\begin{align*}
		\mathbb{E}_X \exp(Y) \geq \exp(\mathbb{E} Y)
	\end{align*}
	choose $Y: y = \log x$, then
	\begin{align*}
		\sum_1^n \alpha_i x_i \geq \prod_1^n x_i^{\alpha_1} 
	\end{align*}
\end{proof}

\begin{lemma}[Holder's inequality]
$1 \leq p, q \leq \infty$, $\frac{1}{p} + \frac{1}{q} = 1$, $f, g$ measurable functions on $X$, we have:
\begin{align*}
	\|fg\|_{L^1} \leq \|f\|_{L^p} \|g\|_{L^q}
\end{align*}
\end{lemma}
\begin{proof}
	As we know from the results above:
	\begin{align*}
		\xi \eta \leq \frac{1}{p}\xi^p + \frac{1}{q}\eta^p
	\end{align*}
	Set $\xi = \frac{|f|}{\|f\|_p}, \eta = \frac{|g|}{\|g\|_q}$ and get the desired result.
\end{proof}
\begin{remark}
To apply the Holder usually we just need to make sure that $p > 1$.	
\end{remark}


\begin{theorem}[Minkowski Inequality]
	For $f, g \in L^1$ and $p \in [1, \infty]$, we have:
	\begin{align*}
		\| f + g \|_p \leq \|f\|_p + \|g\|_p
	\end{align*}
\end{theorem}
\begin{proof}
	\begin{align*}
		\|f + g\|^p_p =
		\int_\Omega |f + g|^p d\mu &= \int_\Omega |f+g|^{p-1} |f + g| d\mu \\
		&\leq \int_\Omega |f| |f+g|^{p-1} + |g| |f+g|^{p-1} d\mu \\
		&\leq (\|f\|_p + \|g\|_p) \|(f + g)^{p-1}\|_q\\
		&= (\|f\|_p + \|g\|_p) \|f + g\|_p^{\frac{p}{q}}\\
		(1)\Rightarrow
		\|f + g\|_p &\leq \|f\|_p + \|g\|_p
	\end{align*}
	In the case where division by infinity is not defined, we need to prove that $\|f + g\|_p$ is finite in (1).
\end{proof}

\begin{theorem}
	For $1 \leq p \leq \infty$ is $L^P$ a normed vector space and for $0 < p < 1$ is 
	\begin{equation*}
		d_p(f, g) := \|f - g\|^p_p
	\end{equation*}
	a metric in $L^p$
\end{theorem}


\subsection{The $L^p$ space}
Till now we constructed the basic topological structure of the $\mathcal{L}^p$ space. We established it's a vector space equipped with a metric(or in the case $p \geq 1$ a norm). Naturally we would like the space to be complete. The tools at hand are about pointwise convergence, with these tool we are able to deduce the $\mathcal{L}^p$ convergence.  

\begin{theorem}[Theorem of Riesz-Fisher]
	For $0 < p \leq \infty$ is $L^p$ a complete metric space w.r.t. $\|\cdot\|_p$. 
\end{theorem}
\begin{proof}
	The proof follows generally:
	\begin{itemize}
		\item Construct a pointwise convergent limit:
			\begin{enumerate}
				\item Use telescope sum to construct a subsequence in the form: $f_{n_k} = f_{n_1} + \sum g_i$. By the $L^p$ convergence of $f_n$ we know there exists a subsequence of $f_n$ such that 
				\begin{align*}
					\|f_{n_k} - f_{n_{k+1}}\|_p \leq 2^k\quad \forall m > n_k				\end{align*}
				define:
				\begin{align*}
					g_k := f_{n_k} - f_{n_{k+1}}
				\end{align*}
				\item $\lim_{k \rightarrow \infty} \sum_{i = 1}^k | g_i| := \hat{f}$ exists a.s. by bounded convergence theory since
				\begin{align*}
					 \|\sum |g_k|\|_p \le \sum \|g_k\|_p \le 1
				\end{align*}
				we have
				\begin{align*}
					\lim \|\sum|g_k|\|_p = \|\lim \sum |g_k|\|_p
				\end{align*}
				the pointwise limit is then well-defined. 
				\item \textbf{Construct the potential pointwise limit} 
					\begin{align*}
						\sum g_k = f_{n_1} - f_{n_k} \overset{a.s.}{\longrightarrow}  \hat{f}
					\end{align*}
					we define our potential pointwise limit
					\begin{align*}
						f:=  f_{n_1} - \hat{f}
					\end{align*}
			\end{enumerate}
		\item Show the convergence: At this point we have a a.s limit of the $L_p$ Cauchy sequence of function. We do not know if the limit is in $L^p$.
		\begin{enumerate}
			\item $\|f_n - f \|_p \rightarrow 0$
			\begin{align*}
				\int |f - f_n|^p d\mu = \int \liminf_k |f_{n_k} - f_m|^p d\mu \leq \liminf \int |f_{n_k} - f_m|^p d\mu 
			\end{align*}
			Since the sequence and the subsequence are $L^p$ convergent, we have that: $\liminf_k \int |f_{n_k} - f_m|^p d\mu \rightarrow 0$
			\item $f \in \mathcal{L}^p$
			\begin{align*}
				\int |f|^p d\mu = \int |f -f_m + f_m|^p d\mu 
				\leq 2^p \int |f-f_m|^p d\mu + 2^p \int |f_m|^p d\mu < \infty
			\end{align*}
		\end{enumerate}
	\end{itemize}
	The monotone convergence theory could be used to ensure the existence of a a.s. limit.
\end{proof}
	
\begin{corollary}[H.Weyl]
	For $0 < p \leq \infty$:
	\begin{enumerate}
		\item Every Cauchy sequence in $\mathcal{L}^p$ admits a a.s. convergent subsequence to $ f \in \mathcal{L}^p$
		\item Every convergent sequence in $\mathcal{L}^p$ admits a subsequence that converges a.s to a $f \in \mathcal{L}^p$.
	\end{enumerate}
\end{corollary}
	
\begin{corollary}
	For $1 \leq p \leq \infty$, $L^p$ is a Banach space and for $p \in (0, 1)$, $L^p$ is a complete metric space
\end{corollary}
\begin{remark}
Notice the for $p \in (0, 1)$, the norm function no longer satisfies the triangular inequality since it's not a convex function.	However $\|f\|_p^p$ is still a metric.
\end{remark}

\begin{example}[$L^p$ convergent yet not pointwise convergent sequence]
\end{example}

\begin{example}[Compactness of $L^p$ space]
	the $L^p$ space is in general not compact, in the sense that every $L^p$ bounded sequence needs not to admit a convergent subsequence.
\end{example}

\begin{theorem}
	Given $\mu(X) < \infty$, $0 < p < p' \leq \infty$, $L^{p'} \subset L^p$ and 
	\begin{align*}
		\|f\|_p \leq C(X, p, p') \|f\|_{p'}
	\end{align*}
	In another word, there exist a continuous embedding $L^{p'} \hookrightarrow L^p$.
\end{theorem}
\begin{proof}
	Write $\int_{M}|f|^p dm $ as $\int_M |f|^p \cdot 1 dm $ and apply the Holder's inequality. 
\end{proof}
\begin{remark}[Counterexample: $\mu(X) = \infty$]
Observe the sequence space $l^p(\mathbb{R}) \subset l^q(\mathbb{R})$ $\forall p \le q$. 
\end{remark}


\subsection{Convergence in measure space }
\begin{definition}[a.s pointwise]
Convergence pointwise except for a null set 
\end{definition}

\begin{definition}[a.s.uniform convergence AND almost uniform convergence]
There 2 concept of uniform convergence in measure space:
\begin{itemize}
	\item \textbf{a.s. uniform convergence}: 
		$\|f_n - f\|_\infty \rightarrow 0$  \textbf{on $X\slash N$} with $N$ a null set
	\item \textbf{almost uniform convergence}: $\forall \delta > 0 \exists E \textbf{ with } \mu(E) < \delta: \|f_n - f\|_\infty \rightarrow 0 \textbf{ on } X \slash E$
\end{itemize}
\end{definition}
\begin{remark}
Notice a.s is stronger since the null set is not dependant on the choice of $\delta$.	It's similar to the concept of uniform continuity. Take a look at the example $x \rightarrow x^n$ on the interval $[0, 1- \frac{1}{n}]$, we know the function is uniformly continuous on every $[0, 1 - \frac{1}{n}]$ but not $\lim [0, \frac{1}{n}]$. After taking the limit, the property of uniform continuity on a compact interval does not work.
\end{remark}

\begin{definition}[Convergence in measure]
There are 2 kinds of convergence in measure:
\begin{itemize}
	\item \textbf{Convergence in measure}: $\forall \epsilon > 0: \mu(\{f_n - f\} \geq \epsilon ) \rightarrow 0$
	\item \textbf{local convergence in measure:} $\forall \epsilon > 0 \forall A, \mu(A) < \infty: \mu \bigg(\{f_n - f > \epsilon\} \cap A \bigg) \rightarrow 0$
\end{itemize}
Notice if we unify the $\epsilon$: $\mu( \{f - g\} \geq \epsilon \} \leq \epsilon$, we may define a metric by $\|f - g\| := \{\inf \epsilon :  \mu( \{f - g\} \geq \epsilon \} \leq \epsilon\}$. We can therefore speak about the Cauchy sequence that converges i.M. 
\end{definition}
\begin{remark}
In stochastic the term is called convergence in Probability.	
\end{remark}


We'll now explore the relation between those convergence concept. First we study the relation between almost uniform and the a.s pointwise convergence. The conclusion is: almost uniform convergence implies a.s. pointwise convergence. The converse is true if we have a finite measure space. This is proved by the Russian mathematician \textbf{Jegorow}.

\begin{lemma}
	Given measurable function $f_n, f : X \rightarrow \mathbb{R}$, we have:
	\begin{itemize}
		\item $f_n \overset{a.s.}{\rightarrow} f \Leftrightarrow \forall \epsilon > 0:
		  \mu( \{x:   \bigcap_{n = 1}^\infty  \bigcup_{k = 1}^\infty |f_{n + k}(x) - f(x)| > \epsilon \}) = 0 $
		  \item $f_n \overset{a.s.}{\rightarrow} f \Leftarrow \forall \epsilon > 0 : \lim_{n \rightarrow \infty} \mu(\{x: \bigcup_{k = 1}^\infty | f_{n+k}(x) - f(x)| \geq \epsilon \}) = 0 $
		  \item $f_n \overset{a.s.}{\rightarrow} f \Leftrightarrow \forall \epsilon > 0  : \lim_{n \rightarrow \infty} \mu(\{x: \bigcup_{k = 1}^\infty | f_{n+k}(x) - f(x)| \geq \epsilon \}) = 0 $ if $\mu(X) < \infty$.		  
	\end{itemize}
\end{lemma}
\begin{proof}
	a) is a reformulation of the definition of convergence, b) follows directly form a) and c) follows form b) and the continuity of measure from above. 
\end{proof}
\begin{remark}
The equivalent definition of a.s. convergence could be applied to Cauchy-sequence as well.	
\end{remark}

\begin{lemma}
	Every \textbf{almost uniformly convergent} sequence $f_n$ admits an \textbf{a.s pointwise convergent} limit $f \in L^0$.
\end{lemma}
\begin{proof}
	Define $E_k$ with $\mu(E_k) < \frac{1}{k}$ and $f_n \overset{\infty}{\rightarrow} f$ on $X \slash E_k$. Observe $E:= \bigcap_{k = 1}^\infty E_k$\\ 
	The existence of $E$ is given by the fact that we have a $\sigma$-field. And 
	\begin{align*}
		\mu(E) = \lim_{k \rightarrow \infty} \mu (\bigcap_{k = 1}^{\infty} E_k) \leq \lim_{k \rightarrow \infty} \frac{1}{k} \rightarrow 0
	\end{align*}
	We can argue by contradiction that $\forall x \in E^c$, the function converges.
\end{proof}
\begin{remark}
From \textbf{uniform convergence on $E_n$} we cannot follow \textbf{uniform convergence on $\bigcup^\infty_1 E_n$}. However, we get \textbf{pointwise convergence on $\bigcup_1^\infty E_n$}.	
\end{remark}

\begin{theorem}[the Jegorow's theorem]
	Given a finite measure space($\mu(X) < \infty$),  \textbf{a.s pointwise convergence} implies \textbf{almost uniform convergence}.
\end{theorem}
\begin{proof}
	Uniform
\end{proof}

\begin{theorem}
	Almost uniform convergence implies convergence i.M.
\end{theorem}

\begin{theorem}
	Every i.M convergent sequence admits a almost uniform convergent subsequence
\end{theorem}
\begin{proof}
	By convergence in measure we have: $\forall n,m \geq n_k$
	\begin{align*}
		\mu(|f_m - f_n| \geq 2^{-k}) \le 2^{-k}
	\end{align*}
	Define:
	\begin{align*}
		A_k := \{ x \in \Omega| f_m - f_n \geq 2^{-k}\}
	\end{align*}
	we have 
	\begin{align*}
		B_l := \cup_{k = l}^\infty A_k ,\qquad 
		\mu(B_l) \le \sum_{k = l}^\infty 2^{-k} = 2^{-l}
	\end{align*}
	Then we have $\forall \mu(B_m^C) \le \delta), \forall k > l > m$:
	\begin{align*}
		|f_{n_k} - f_{n_l}| \le \sum_{i = l+ 1}^k |f_{i} - f_{i -1}| \le 2^{-m}
	\end{align*}
\end{proof}

\begin{theorem}
	Given a $\sigma$-finite measure space, every locally i.m convergent sequence admits an a.s. pointwise convergent subsequence.
\end{theorem}
\begin{proof}	
\end{proof}


\subsection{Convergence in $L^p$ space}
\begin{definition}[uniformly integrable]
A measurable sequence of function $f_n$ is uniformly integrable if:
\begin{itemize}
	\item $\forall \epsilon > 0 \exists E  \in \mathfrak{A} \textbf{ with } \mu(E) < \infty: \int_{E^c} |f|^p d\mu < \epsilon$
	\item $\forall \epsilon > 0 \exists \delta > 0 \forall E \textbf{ with } \mu(E)< \delta: \int_E |f|^p d\mu < \epsilon $ 
\end{itemize}
\end{definition}
The conditions could be summarised as:\par 
The functions are supported \textbf{arbitrarily well} on a finite-measure set. Therefore the integral cannot explode on a small set.\par 
For all the functions, the integral \textbf{is continuous at 0} with respect to the measure at 0. Therefore the measurable set supporting an arbitrary small integral cannot explode.

\begin{theorem}[The Vitali convergence theorem]
	For $0< p < \infty$ and $f_n, f$ measurable in $L^p$, if:
	\begin{itemize}
		\item $f_n \textbf{  u.i.}$
		\item $f_n \longrightarrow f$ locally in measure
	\end{itemize}
	Then we have:
	\begin{align*}
		f_n \overset{L^p}{\longrightarrow} f 
	\end{align*}
\end{theorem}
\begin{proof}.\\
	"$\Rightarrow$":\\
	"$\Leftarrow$"\\
	First we prove if $f_n \overset{a.s}{\longrightarrow} f$, the claim holds. 
	Select $E$ as in (1) of the definition of u.i. and $E \slash B$ as in the (2). We have:
	\begin{align*}
		\int_\Omega |f - f_n|^p d\mu \leq 2^p \int_{E^c} |f|^p + |f_n|^p d\mu 
		+ 2^{p-1} \int_{E \slash B} |f|^p + |f_n|^p d\mu + \int_B |f - f_n|^p d\mu 
	\end{align*}
	By Fatou:
	\begin{align*}
		\int_{E^c} |f|^p d\mu \leq \liminf \int_{E^c} |f_n|^p d\mu \leq \epsilon\\
		\int_{E\slash B} |f|^p d\mu  \leq \liminf \int_{E^c} |f_n|^p d\mu  \leq \epsilon
	\end{align*}
	Together with the definition of u.i., we conclude that the first terms approach to 0. And by the Jegorow's theorem we know on $B$ $f_n$ convergence uniformly therefore in $L^p$.\par 
	Now argue by contradiction, assuming 
	\begin{align*}
		\exists \epsilon >0 \forall k \in \mathbb{N} : \|f_{n_k} - f\|_p > \epsilon
	\end{align*}
	$f_{n_k}$ is a subsequence of a convergent(in measure) sequence therefore converges in measure. We get an a.s. pointwise convergent subsequence of $f_{n,k}$ and it convergence in $L^p$ to $f$, which contradicts the assumption.
\end{proof}
\begin{remark}
The technique $(f \pm g)^p \leq (\frac{2|f| + 2|g|}{2})^p \leq \frac{2^p |f|^p + 2^p|g|^p}{2} = 2^{p-1}(|f|^p + |g|^p)$ is often used in the context of $L^p$ convergence.	
\end{remark}
\newpage



\section{Absolute Continuity}
\begin{definition}[Absolute continuous function]
A function $F$ is absolute continuous if:$\forall \epsilon > 0 \exists \delta > 0$, such that for all finite partition $(a_i, b_i)$: 
\begin{equation*}
	\sum_{i = 1}^N b_i - a_i < \delta \Rightarrow \sum_{i = 1}^N |F(b_i) - F(a_i)| < \epsilon
\end{equation*}

\end{definition}
\begin{remark}
\end{remark}

\begin{definition}[Locally integrable function]
	\begin{equation*}
		L^p_{loc}(\Omega) := \{ f: \Omega \rightarrow V | f|_K \in L^p(\Omega) \forall K \subset \Omega \text{ compact} \}   \}
	\end{equation*}
\end{definition}

\begin{definition}[Lebesgue Point]
	For $f \in L_{loc}^1(\mathbb{R}^n), x \in \mathbb{R}$ is called a Lebesgue point of $f$ if:
	\begin{equation*}
		\lim_{r \rightarrow 0^+} \frac{1}{m(B_r(x))} \int_{B_r(x)} |f(x) - f(y)|dy = 0
	\end{equation*}
\end{definition}

\begin{definition}[Maximal Function]
	For $ f \in L^1_{loc}(\mathbb{R}^n)$ the \textbf{maximal function} $Mf: \mathbb{R}^n \rightarrow [0, \infty]$ is defined as:
	\begin{equation*}
		Mf(x) = \sup_{r > 0} \frac{1}{m(B_r(x))} \int_{B_r(x)} |f|dm
	\end{equation*}
\end{definition}

\begin{lemma}
	We define $f^0 := \limsup_{r \rightarrow 0^+} \frac{1}{m(B_r(x))} \int_{B_r(x)} |f|dm$.For every $N \in \mathbb{N}$, 
	\begin{equation*}
		A_N := \{ x \in \mathbb{R}^n | f^0(x) > \frac{1}{N} \}
	\end{equation*}
	Then $A_N$ is contained in a Lebesgue measurable set with measure less than $\frac{1}{N}$.
\end{lemma}

\begin{theorem}[Lebesgue Differentiation theorem]
	For $f \in L_{loc}^1(\mathbb{R}^n)$. Almost every point is a Lebesgue point.
\end{theorem}
\begin{remark}
To see why it is called the differentiation theorem, observe:
\begin{equation*}
	\frac{F(x + h) - F(x)}{h} - f(x) = \frac{1}{h}\int_x^h f(t) - f(x) dt \leq 2 \frac{1}{m(x-h, x + h)} \int^{x+h}_{x-h} |f(t) - f(x)| dt \rightarrow 0 
\end{equation*}	
\end{remark}


\begin{corollary}
	$\forall f \in L^1([a, b])$, $F := \int_a^b f(t) dt$ is differentiable a.e. and $F' = f$ 
\end{corollary}

\begin{theorem}[Fundamental theorem of Calculus for the Lebesgue Integral]
\end{theorem}

\begin{theorem}[Radon-Nikonym]
	Given $\mu$ a $\sigma$-finite measure and $\nu \ll \mu$ a signed measure on $\mathcal{A}$, then there exist a quasi-integrable function $f: X \rightarrow \bar{\mathbb{R}}$, such that $\nu = f \circ \mu$ . 
\end{theorem}
\begin{proof}
	w.l.o.g. we only proved the case of positive measure.
	\begin{itemize}
		\item \textbf{Prove the lemma:} If $\mu, \nu$ are finite measure and $\nu \leq \mu$, then there exists a $\rho$, such that $\nu = \rho \circ \mu$ with $\rho: X \rightarrow [0, 1]$. And $h$ is unique in $L^2$. \par
			 
		\begin{enumerate}
			\item For a partition $P_1$ of $X$, define the stepfunction:
				\begin{equation*}
					p_1(x) =\left\{ \begin{aligned}
						\frac{\nu(A)}{\mu(A)} \quad &\mu(A) \neq 0 \\
						0 \quad &\mu(A) = 0
					\end{aligned}
					\right.
				\end{equation*}
			Given a finer partition $P_2$, we know $\forall A \in P_1:$
				\begin{equation*}
					\int_A p_1 d\mu = \int_A p_2 d\mu 
				\end{equation*} 
			\item Define:
				\begin{equation*}
					\alpha = \sup\{\int_X p^2 d\rho |\text{ $P$ a partition of X} \}
				\end{equation*}
				and $p_n$ such that:
				\begin{equation*}
					\alpha \geq \int_X p_n^2 d\rho \geq \alpha - \frac{1}{n} 
				\end{equation*}
				then for the finest common partition:
				\begin{equation*}
					\int_X (p_n - p_m)^2 d \mu = \int_X p_n^2 - p_m^2 d\rho \leq \frac{1}{m}
				\end{equation*}
				which makes it a Cauchy sequence. Since $L^2$ is complete, we know $\|p_n - h\|_2 \rightarrow 0$ for a $h \in L^2$
		\end{enumerate}
		\item \textbf{Prove the case: $\nu$ and $\mu$ are finite.}\par 
		We define $\tau = \mu + \nu$, sine now $mu \leq \tau$ and $\nu \leq \tau$, we can apply the lemma above, and get $\tau = g \circ \mu$ and $\tau = h \circ \nu$. Now define $N:= \{M \subset X | \mu(M) = 0\}$ and $f$:
		\begin{equation*}
			f(x) = \left\{ \begin{aligned}
				\frac{h(x)}{g(x)} \quad & x \in N^c\\
				0 \quad & x \in N
			\end{aligned}\right.
		\end{equation*}
		then $\forall A \in \mathcal{A}$( $N$ excluded):
		\begin{equation*}
			\nu(A) = \int_A h d \tau = \int_A f \circ g d\tau = \int_A f d \mu
		\end{equation*}
		\item Prove the case: $\mu$ is finite, $\nu$ arbitrary. 
		\item Prove the case: $\mu$ $\sigma$-finite, $\nu$ arbitrary. 
	\end{itemize}
\end{proof}


\subsection{Hahn decomposition}
\begin{definition}[the signed measure]
	Given $\mathcal{A}$ the $\sigma$-algebra, the function $\nu: \mathcal{A} \rightarrow \bar{\mathbb{R}}$ is called the \textbf{signed measure} if 
	\begin{enumerate}
		\item $\nu(\emptyset) = 0$,
		\item $\nu(\mathcal{A}) \subset (-\infty, \infty]$ or $\nu(\mathcal{A} \subset [-\infty, \infty)$,
		\item $\sigma$-additivity. 
	\end{enumerate}
\end{definition}


\begin{theorem}[Hahn decomposition]
Let $\nu$ be a signed measure $(X, \mathcal{A}) \rightarrow \bar{\mathbb{R}}$. For every $\mu$ there exists a decomposition of $X$: $X = P \cup N$ ($P, N \in \mathcal{A}$) such that 
\begin{enumerate}
	\item $\forall A \in \mathcal{A}, A \subset P:\quad \nu(A) \ge 0$, i.e., $P$ is $\nu$-positive. 
	\item $\forall B \in \mathcal{A}, B \subset N: \quad \nu(B) \le 0$, i.e., N is $\nu$-negative. 
	\item $P$ and $N$ are unique outside a $\nu$-null set. 
\end{enumerate}
\end{theorem}




\section{Measure in topological spaces}
\begin{definition}
	Given a measurable space $(X, \mathcal{F})$, a measure $\mu: \mathcal{B} \rightarrow [0, \infty]$ is called 
	\begin{enumerate}
		\item \textbf{locally finite measure (Borel measure)}: 
			\begin{equation*}
				\forall x \in X \exists U \ni x: \mu(U) < \infty
			\end{equation*}
		\item \textbf{Inner \& outer regular measure}: $\mu$ is inner regular if:
			\begin{equation*}
				\forall A \in \mathcal{B}: \mu(B) = \sup_{K \in \mathcal{B} \text{compact}} \mu(K) 
			\end{equation*}
		\item \textbf{Radon-measure:} a inner regular Borel-measure.
		\item \textbf{tight}
	\end{enumerate}
\end{definition}

\begin{theorem}[The theorem of Ulam]
\end{theorem}

\begin{proof}
	\begin{itemize}
		\item \textbf{The lemma:} Every finite measure on a $\sigma$-Algebra, the set:
			\begin{equation*}
				\{A \subset \mathcal{A}: A \text{ is  $\mu$-regular}\}
			\end{equation*}
			is a $\sigma$-ring
	\end{itemize}
\end{proof}

\begin{definition}[$\mu$-boundry-less set]
In the measure space $(X, \mathcal{B}, \mu)$, a set $B$ is called $\mu$-boundary-less, if:
\begin{equation*}
	\mu(\partial B) = \mu(\bar{B} \setminus \overset{\circ}{B}) = 0
\end{equation*}
\end{definition}

\begin{theorem}[Portmanteau-Theorem]
	For a metric space $X$, $\mu_n, \mu$ finite Borel measure, then the followings are equivalent:
	\begin{enumerate}
		\item $\mu_n \overset{w}{\rightarrow} \mu$
		\item $\forall F \subset X \text{ closed}: \limsup \mu_n(F) \leq \mu(F)$
		\item $\forall U \subset X open: \liminf \mu_n(U) \geq \mu(F)$
	\end{enumerate}
\end{theorem}
\newpage





\chapter{Linear PDE}
\section{Introduction}
\subsection{Modelling, origin, and the divergence structure}

\begin{example}
The usual PDE includs
\begin{itemize}
	\item the Poisson equation: $\Delta u = f$. 
	\item the Heat equation(diffusion equation): 
\end{itemize}	
\end{example}

\begin{example}[Heat Equation]
A quantity is changing by moving or by producing(consumption)
\begin{equation*}
	\frac{\partial}{\partial t} \int_\Omega \rho(t, x) dx = \int_\Omega h(t, x) dx  - \int_{\partial \Omega} j(t, x) \cdot \nu dx 
\end{equation*}
applying integral by parts and DUIS
\begin{equation*}
	\int_\Omega \frac{\partial}{\partial t} \rho - h + div(j) dx =0
\end{equation*}
and we get
\begin{equation*}
	\rho_t = -div(j) +h
\end{equation*}
In the case of heat conduction, where we have $\rho$ as temperature and $ j$ the heat flux density. The Fourier law implies $j = -\nabla \rho$(implying heat flow from hotter to cooler spot). We get the heat equation:
\begin{equation}
	\rho_t = div(\nabla \rho)
\end{equation}
To be more precise, we define:
\begin{itemize}
	\item The internal energy: $\rho = c(x)\theta(t, x)$
	\item The heat flux:  $ q = -K(x)\nabla \theta $ according Fourier's law, where $K(x)$ a SPD matrix.
\end{itemize}
\end{example}

\subsection{Superposition principle and separation of variables}
We observe the wave equation:
\begin{equation*}
	\Delta u = 0
\end{equation*}
\subsubsection*{Superposition}
Superposition means linearity, indicating if $x_1, x_2$ satisfies the PDEs then the linear combination must also satisfies the PDEs. 

\subsubsection*{Separation of variables}
We guess the solution of the wave equation could be written as a multiplicative form(due to superposition principle the additive ansatz is not very useful):
\begin{equation*}
	u(x) = \varphi(x_1)\psi(x_2), \qquad x \in \mathbb{R}^2
\end{equation*}
and the wave equation could be written as 
\begin{align*}
	0 &= \varphi''(x_1)\psi(x_2) + \varphi(x_1)\psi''(x_2)\\
	0 &= \frac{\varphi''(x_1)}{\varphi(x_1)} + \frac{\psi''(x_2)}{\psi(x_2)}(\textbf{assuming } \varphi, \psi \neq 0)
\end{align*}
From above we set 
\begin{equation*}
	\frac{\varphi''(x_1)}{\varphi(x_1)} = - \frac{\psi''(x_2)}{\psi(x_2)}:= \lambda
\end{equation*}
It remains to solve the ODE with the free variable $\lambda$ :
\begin{align*}
	\varphi''(x_1) &= \lambda \varphi(x_1)\\
	\psi''(x_2) &= -\lambda \psi(x_2)
\end{align*}

\begin{remark}[What kind of linear homogeneous PDE can't be solved by separation of variable?]
\end{remark}

\subsection{Separation of Variable:1-D wave equation}
Observe the simplified 1-dim(homogeneous) wave equation\\
\begin{equation}\label{WE}
\left\{
\begin{aligned}
	&u_{xx} = u_{tt} \\
	&u(0, x) = u_0(x) \qquad (t, x) \in \mathbb{R} \times \mathbb{R}\\
	&u_t(0, x) = u_1(x)
\end{aligned}
\right.
\end{equation}
There are many ways to derive this PDE. One of the most common is based on Hook's Law, where the second time derivative represents the kinetic force by Newton and the Laplacian represents the elastic force. And we have the initial position and the initial speed. Now we solve the PDE: \\
1. Ansatz: the special wave/transport solution:
\begin{equation*}
	u(t, x) = G(x - ct)
\end{equation*}
(The form of solution is an educated guess since the object moves with a constant speed, we need to prove it solves the PDE)
we have
\begin{align*}
	u_{tt} &= c^2 G''(x - ct)\\
	u_{xx} &= G''(x-ct)
\end{align*}
The ansatz works when $c^2 = 1$. Assuming(superposition principle applied)
\begin{equation*}
	u(t, x) = G_1(x - t) + G_2( x + t)
\end{equation*}
2. \textbf{Tranform this to the one-dimensional transport equation.} The PDE $u_{tt} = u_{xx}$ could be decomposed as
\begin{equation*}
	u_{tt} - u_{xx} = (\partial_t + \partial_x)(\partial_t - \partial_x)u 
\end{equation*}
Define
\begin{equation*}
	v:= \partial_t u - \partial_x u
\end{equation*}
and the PDE becomes a homogeneous transport equation
\begin{equation*}
\left\{
\begin{aligned}
	&\partial_t v = - \partial_x v\\
	& v(0, x) =	u_1(x) - u_0'(x)
\end{aligned} 
\right.
\end{equation*}
By the result in 2.2 we get a solution
\begin{equation*}
	v(t, x) = u_1(x - t) - u_0'(x - t)
\end{equation*}
Now we use $v$ to solve $u$ by solving
\begin{equation*}
\left\{
\begin{aligned}
	&u_t = u_x + v(t, x)\\
	&u(0, x) = u_0(x)
\end{aligned}
\right.
\end{equation*}
Apply the result in 2.2 \ref{TE: solution} we get
\begin{align*}
	u(t, x) &= u_0(x + t) + \int_0^t v(x + (t-s))ds\\
			&=u_0(x + t)  - u_0(x) + \int_{x - t}^0 u_1(y)dy
\end{align*}
Combine $c = 1$ and $c = -1$ we get the d'Alembert formula:
\begin{equation} \label{d'Alembert}
	u(t, x) = u_0(x-t) + u_0(x+t) + \int_{x-t}^{x+t} u_1(y)
\end{equation}

\subsection{Uniqueness of the initial value boundary problems}
The integral is a linear operator. For a linear problem, to show uniqueness is to show the positive definiteness of the linear operator.
\subsubsection*{Transport Equation}
\begin{theorem}[Uniqueness of solution of the Transport Equations]
	Given $f \in C^0(\mathbb{R} \times \mathbb{R}^d)$, $u_0 \in C^1(\mathbb{R}^d)$. Then the transport equation \ref{TE}  admits an unique solution given by the form \ref{TE: solution} .
\end{theorem}
\begin{proof}
	The proof has the following steps:
	\begin{enumerate}
		\item Introduce the traucated function $\varphi \in C^1(\mathbb{R} \times \mathbb{R}^d)$
		\begin{equation*}
			\varphi = 
			\left\{
			\begin{aligned}
				&1 \qquad x \in [-R, R]\\
				&0 \qquad x \in (-\infty, -R-1] \cup [R + 1, \infty)\\
				& \geq 0 \qquad elsewhere
			\end{aligned}
			\right.
		\end{equation*}
		\item 
		define $w := u - \tilde{u}$, where $u$ and $\tilde{u}$ are 2 solutions of the PDE and we get rid of the inhomogeneity    and obtain trial initial condition. By superposition it's also a solution
		\item then define the bilinear form of $w$:
		\begin{align*}
			E(t) := \int_{\mathbb{R}^d} \varphi(x + ta) w^2(t, x) dx
		\end{align*}
		to show $w \equiv 0$ we show $E(w) \equiv 0$. $E$ is well-defined since it has support of $\hat{B}:= B_{R+1}(-ta)$		
		\item Observe that
		\begin{equation*}
			E(0) = 0
		\end{equation*}
		\begin{align*}
			\dot{E}(t) &= \frac{d}{dt} \int_{B_{R+1}(-ta)} \varphi(x + ta) w^2(t, x)dx\\
		 (1)&= \int_{\hat{B}_R(0)} \frac{d}{dt} (\varphi(x+ ta) w^2(t, x)) dx\\
			&= \int a\varphi_{x+ta}(x + ta) w^2(t, x) + 2w w_t(t, x) \varphi(x + ta) dx\\
		 (2)&= \int a \varphi_{x + ta}(x+ta)w^2(t, x) + 2w (aw_x(t, x))\varphi(x+ta) dx\\
			&= \int a \nabla(\varphi w^2)dx\\
		 (3)&= \int_{x \in \partial \hat{B}} a \varphi w^2 \cdot \nu dx -  \int_\Omega \varphi \omega^2 div(a) dx\\
		 (4)&= 0 \qquad \forall R \in \mathbb{R}^d
		\end{align*}
		where (1) we use DUIS, since we have the truncated function the integrand admits support of a closed ball and therefore integrable. And given any $x$ the integrand is continuous with respect to $t$. 
		\\ (2) the PDE and (3) the integration by parts.\par 
		(4) works since for all $R$ the boundary value of $\varphi$ is 0  and $div(a)$ is 0. We conclude the last line.
		Therefore $E(t) = 0$ $\forall t \in \mathbb{R}$. Together we have
		\begin{align*}
			&\forall (t, x) \in \mathbb{R} \times \mathbb{R}^d: E(t) = 0\\ \Rightarrow &\forall (t, w) \in \mathbb{R} \times \mathbb{R}^d: w(t, x) = 0
		\end{align*}
	\end{enumerate}
\end{proof}

\subsubsection*{Heat Equation}
Given $\Omega$ a bounded domain with $C^1$ boundary. Observe the heat equation with initial boundary condtion
\begin{equation}\label{HE_Dir}
	\left\{
	\begin{aligned}
		&\partial_t u = div(K(x) \nabla u(t, x))\\
		&u(0, x) = u_0(x)\\
		&u(t, x) = g(t, x) \quad x \in \partial \Omega
	\end{aligned}
	\right.
\end{equation}
Analogously we observe the homogeneous solution $w$ and $E(t)$. 
\begin{align*}
	\dot{E}(t) &= \frac{d}{dt} \int w^2(t, x)dx\\
	(1)&= \int \frac{d}{dt}(w^2(t, x)) dx\\
	   &= \int 2w \cdot w_t dx \\
 	   &= \int 2w \cdot div(K \nabla w) dx\\
	   &= \int_{x \in \partial \Omega} (w \cdot K \nabla w) \cdot \nu dx - \int \nabla w \cdot K \nabla wdx \\
	(2)&= 0  -  non.neg\\
	   &\leq 0
\end{align*}
Also $E(t) \geq 0$ and $E(0) = 0$ by construction. Therefore we conclude $\dot{E}(t) = 0$ and $E(t) = 0$ for all $t \in \mathbb{R}$.
The Initial-Boundary value problem \ref{HE_Dir} is one of the boundary value types. Through the proof we noticed that to obtain $0$ in the first term in step $(2)$, \textbf{we can also let $K \nabla w$ be 0}. We introduce the second type of \textbf{Heat equation} with initial boundary data:

\begin{equation}\label{HE_Neu}
	\left\{
	\begin{aligned}
		&\partial_t u = div(K(x) \nabla u(t, x))\\
		&u(0, x) = u_0(x)\\
		&u(t, x) = g(t, x) \quad &x \in \Gamma_{Dir}\subset \partial \Omega\\
		&K(x) \nabla u(t, x) = h(x) \quad &x \in \Gamma_{Neu} \subset \partial \Omega
	\end{aligned}
	\right.
\end{equation}

\subsubsection*{Wave Equation}
The wave equation with initial-boundary conditions has the form:
\begin{equation}
\left\{
	\begin{aligned}
		&\rho(x)\partial^2_t u(t, x) = \partial_x (A(x) \nabla u) + f(t, x)\\
		&u(0, x) = u_0(x)\\
		&\partial_t(0,x) = u_1(x)\\
		& u(t, x) = h(t, x) \quad &\forall x \in \Gamma_{Dir}\\
		& (A(x) \nabla u(t, x)) \cdot \nu = g(t, x) \quad &\forall x \in \Gamma_{Neu} 
	\end{aligned}
\right.
\end{equation}
\begin{theorem}[Uniqueness of the Wave Equation]
	Consider a wave equation above with the form above and $\partial_t A = 0, \partial_t \rho = 0$. The homogeneous solution of the wave equation in $C^2(\mathbb{R} \times \mathbb{R}^d)$ satisfies the energy conservation law
	\begin{equation}
		\int_\Omega \frac{\rho}{2}u_t^2 + \frac{1}{2}\nabla u A \nabla u dx = \int_\Omega \frac{\rho}{2} u_1^2 + \nabla u_0 A \nabla u_0 dx 
	\end{equation}
\end{theorem}
\begin{proof}
	\begin{align*}
		\dot{E}(t) &= \frac{d}{dt} \int_\Omega \frac{\rho}{2} u_t^2 + \frac{1}{2} \nabla u A \nabla u dx\\
		&= \int \rho u_t u_{tt} + A \nabla u \cdot \nabla u_t dx\\
		&= \int \partial_x (A \nabla u) u_t + (A\nabla u) \nabla u_t + f u_t dx\\
		&= \int_{\partial \Omega} (A \nabla u) u_t \cdot \nu dx + \int_\Omega (A \nabla u) \nabla u_t - (A \nabla u) \nabla u_t + f u_t dx\\
	 (1)&= \underline{ \int_{\partial \Omega} (A \nabla u) \cdot \nu \cdot u_t} + \int_\Omega f u_t  dx
	\end{align*}
	In the homogeneous case we have $f\equiv 0$ and $g \equiv 0, h \equiv 0$. Therefore $\dot{E} \equiv 0$.
\end{proof}
\begin{remark}
Notice in the last line, to prove the uniqueness of wave equation, we only need the product of the flux and the boundary value to be zero. 	
\end{remark}


\begin{corollary}
	Given the same inhomogeneity and boundary condition with different initial data. The solution of the wave equation is continuously dependent on the initial data in the sense that
	\begin{equation}
		E(u - \tilde{u}) \geq \frac{\inf \rho}{2} \|u - \tilde{u}\|_2^2 + \alpha_{\min}^A \|\nabla (u - \tilde{u})\|_2^2
	\end{equation}
\end{corollary}



\subsection{Exercises}
\begin{exercise}[Integration by parts]
	Given $\alpha \in C^1(\Omega)$, $\beta \in C^1(\Omega, \mathbb{R}^d)$, 
	prove
\begin{equation*}
	\int_\Omega \alpha \textbf{ div}\nu dx = \int_{\partial \Omega} \alpha \beta \cdot n da  -\int_\Omega \nabla \alpha \cdot  \beta dx
\end{equation*}
\end{exercise}

\begin{exercise}[$|5|$ Gronwell's inequality(estimating the energy functions)]
Assume that $y \in W^{1, 1}([0, \infty))$ and $\phi \in L^1([0, \infty))$. and 
\begin{align*}
	\forall t \geq 0: \frac{d}{dt} y^2 \leq 2 \phi y 
\end{align*}
show
\begin{align*}
	\forall t : y(t) \leq |y(0)| + \int_0^\infty \phi(t) dt
\end{align*}
	
\end{exercise}


\begin{exercise}[1-dimensional homogeneous transport equation]
Given the transport equation($u \in C^1(X, T)$, $f \in ?$:
\begin{equation*}
\left\{
\begin{aligned}
	&u_t = bu_x + f(x, t)\\
	&u(0, x) = u_0(x)
\end{aligned}
\right.
\end{equation*}
Prove that the solution is given by the form:
\begin{equation*}
	u(t, x) = u_0(x + bt) + \int_0^t f(s, x + (t - s)b) ds
\end{equation*}
\begin{proof}
	Define $a(t) := x + b(t - s)$, then
\begin{align*}
		u_t &= u_0'(x + bt) \cdot b + \int_0^t b \cdot  \partial_2 f(s, x + (t - s)b)  ds + f(t, x) \\
		u_x &= u_0'(x + bt) + \int _0^t  \partial_2 f(s, x + (t -s)b) ds
	\end{align*}
\end{proof}
\end{exercise}

\begin{exercise}[Differentiation under the integral sign]
\end{exercise}

\begin{exercise}[$|7|$ Parabolic equation with Robin boundary condition]
	Consider the parabolic equation:
	\begin{align*}
		\partial_t u = \partial_x (k(t, x) \partial_x u) + b(t, x) u + f(t, x)
	\end{align*}
\end{exercise}

\begin{exercise}[$|8|$ Characteristic Method(Linear transport problems on a finite interval)]
Given the transport equation:
\begin{equation}
\left\{
\begin{aligned}
u_t + au_x &= bu \\
u(0, x) &= u_0(x)\\
u(t, l) &= g(t)	
\end{aligned}
\right.	
\end{equation}
and in the overlapping region the compatibility condition must hold. The tasks are:
\begin{enumerate}
	\item find a non-trivial solution
	\item prove the uniqueness when $a < 0$
	\item construct a non-trivial solution to the homogeneous problem for $a > 0$. 
\end{enumerate}


\begin{proof}
Ansatz: by inserting $t$ into the location variable we can get 2 terms by taking the derivative of $t$:\par 
Define: $	\tilde{u}(t, y) = u(t, y + at)$, then PDE could be transfomed to :
\begin{equation}
\left\{
\begin{aligned}
	\tilde{u}_t &= b\tilde{u}\\
	\tilde{u}(0, y) &= u_0(y)\quad 0 < y + at < l \\
	\tilde{u}(t, l - at) &= g(l - at)
\end{aligned}
\right.	
\end{equation}
we transformed the PDE to 2 set of ODEs with boundary conditions. Solving the ODEs we get:
\begin{align*}
	\tilde{u} &= u_0(y) e^{bt} \quad y \in (0, l) \\
	\tilde{u} &= g(T(y)) e^{b[t - T(y)]} \quad T(y) = \frac{l - y}{a} \geq 0
\end{align*}
Change the variable by $y + at = x$ we get:
\begin{align*}
	u(t, x) = u_0(x - at)
\end{align*}
\end{proof}
\begin{proof}
\end{proof}

\begin{proof}
	We try to construct a solution that is meaningful. Given semi-trivial boundary condition:
	\begin{align*}
		u_0 &\equiv 0 \qquad x >0\\
		u_0 &> 0  \qquad x \in (-l, 0)
	\end{align*}
	then by characteristic we have:
	\begin{align*}
		u(t, x) = u_0(x - at) e^{bt}
	\end{align*}
	For $t \in (0, \frac{l}{a})$ we get a non-trivial solution even for the trivial boundary condition.
\end{proof}
\textbf{For the 1st order PDE, the method of characteristic finds curves, along which the PDE becomes ODE. After solving the ODE, it can be transformed to the original PDE}
\end{exercise}



\section{The Method of Characteristic}


\subsection{Method of Characteristic}
We observe the 2d semi-linear PDE
\begin{equation}
	a(x, y, u)u_x + b(x, y, u)u_y  = c(x, y, z)
\end{equation}
It's important to understand the \textit{characteristic curve.} Suppose $u$ is a solution and we observe the integral surface defined by
\begin{align*}
	\mathcal{S} = \bigg(x, y, u(x, y) \bigg) \subset \mathbb{R}^3
\end{align*}
The normal vector of $\mathcal{S}$ is define by
\begin{align*}
	\bigg(u_x \frac{dx}{dt} , u_y \frac{dy}{dt} , -1\bigg) 
\end{align*}


\begin{definition}[Method of characteristic]
Observe the PDE:
\begin{align*}
	\left\{
	\begin{aligned}
		&F(Du, u, x) = 0\\
		&u = g \quad \texttt{on } \Gamma
	\end{aligned}
	\right.
\end{align*}
$F$ possibile non-linear. $F, g \in C^2(\mathbb{R}^n \times \mathbb{R} \times \Omega)$. Introducing the time variable $s$ and define
	\begin{align*}
		z(s) &:= u(x(s))\\
		p(s) &:= Du(x(s))\\
		\gamma(s) &:= x(s)
	\end{align*}
Then the ODE satisfies $ (p(.), z(.), \gamma(.)) \in C^1(T; \mathbb{R}^n \times \mathbb{R} \times \Omega)$ and :
	\begin{align*}
		\begin{bmatrix}
			\dot{p}(s)\\
			\dot{z}(s)\\
			\dot{\gamma}(s)
		\end{bmatrix}
	= 
		\begin{bmatrix}
			- F_z p - F_x \\
			p \cdot F_p\\
			F_p
		\end{bmatrix}
	\end{align*}
	where we abbreviates 
	\begin{align*}
		&\left\{
	\begin{aligned}
		&\dot{p}(s) =-D_x F - D_z F \cdot p \\
		&\dot{z}(s) = p(s) \cdot F_p\\
		&\dot{\gamma}(s) = F_p(s)
	\end{aligned}
	\right.\\
	&\texttt{with}\\
	&p(0) = \nabla g(x_0), \quad z(0) = g(x_0), \quad \gamma(0) = x_0, \quad x_0 \in \Gamma
	\end{align*}
\end{definition}
\begin{remark}
Notice it's a $2n + 1$ ODE system. We sometimes call $\gamma$ the \texttt{projected characteristic}, meaning it's the projection of the full characteristic in $\mathbb{R}^{2n + 1}$ onto the physical space  $U \subset \mathbb{R}^n$.
\end{remark}

 

\begin{theorem}
	Let $u \in C^2(\Omega)$ solve the nonlinear first-order PDE at the start of the chapter. 
	Suppose $\gamma \in C^1(\xi, \Omega)$ for an open interval of the form above
	\begin{align*}
		\dot{\gamma}(t) = F_p (Du(\gamma(t)), u(\gamma(t)), \gamma(t)) 
	\end{align*}
	solves the $2n+1$ ODE characteristic system, then u is a solution of:
	\begin{align*}
		F(Du, u, x) = 0
	\end{align*} 
\end{theorem}
\begin{proof}
Observe:
\begin{align*}
	\dot{p}^i(s) = \sum_{j = 1}^n u_{x_i x_j}(x(s)) \dot{x}^j(s)
\end{align*}
By $F = 0$ we know:
\begin{align*}
	\frac{\partial F}{\partial x_i} = \sum_{j = 1}^n F_{p^j}u_{x_i x_j} + F_z u_{x_i} + F_{x_i} = 0
\end{align*}
\textbf{Assume:}
\begin{align}
	\dot{x}^j(s) = F_{p_j}(s)
\end{align}
Then we have:
\begin{align*}
	\sum_{j = 1}^n u_{x_i x_j} \dot{x}^j + F_z p^i + F_{x_i} = 0\\
	\Rightarrow \dot{p}^i = - F_z p^i - F_{x_i}
\end{align*}
Differentiate $z$ and $x$ w.r.t. $s$ we get:
\begin{align*}
	\dot{z}(s) = \sum_{j = 1}^n z_{x_j} \dot{x}_j = \sum_{j = 1}^n p_j F_{p_j} \quad\texttt{(assumption above)}
\end{align*}
\end{proof}



\begin{example}[linear first-order ODE.]
\begin{align*}
	F(Du, u, x) = b(x)u_x + c(x)u = 0
\end{align*}
Calculate the characteristic curve:
\begin{align*}
	\dot{x}(s) &= b(x(s)) \cdot p(x(s))\\
	\dot{z}(s) &= - c(x(s))  \cdot p(x(s))
\end{align*}
Altogether:
\begin{align*}
	\dot{x}(s) &= b(s)x(s) \\
	\dot{z}(s) &= -c(s)z(s) 
\end{align*}
Problem: $u$ is discontinuous at $0$ unless $g(0) = 0$. 
\end{example}

\subsubsection*{Local existence}
Continue towards the local existence for the general BVP. Assume:
\begin{align*}
	F \in C^2(\mathbb{R}^n \times \mathbb{R} \times \mathbb{R}^n)
\end{align*}
and $\Omega \subset \mathbb{R}^n$ open with boundary $\partial \Omega$. that is $C^2$ in a neighborhoods $x_0 \in \Gamma$ with the unit outer normal
\begin{align*}
	\nu(x_0) \in \mathbb{R}^n, |\nu(x_0)| = 1
\end{align*}
The \textbf{(BVP)}
\begin{align*}
	F(Du, u, x) = 0 &\quad x \in \Omega\\
	u = g &\quad \texttt{on } \Gamma
\end{align*}
\textbf{Motivation:}
tangential derivative along $\Gamma$ asserts 
\begin{align*}
	\nabla u - \nabla g \| \nu \qquad x \in \Gamma
\end{align*}
Recall the abbreviation $p = \nabla u$

\begin{definition}[admissible points]
	Given $x_0 \in \Gamma$ with the normal vector $\nu(x_0)$, the triple $(p_0, z_0, z_0) \in \mathbb{R}^n \times \mathbb{R} \times \mathbb{R}^n$ is called admissible, if:
	\begin{itemize}
		\item $z_0 = g(z_0)$
		\item $p_0 - \nabla g(x_0) \| \nu(x_0)\texttt{(compatibility condtition)}$ on $\partial\Omega$
		\item $F(p_0, z_0, x_0) =0$
	\end{itemize}
	The compatibility condition is due to that $u_{x_i}$ must satisfy $u_{x_i}(x_0) = g_{x_i}(x_0)$ as well as $F(u_{x_i}(x_0), \dots) = 0$. 
\end{definition}

Assuming the boundary data could be smoothly reparametrized. 
\begin{theorem}[local existence]
	Suppose $(p_0, z_0, x_0)$ is admissible. And
	\begin{align*}
		\nu(x_0) \cdot F_p \neq 0 \texttt{ or }(F_{p_n} \neq 0) \quad (\texttt{noncharacteristic condition)}
	\end{align*}
	Let $u \in C^1$, then there exists $\epsilon > 0$ such that 
	\begin{align*}
		F(\nabla u, u, x) = 0 \quad &x \in \Omega \cap B_\epsilon(x_0)\\
		u(x) = g(x) \quad &x \in \Gamma \cap B_\epsilon (x_0)
	\end{align*} 
	admits an unique solution $u \in C^2(B_\epsilon(x_0) \cap \Omega)$
\end{theorem}
\begin{proof}\quad\\
\textbf{2. Claim:} $\exists \epsilon, \delta > 0 \forall x \in \Gamma \cap B(0, \epsilon) \exists ! q(y) \in B(p_0, \delta) \subset \mathbb{R}^n$:
\begin{align*}
	&q_j(y) = g_{x_j}(y)  \quad j = 1, \dots, n-1\\
	&F(q(y), g(y), y) = 0
\end{align*}
\textbf{1. straighten the boundary}.	[Evens p.103] via smooth transformation. Therefore we can assume w.o.l.g 
\begin{align*}
	\Gamma= \{x_n = 0\}
\end{align*} \\
\textbf{2. noncharacteristic boundary condition}: 
w.l.o.g. $\Omega = \mathbb{R}^{n - 1} \times \mathbb{R}_+$, $x_0 \in \mathbb{R}^{n-1}$,  $p_0 = (p_0^1, \dots, p_0^n) \subset \mathbb{R}^n$, $z_0 = g(x_0)$ is given by the initial data $g$. Therefore 
\begin{align*}
	(g_{x_1}(x_0), \dots, g_{x_{n-1}}(x_0)) = (p_0^1, \dots, p_0^{n-1})
\end{align*}
For the PDE to hold we also need $F(x, u, Du) = 0$, therefore we formulate the \textbf{compatibility condition:}
\begin{align*}
	&g_{x_i}(x_0) = p_0^i \\
	&F(p_0, z_0, x_0) = 0
\end{align*}
Now the problem arises that since $x_n \equiv 0$ on $\Gamma$ we cannot perturb $x_n$ to get the partial derivative. We need to verify that $\partial_n u$ is well-defined. To this end we observe:
\begin{align*}
	&g_{x_i}(y) = q^i(y)\\
	&F(g(y), q(y), y) = 0
\end{align*} 
Apply the implicit function theorem. Define the function:
\begin{align*}
	&G: \mathbb{R}^n \times \mathbb{R}^n \rightarrow \mathbb{R}^n \\
	&\left\{
	\begin{aligned}
		G_i(x, p) &= g_i(x) - p^i\\
		G_n(x, p) &= F(p, g(x), x) = 0 
	\end{aligned}
	\right.
\end{align*} 
We would like to prove that $G(x_0, p_0) = 0$, $x_0 \in \Gamma$ and $D_2 G(x_0, p_0) \in Lis(\mathbb{R}^n, \mathbb{R}^n)$. Observe
\begin{align*}
	D_2G(x_0, p_0 ) =
	\begin{bmatrix}
		1 & 0 & \dots &|0\\
		0 & 1 & \dots &|0\\
		\dots\\\
		F_{p_1} & F_{p_2} & \dots & F_{p_n} 
	\end{bmatrix}
\end{align*}
where $F_{p_n}$ is not $0$ by the assumption, thus $D_2G(x_0, p_0) \in Lis(\mathbb{R}^n, \mathbb{R}^n)$.   The implicit function theorem leads to $\delta, \epsilon > 0$ and a implicit function 
\begin{align*}
	q: B_{\epsilon}(x_0) &\rightarrow B_\delta(p_0),\quad \\
	x &\mapsto q(x)  \qquad q \in C^1(\mathbb{R}^n, \mathbb{R}^n)
\end{align*}
Therefore locally around $x_0$, the graph(vector field $p$) could be described implicitly by $q$. In particular $g_n(x)$. The projected characteristic is not parallel to the boundary. \\
\textbf{3. Local invertibility of $x$ :} claim: there exists an interval $I \subset \mathbb{R}, W \subset \Gamma \subset  \mathbb{R}^{n-1}, V \subset \mathbb{R}^n$, such that $\forall x \in V \exists s \in I, y \in \Gamma:$
\begin{align*}
	x = \textbf{x}(y, s) \quad (y \in \Gamma) 
\end{align*}
and the mappings $x \mapsto (s, y)$ are $C^2$. \\
Proof by the inverse function theorem: we have $\textbf{x}(x_0,0) = x_0$. We would like to prove $D \textbf{x}(x_0, 0) \in Lis(\mathbb{R}^n, \mathbb{R}^n)$. Indeed,
\begin{align*}
	\frac{\partial \textbf{x}_i }{\partial y_j}(x_0, 0) &= \delta_{ij} \quad j\in (1, \dots, n-1) \\
	\frac{\partial \textbf{x}_i }{\partial s}(x_0, 0) &= F_{p_i}(p_0, z_0, 0)
	\quad \texttt{(by the assumption of characteristic)}
\end{align*} 
Together we have:
\begin{align*}
D\textbf{x}(x_0, 0)= 
	\begin{bmatrix}
		1 & 0 & \dots &|F_{p_1}\\
		0 & 1 & \dots &|F_{p_2}\\
		\dots\\\
		0 & 0 & \dots & F_{p_n} 
	\end{bmatrix}
	\in \mathbb{R}^{n \times n}
\end{align*}
By the noncharacteristic condition we know the inverse function exists. \\
\textbf{4.} Define $u(x):= z(y, s)$ for $x = \gamma(y, s) \in B(0, \epsilon)$
\begin{align*}
	Q(x):= p(y, s)
\end{align*}
\textbf{claim: $F(Q(x), u(x), x) = 0$} for $x \in B(0, \epsilon)$. \\
Define:
\begin{align*}
	G(y,s):= F(p(y, s), z(y, s), \gamma(y, s)) \overset{!}{=} 0
\end{align*}
Observe:
\begin{align*}
	G_t &= F_p \cdot \dot{p} + F_z \dot{z} + F_\gamma \dot{\gamma}\\
		&= - F_p \cdot  F_x - F_z F_p \cdot p+ F_z F_p \cdot p + F_x \cdot F_p \\
		&= 0
\end{align*}
\textbf{5: $Du = Q$} in $B(0, \epsilon)$. $\forall k = 1, \dots, n$:
\begin{align*}
	u_{x_k} &= \sum_{j = 1}^{n - 1} z_{y_j} (y_j)_{x_k} + z_s s_{x_k}\\
	&= 
	F_p \cdot p\\
	&= 
	\dot{\gamma} \cdot p
\end{align*}

\end{proof}

limitation: we typically don't know the domain of the unique existence.


\subsection{The Cauchy-Kowaleski theorem}

\subsubsection{Normal derivatives to the Cauchy Hypersurface}
The Question now need to ask is: Given the boundary data given as the form above, which further derivative are we able to compute. \par 
Given a quasilinear PDE of the form
\begin{align*}
	\sum_{|\alpha| = k} a_\alpha(\nabla^{k -1}u, \dots, u, x) \partial^\alpha u + a_0(\nabla^{k -1} u, \dots, , x)  = 0
\end{align*}
where $\nabla^j$ is defined as all the partial derivative of total order less than $j$.\par 
The Cauchy problem equipped the general quasilinear PDE with the boundary data
\begin{align*}
	u = g, \quad \frac{\partial u}{\partial \nu} = g_1, \quad \dots \frac{\partial u}{\partial \nu^{k-1}} = g^{k -1} \qquad \forall x \in \Gamma
\end{align*}
where we define the $j-th$ \textbf{normal derivative} of $u$ at $x \in \Gamma$ as:
\begin{align*}
	\frac{\partial^j u}{\partial \nu^j}:= \sum_{|\alpha| = j} \nabla^\alpha_x u : \nu^j = \sum_{\alpha = j} \partial_x^\alpha u \cdot  \nu^\alpha
\end{align*}

\subsubsection{The flat boundary}
In the case of the flat boundary, we have $\Gamma = \{ x_l = 0\}$. Accordingly the boundary condition read:
\begin{align*}
	u(x_1, \dots, x_{l-1}, 0) = g(x_1, \dots, x_{l-1})\\
	\partial_{x_l} u = g_1\\
	\partial^2_{x_l} u = g_2\\
	\dots\\
	\partial^{k-1}_{x_l} u = g_{k-1}
\end{align*}
since all the derivative in the direction of $x_i, i \in {1, \dots, k-1}$ are given by $g_0$. To calculate the $k$-the derivative in the direction of $x_l$ we need to represent it with lower order derivative in other directions. 


\begin{definition}[Non-characteristic hypersurface PDE]
	Given the hypersurface $\Gamma = \{ x_l \} = 0$ and the boundary data and the PDE. If at $x \in \Gamma \cap U \subset \mathbb{R}^l$, the highest derivative could be expressed as the lower order derivative. Then it is non-characteristic on $\Gamma \cap U$.  
\end{definition}





In the spirit of method of characteristic we observe the higher-order implicitly defined non-linear PDE:
\begin{align*}
	& \partial_t^k u =  F\bigg(x, t, (\partial_x^\alpha \partial_t^j u)_{ |\alpha| + j \le k, j < k } \bigg) \qquad S \subset \mathbb{R}^n \times \mathbb{R}_+ \\
	& \partial_t^j  u (0, x) = g_j(x)   \quad 0 \le j < k
\end{align*}
where $S$ is assumed to be a smooth hypersurface \\
Recall the definition of the analytic function: the Taylor expansion of which converges absolutely. We may bound the domain on small ball $B(0, r)$ such that the multinomial could be bounded by a monomial. If the coefficients are uniformly bounded, we get a geometric series multiplied by the upper bound of the coefficient. 

\begin{lemma}
\end{lemma}

\begin{theorem}[Cauchy-Kowaleski theorem]
	Given the Cauchy problem above, assuming $g_j$ and the hypersurface $S \subset \mathbb{R}^{n + 1}$ are analytic. Then there exists a \textbf{unique analytic solution} to the Cauchy Problem.
\end{theorem}
\begin{proof}\quad\\
\textbf{1. Smooth transformation of the boundary:} since the hypersurface admits enough regularity we transform the boundary to the a real line. By the non-characteristic condition the characteristic is transversal, i.e,  the highest order partial derivative w.r.t. the boundary could be described by the initial Cauchy data. \\
\textbf{2. Transform the implicit higher-order PDE to a first order explict PDE.}: we do this by introducing the system: we replace $\partial_x^\alpha \partial_t^j u$ by $Y_{\alpha, j}$. 
\begin{align*} \textbf{CP(b):} \qquad
 	&\partial_t Y_{\alpha, j} = Y_{\alpha, j+1} \quad \forall |\alpha| + j \leq m - 1 \\
 	&\partial_t Y_{\alpha, j} = \partial_{x_{k(\alpha)}} Y_{\alpha - e_{k(\alpha)}, j} \quad \forall  |\alpha| + j = m, j \neq m \\
 	& \partial_t Y_{0, m} = \partial_t F + \sum_{|\alpha| + j \le m -1 }  \partial_{Y_{\alpha, j}} F \cdot  Y_{\alpha, j + 1} + \sum_{|\alpha| + j = m, j \neq m} \partial_{Y_{\alpha, j}} F \cdot  \partial_{x_{k(\alpha)}} Y_{ \alpha - e_{k(\alpha), j + 1}}
\end{align*}
with the initial condition
\begin{align*}
	 &Y_{\alpha, j}(0, x) = \partial^\alpha_x g_j(x) \quad \forall j \le m - 1, |\alpha|  + j \le m \\
	 &Y_{0, m}(0, x) = F \quad 
\end{align*}
\textbf{3. Transform into a 1st linear system with trivial initial condition and t-independent coefficients}:
\begin{align*}
\textbf{CP(c)} \quad
	&\partial_t \mathbb{Y} = \pmb{A}(x, \mathbb{Y}) \mathbb{Y} + \pmb{B}(x, \mathbb{Y}) \\
	&\mathbb{Y}(0, x) = 0
\end{align*}
From \textbf{CP(b)} we set $\mathbb{Y} = Y - g(x)$ and set $t = Y_m$ with $m = |B| + 1$. \\
\textbf{4. Transform $\mathbb{Y}$ into a Taylor series.}
df
\\
\textbf{5. Find a majorized serie.} Observe
\end{proof}









\section{Sobolev Space and General Solution Concept}
\subsection{The Divergence Theorem}
\subsubsection{The Lipschitz domain}
In the following we observe the domain $\Omega \subset \mathbb{R}^n$ and the function defined on it. We often observe the open domain $\Omega$ since the definition of derivative on $x \in \partial \Omega$ is not clear. 
\begin{definition}[Lipschitz domain]
A domain is an open, connected and non-empty subset of $\mathbb{R}^n$. We say $\Omega \subset \mathbb{R}^n$ is a Lipschitz domain if there exists $\forall x_0 \in \partial\Omega \exists w \in \mathbb{R}^n, \epsilon > 0$, function $g$ such that
\begin{enumerate}
	\item \textbf{coordinate system:} $H:= \{ y \in \mathbb{R}^n| \langle y - x_0, w \rangle = 0\}$ is a 1-dimensional axis. 
	\item \textbf{parametrisation of the boundary:} the boundary $\partial\Omega$ can be locally parametrised by a Lipschitz function $g:H \rightarrow \mathbb{R}^{n -1}$ . 
\end{enumerate}
\end{definition}
\begin{remark}
\textbf{1. choice of $w$:} Ideally normal vector but could also be something else. \\
\textbf{2. The partition of unity around boundary:}
 For bounded $\Omega$ we know $\partial \Omega$ is closed and bounded hence compact, therefore we can find a finite open cover of $\partial\Omega$:
\begin{align*}
	\partial\Omega \subset \cup_{j = 1}^N U_j, \quad U_j \texttt{ cylinder}.
\end{align*}	
Adding a subset $U_0$ with positive distance to $\partial\Omega$(writing as $U_0 \subset  \subset \Omega$). We have 
\begin{align*}
	\overline{\Omega} = \cup_{j = 1}^N U_j.
\end{align*}
We can therefore define the partition of unity $\eta_j \in C^\infty(U_j, \mathbb{R})$ with
\begin{align*}
	\sum_{j = 0}^N \eta_j(x) = 1 \qquad \forall x \in \bar{\Omega}
\end{align*}
\end{remark}

\subsubsection{The Boundary Integral}
Question: Given $f \in L^p(\Omega)$, how can we define 
\begin{align*}
	\int_\Gamma f.
\end{align*}
By the last remark we have
\begin{align*}
	\int_\Gamma f  = \sum_{j = 0}^M \int_{U_j\cap \Gamma} f\eta_j.
\end{align*}
Given the coordinate system defined by the $H_{x_0, w}$ we may locally define the graph
\begin{align*}
	\Gamma := \texttt{graph}(g) = \{(f(g(y), y), y \in H\}
\end{align*}  
Since $\Omega$ is Lipschitz domain, by the theorem of Rademacher $\nabla g$ is well-defined a.e.,  we define the \textbf{boundary integral:}
\begin{equation}
	\int_\Gamma f d\mathcal{H}^{n -1} = \int_\Gamma f(g(y), y) \sqrt{1 + \nabla g(y)^2} \mathcal{L}^{n -1}. 
\end{equation}


\subsubsection{The Divergence Theorem}
\begin{definition}[the normal vector]
	Given $\Omega \subset \mathbb{R}^n$ a bounded and Lipschitz domain parametrised by $\Gamma = \{(y, g(y)): y \in \mathbb{R}^n \}$. We define the \textbf{normal vector} at $x_0 \in \partial\Omega$ as 
	\begin{equation}
		\nu(x_0) = \frac{(- \nabla g(x_0), 1) }{\sqrt{|\nabla g(x_0)|^2 + 1}.}
	\end{equation}
\end{definition}

\begin{theorem}[Divergence Theorem]
	Given $\Omega$ Lipschitz domain, $u \in C^1_c(\mathbb{R}^n)$,  we have
	\begin{equation}
		\int_\Omega \partial_i u  = \int_{\partial\Omega} u(e_i \cdot \nu)
	\end{equation}
	Integration by parts we have
	\begin{equation}
		\int_\Omega \partial_i(uw) = \int_{\partial\Omega} (uw)(e_i \cdot \nu)
	\end{equation}
\end{theorem}
\begin{remark}
\textbf{Notation}: We write 
\begin{equation}
	\nabla u \cdot \nu := \partial_\nu u. 
\end{equation}	
\end{remark}




\subsection{Weak derivative, Weak Solutions and Sobolev space}
\begin{definition}[weak(distributional) partial derivative of $L^p$ functions]
	A function $w \in L^p(\Omega)$ is called a weak the weak derivative with respect to $x_j$ to $u \in L^p(\Omega)$ if :
	\begin{equation*}
		\forall \varphi \in C^\infty_C(\Omega): \int_\Omega w \varphi d\mu = -\int_\Omega u \partial_{x_j} \varphi d\mu  
	\end{equation*}
	more generally, $w \in L^p(\Omega)$ is $\alpha$ order derivative of $u$, if
	\begin{equation*}
		\forall \varphi \in C^\infty_C(\Omega): \int w \varphi d\mu = |-1|^\alpha \int u D^\alpha \varphi d\mu
	\end{equation*}
\end{definition}

\begin{remark}.
\begin{enumerate}
	\item All terms are well-defined
	\item For $u \in C^1(\Omega)$, $w = \partial_{x_j} u$
	\item The weak derivative is uniquely defined
\end{enumerate}	
\end{remark}

Later on we would see the implication behind this definition. When we are given a Hilbert space with infinite dimension, Balzano-Weierstrass is not applicable and for a bounded sequence to have a convergent subsequence(compactness) with respect to the norm induced by the scalar product. But later on, Banach-Alaoglu suggests in the dual of a Banach space, every bounded sequence has a weak* convergent subsequence.

\begin{definition}[Sobolev space]
For $\Omega \subset \mathbb{R}^d$, $\alpha \in \mathbb{N}$ we define
\begin{equation*}
	W^{k, p}:= \{u \in L^p(\Omega): \forall |\alpha| \leq k: D^\alpha u \in L^p(\Omega)\}
\end{equation*}
as the Sobolev space with integrability power $p$ and differentiation index $\alpha$.(Here "$\subset \subset$" mean compact containment)
\begin{equation*}
	W^{k, p}_{loc}:= \{ u \in L^p(\Omega)| \forall V \subset \subset \Omega \forall |\alpha| \leq k  :D^\alpha u \in L^p(V) \}
\end{equation*} 
To control the boundary value, we define
\begin{equation*}
	W_0^{p, k}:= \overline{ C_C^\infty(\Omega)}^{W_{k, p}}
\end{equation*}
	The function of $C^\infty_0(\Omega)$ that are certainly 0 on the boundary $\partial \Omega$.
\end{definition}

\begin{example}
By solving the Dirichlet problem we're going to encounter the space:
\begin{align*}
	H^1_0:= \overline{ C_0^\infty(\Omega)}^{\|\cdot\|_{H_1}}
\end{align*}
where the $\|\|_{H_1}$ is defined as:
\begin{align*}
	\|u\|_{H_1}^2 := \int |u|^2 + |\nabla u \cdot \nabla u|  d\mu
\end{align*}
\end{example}

\begin{example}
Consider $B_R(0) \subset \mathbb{R}^d$, we define:
\begin{equation*}
u(x) = \left\{ 
	\begin{aligned}
		&\frac{\pi^2}{17} \quad &x = 0\\
		&|x|^\alpha \quad \quad &n \neq 0
	\end{aligned}
	\right.
\end{equation*}	
We are interested in that for which $p$ we have $u \in W^{1, p}$. 
\begin{align*}
	\|u\|_p^p &= \int_{B_R(0)} |u(x)|^p dx \\
	&= \int_0^R r^{\alpha p} dr \int_{\partial B(0, r)} dS  \\
	&= \omega_{d-1} \int_0^R r^{\alpha p} r^{d-1}dr
\end{align*}
By using the spherical integral
\begin{equation}
\begin{aligned}
	\int_{\partial B(0, r)} 1 dS &= \omega^{n-1} r^{n-1}
\end{aligned}
\end{equation}
where $\omega^n$ is the \textbf{surface integral of an n-sphare}
\begin{itemize}
	\item The condition for $u \in L^p$ is 
\begin{equation*}
	\alpha p + d-1 > -1 \Leftrightarrow \alpha p > -d
\end{equation*}
	\item $\forall x \neq 0$ we can define the classical partial derivative of $u$: $w_j = \alpha |x|^{\alpha - 2} x_j$. Notice the original $u$ has a singularity, we avoid this by applying the DCT. We have: $\forall \varphi \in C_c^\infty(B_R(0)$:
	\begin{align*}
		\int_{B_R(0)} w_j \varphi d\lambda &= \lim_{\epsilon \rightarrow 0} \int_{B_R(0)\slash B_\epsilon (0)} w_j \varphi d\lambda \\
		&= \lim_{\epsilon \rightarrow 0}- \int_{B_R(0)\slash B_\epsilon (0)} u \partial_j \varphi  d\lambda
	\end{align*} 
	Given the condition that $w_j$, $u$ are in $L^p$. Notice that $u \in L^p$ implies that $w_j \in L^p$.
	\item All we need to do now is to verify that $\forall x \in B_R(0) \slash B_\epsilon(0)$, $w_j$ satisfies the definition for weak partial derivative. Indeed, applying integration by parts we have
	\begin{align*}
		\int_{B_R(0)\slash B_\epsilon (0)} w_j \varphi d\lambda &= \int_{\partial B_R(0) \cup \partial B_\epsilon(0)} u \varphi \cdot \nu  d\lambda - \int_{B_R(0)\slash B_\epsilon (0)} \partial_j \varphi \omega d\lambda\\
		\textbf{(1)}&=  - \int_{B_R(0)\slash B_\epsilon (0)} \partial_j \varphi \omega d\lambda
	\end{align*}
	(1) follows from: for $\partial B_R(0)$ we conclude from the zero-boundary condition of $\varphi$. For $\partial B_\epsilon(0)$ we have:
	\begin{align*}
		\int_{\partial B_\epsilon(0)}u \varphi \cdot \nu d \lambda &\leq \|\varphi\|_\infty \omega_{d-1} \int_0^\epsilon |r|^{\alpha + d -1} dr \rightarrow 0
	\end{align*}
\end{itemize}
\end{example}


\subsubsection{Weak Solutions}
\begin{definition}[Types of solutions of Wave Equations]
	Given a general PDE of the form 
	\begin{align*}
		F(u, \nabla u, \dots,  D^\alpha u)
	\end{align*}
	we define
	\begin{itemize}
		\item \textbf{The classic solution}: $u$ is a classic solution if $u \in C^{|\alpha|}(\Omega)$
		\item \textbf{The strong solution}: $u$ is a strong solution if $u \in W^{|\alpha|, p}$
		\item \textbf{weak solution}: $u$ is a weak solution if
		\begin{equation*}
		\textbf{(WF)} \qquad
		 	\forall \varphi \in C_c^\infty(\Omega):
		 	\int_\mathbb{R} \int_{\mathbb{R}^d} - \partial_t \varphi \partial_t(\rho u)  + div(\varphi) (A \nabla u) - f \varphi dxdt	= 0
		\end{equation*}
	\end{itemize}
\end{definition}
\begin{remark}.
\begin{itemize}
	\item the types of solutions are ordered.
	\item The test function in general doesn't assume $0$ boundaries. 
\end{itemize}	
\end{remark}

\begin{theorem}
	If $u$ is a weak solution of the WE, then:
	\begin{itemize}
		\item if $u \in W^{2, p}_{loc}(\Omega)$  Then $u$ is a strong solution.
		\item if $u \in C^2(\mathbb{R}^d, \mathbb{R})$, then $u$ is a classic solution.
	\end{itemize} 
\end{theorem}
\begin{proof}\quad\\
\textbf{1)}By the definition of weak solution we have:
	\begin{align*}
		\forall \varphi \in C^\infty(\Omega):& \int_{\mathbb{R}^d} \int_\mathbb{R} - \partial_t \varphi \partial_t(\rho u) + div(\varphi)(A \nabla u) - f \varphi dx dt = 0\\
		(u \in W^{2, p}) \quad \Rightarrow \forall \varphi \in C^\infty(\Omega): & \int_{\mathbb{R}^d} \int_\mathbb{R} \varphi \{\partial^2_t (\rho u)   -  div(A \nabla u) -  f \}dxdt = 0
	\end{align*}
	By \textbf{the fundamental lemma of calculus variations} we know $\partial_t^2(\rho u) - div(A \nabla u) = f$ a.e. Hence $u$ satisfies the strong form. \\
	Analog for $2)$ we get the same results. 
\end{proof}
\begin{remark}
Observe the difference in regularities of strong solution and weak solution. The weak solution need only to be in $W^{1, p}(\Omega)$. 
\end{remark}

\begin{definition}[distribution(super weak) solution]
\end{definition}


\subsection{Embedding in Sobolev spaces}
\begin{definition}[Embedding]
	Given 2 Banach spaces $(X, \|\cdot\|_x), (Y, \|\cdot\|_Y)$, we say $i: X \rightarrow Y$ is an embedding if $i$ is \textbf{injective} and \textbf{linear}. Furthermore the embedding is continuous, if $i$ is continuous(denoted by $X\hookrightarrow Y$), i.e 
	\begin{align*}
		\exists C \forall u \in X: \|i(u)\|_Y \leq C \|u\|_X
	\end{align*}
	\textbf{compact embedding: }
	An embedding $X \hookrightarrow Y$ is \textbf{compact} if every bounded sequence $u_n$ in $X$ admits an strongly convergent subsequence in $Y$. 
\end{definition}

\begin{theorem}[$|21|$Rellich's compact embedding]
Given an open and bounded domain $\Omega \subset \mathbb{R}^n$, we have 
\begin{align*}
	H_0^1(\Omega) \overset{c}{\hookrightarrow} L^2(\Omega)
\end{align*}
	
\begin{proof}
	Given $f_k \in H_0^1(\Omega)$ a bounded sequence, we would like to show it admits a convergent subsequence in $L^2(\Omega)$. By the boundedness of $f_k$ we know by Banach-Anauglu it admits a weak convergent subsequence, w.o.l.g. assuming the $f_k$ being weak convergent. We prove it's also $L^2$ convergent. Define $\mathcal{F}$ as the Fourier transform, by Parseval:
	\begin{align*}
		\|f_j - f_k\|_{L^2(\Omega)} &= \|\mathcal{F}f_j - \mathcal{F}f_k\|_{L^2(\mathbb{R}^n)}
		= \int_{\xi \in \mathbb{R}^n} |\mathcal{F}f_j - \mathcal{F} f_k|^2 d\xi
	\end{align*}
	Observe
	\begin{align*}
		\int_{|\xi| < R} |\mathcal{F}f_j - \mathcal{F} f_k|^2 d\xi + \int_{|\xi| \ge R} |\mathcal{F}f_j - \mathcal{F} f_k|^2 d\xi
	\end{align*}
	For the second part, since the domain is unbounded, the $L^2$ convergence doesn't imply the $L^1$ convergence. We leverage that the Fourier transform of the derivative is the Fourier transform multiplied by a monoid, we have:
	\begin{align*}
		\|\xi_l \mathcal{F} f\|_{L^2(\mathbb{R})} = \bigg\|\mathcal{F} \bigg(\partial_{x_l} f\bigg)\bigg\|_{L^2(\mathbb{R})} = \|\partial_{x_l} f\|_{L^2(\mathbb{R})} < \|f\|_{H^1(\mathbb{R}^1)}
	\end{align*}
	\begin{align*}
		\int_{|\xi| > R} |\mathcal{F}f_j - \mathcal{F}f_k|^2 d\xi 
		& = \int_{|\xi| >R} \bigg| \underbrace{ \sum_{1 \le i \le n} \xi^i (\mathcal{F}f_j - \mathcal{F} f_k)}_{ \le \sum_{ 1 \le i \le n } (\mathcal{F} [\partial_{x_i} f] - \mathcal{F}[\partial_{x_j} f])  }  \times \underbrace{ \frac{1}{\sum \xi^i}}_{\le R} \bigg|^2 d\xi \\
		&\le  \frac{1}{R^2} \|f_j - f_k\|_{H^1} \\
		&\rightarrow 0
	\end{align*}
	For the first part, since the \textbf{domain is bounded} and the sequence is $H^1$ bounded therefore $L^2$ bounded, we know the sequence is $L^1$ bounded. Sending $R \rightarrow 0$ and apply the dominated convergence we know the 2nd part convergence to 0.
\end{proof}
\end{theorem}

\begin{remark}
If $\Omega$ is bounded we can apply Mazur and the $H^1$ weak convergent subsequence is (up to a subsequence) strongly convergent and therefore $L^2$ convergent. 
\end{remark}

\begin{theorem}[The Sobolev Embedding]
	Given a bounded domain $\Omega \subset \mathbb{R}^d$ with Lipschitz boundary $\partial\Omega \in C^{0, 1}(\partial\Omega)$ , we have
	\begin{align*}
		W^{1, p}(\Omega) \hookrightarrow L^{ps}(\Omega)
	\end{align*} 
	provided $p$ satisfies
	\begin{equation}
	ps = \left\{
	\begin{aligned}
		 \frac{dp}{d - 1} &\texttt{ if } p < d\\
		\texttt{large but finite} \quad &\texttt{ if } p = d\\
		+ \infty \quad  &\texttt{ if } p > d
	\end{aligned}
	\right.
	\end{equation}
\end{theorem}

\begin{definition}[The fractional Sobolev spaces]
	$\forall \sigma \in (0, 1)$, $\forall p \in [1, \infty)$, we define 
	\begin{align*}
		W^{\sigma, p}(\Omega):= \bigg\{ u \in L^p(\Omega): \frac{|u(x) - u(y)|}{|x - y|^{\frac{p}{n} + s}} \in L^p(\Omega \times \Omega) \bigg\}
	\end{align*}
\end{definition}

\begin{theorem}
Given a regular domain $\Omega$. We define the fractional Sobolev space as:
\begin{align*}
	W^{\sigma, p}(\Omega):= \{ u \in L^p(\Omega) : \|u\|_p + |u|_{\sigma, p} < \infty \}
\end{align*}
given $0 < \sigma < \sigma' < 1$,  show that:
\begin{align*}
	W^{1, p}(\Omega) \hookrightarrow W^{\sigma, p} \hookrightarrow L^p(\Omega)
\end{align*}
\end{theorem}
\begin{proof}
	The goal is to prove that, there exists $i$:
	\begin{align*}
		\forall u \in W^{\sigma, p}(\Omega): \|i(u)\|_{\sigma, p} \leq C
		\|u\|_{1, p}
  	\end{align*}
  	we take the canonical embedding, and observe
  	\begin{align*}
  		\|u\|_{\sigma, p}^p &= \|u\|_p^t + |u|_{\sigma, p}^p\\
  		&= \|u\|_p + \int_x \int_y  \frac{|u(x) - u(y)|^p}{|x - y|^{n + \sigma p}} dx dy\\
  		&= \|u\|_p^p + \int_{\Omega}\int_{|x-y| < 1} \frac{|u(x) - u(y)|^p}{|x - y|^{n +\sigma p}} dx dy + \int_{\Omega} \int_{|x-y| > 1}\frac{|u(x) - u(y)|^p}{|x - y|^{n +\sigma p}} dx dy
  	\end{align*}
  	define the first integral as $A$ and the second as $B$, 
  	\begin{align*}
  	 B & \overset{(1)}{\leq} \int_\Omega \int_{|x -y| > 1} 2^{p-1} (|u(x)|^p + |u(y)|^p) \cdot \frac{1}{|x - y|^{n + \sigma p}} dxdy\\
  	   &\leq 2^p \int_\Omega \int_{|z| > 1} |u(x)|^p \frac{1}{|z|^{n + \sigma p}}dx dz\\
  	   &\overset{(2)}{=} 2^p C_p \int_{\Omega} |u(x)|^p dx  \\
  		   &\leq \tilde{C_p} \|u\|_p^p
  	\end{align*}
  	(2) follows from that $\frac{1}{|z|^m}$ is integrable in $\mathbb{R}^n$ as long as $m > n$.\\
  	For the part $A$, under the regularity condition, we have $\tilde{u}: \mathbb{R}^n \rightarrow \mathbb{R}$:
  	\begin{align*}
  		\|\tilde{u}\|_{W^{1, p}(\mathbb{R}^n)} \leq C\|u\|_{W^{1, p}(\Omega)}
  	\end{align*}
  	\begin{align*}
  		A &= \int_\Omega \int_{|z| < 1} |u(x) - u(x + z)|^p \frac{1}{|z|^{n + \sigma p}} dx dz\\
  (Jensen)&\leq  \int_\Omega \int_{|z| < 1} \int_0^1 |\nabla u(x + tz)|^p 
  			\frac{1}{|z|^{n + \sigma p}} dt dx dz\\
  		& \leq \int_{\mathbb{R}^n} \int_{|z| < 1} \int_0^1 |\nabla \tilde{u}(x + tz)|^p 
  			\frac{1}{|z|^{n + \sigma p}} dt dx dz \\
  		& = \int_{|z| < 1} \int_0^1 \|\nabla \tilde{u}\|_p^p \frac{1}{|z|^{n + \sigma p}} dz dt \\
  		& = C_1(n, \sigma, p) \|\nabla \tilde{u}\|^p_p\\
  		& \leq C_2(n, \sigma, p) \|u\|_{W^{1,p}(\Omega)}^p
  	\end{align*}
\end{proof}
$\quad$\\
\textbf{Question:}
For which $\sigma$ are the Heaviside functions in $W^{\sigma, p}(-1, 1)$


\subsection{Embedding of the trace}
\begin{theorem}[Existence of traces]
	Given a d-dimensional domain $\Omega$ with Lipschitz boundary $\partial \Omega$. $\forall p \in [1, \infty] \forall u \in W^{1, p}:$ we have a unique mapping $\gamma$
	\begin{equation*}
	\left\{
	\begin{aligned}
		W^{1, p}(\bar{\Omega}) &\rightarrow L^p\\
		u &\mapsto \gamma(u) 
	\end{aligned}
	\right.	
	\end{equation*}
	such that $\gamma$ extends the classical trace mapping from $C^1_c(\bar{\Omega})$ to $C(\partial\Omega)$.
\end{theorem}
\begin{proof}(Ben Schweizer, 3.20)
	
\end{proof}


\begin{exercise}[$|12|$Preparing for trace].
$0<\alpha<1$, $p \in [0, \infty]$, $\Omega = \mathbb{R}^n$, find the exact relation of $p$ and $\alpha$ such that embedding:
\begin{align*}
	W^{1, p}(\Omega) \hookrightarrow C^{\alpha}(\Omega)
\end{align*}
exists
\begin{proof}
\begin{equation}
\int_{B(x, r)} |u(y) - u(x)|dy \leq C \int_{B(x, r)} \frac{|Du(y)|}{|y-x|^{n-1}}dy	
\end{equation}
$\forall x, y \in \Omega \exists \omega \in \partial B_0(1), t \in [0, \infty] $:
\begin{align*}
	u(x) - u(y) &= u(x) - u(x + t\omega) \\
	&= \int_0^t \frac{d}{ds} u(x + s\omega) ds \\
	&= \int_0^t D u(x + s\omega) \cdot \omega ds\\
	&\leq \int_0^t Du (x + s\omega) ds
\end{align*}
Plug in we get:
\begin{align*}
	\int_{B_x(R)}| u(x) - u(y)|dy &= S_{n -1}(u) \int_0^r R^{n-1} dR\\
	&=  \frac{r^n}{n} \int_{\partial B_x(1)} |u(y) - u(x)| dy \\
	&= \frac{r^n}{n} \int_{\partial B_x(1)} |u(x + s\omega) - u(x)|d\omega \\
	&= \frac{r^n}{n} \int_{\partial B_x(1)} \int_0^s Du(x + t\omega)dt d\omega \\
	(y = x + t\omega, t = |x -y|)(1)
	&= \frac{r^n}{n}  \int_{\partial B_x(s)} Du(y)|x - y|^{-(n-1)}dy \\
	&\leq C \int_{B_x(r)} \frac{Du(y)}{|x - y|^{n-1}}dy
\end{align*}
From above we get(see exercise)
\begin{equation}
|u(x)| \leq C \|u\|_{W^{1,p}(\Omega)}	
\end{equation}
therefore
\begin{align*}
	\|u\|_\infty = \sup_x |u(x)| \leq C \|u\|_{W^{1, p}(\Omega)}
\end{align*}
Analog we can prove the result of the semi-norm $|u|_{C^{\alpha}}$(see exercise). In the end, we can show:
\begin{align*}
	\alpha = 1 - \frac{n}{p}
\end{align*} 
\end{proof}
\end{exercise}

\begin{theorem}
	Given a Lipschitz domain, for $1 \le p \le \infty$, we have
	\begin{align*}
		&u_n \rightharpoonup u \texttt{ in } L^p(\Omega)\\
		\Rightarrow &u_n \rightarrow u \texttt{ in } L^p(\partial\Omega)
	\end{align*}
\end{theorem}


\subsection{ $W^{1, ^\infty}$ and Lipschitz functions}



\subsection{The Difference Quotient}
Before proving the regularity condition, we first define the tool: the \textbf{difference quotient.}

\begin{definition}[Difference Quotient]
	Given $e_j$ the shift along the $j-th$ coordinate, we define the difference quotient:
	\begin{align*}
		\partial^h_j u := \frac{T_{he_j} u - u}{h} &\in L^p(\mathbb{R})\\
		\text{where}: T_{he_j}: L^p(\mathbb{R}) &\rightarrow L^p(\mathbb{R}), \\
		u &\mapsto u(\cdot + he_j)
	\end{align*}
\end{definition}

\begin{proposition}
	Given $\Omega = \mathbb{R}^d$, $u \in L^p(\Omega) $, $p \in (0, \infty]$. If for any $\Omega' \subset \subset \Omega, h < dist(\Omega', \partial\Omega)$,  we have $D_i^h u \in L^p(\Omega')$ and  :
	\begin{align*} 
		\|D_i^h u \|_{L^p(\Omega')} \le  C_i  < \infty
	\end{align*}
	Then we have $u \in W^{1, p}(\Omega)$, we can estimate the $L^p$ norm of the partial derivative over $\Omega'$ by the same constant. 
	\begin{align*}
		\|u_{x_i}\|_{L^p(\Omega)} \le C_i
	\end{align*}
	Conversely, if $u \in W^{1, p}(\Omega)$, then 
	\begin{align*}
		\|D^h_k u \|_{L^p(\Omega')} \le \|u_k\|_{L^p(\Omega)} \quad \forall \Omega'\subset \subset \Omega
	\end{align*}
\end{proposition}
\begin{proof}
	By density we only need to show the argument for all $u \in C^\infty_C(\Omega)$.\\
	$\Rightarrow"$: 
	\begin{align*}
		\|\partial_j^h u \|_{L^p}^p &= \int_\Omega |\partial_j^h u|^p dx \\
		(1)&= \int_\Omega | \int_0^1 \partial_j u(x + \theta h e_j)d\theta|^p dx\\
		(2)&\leq \int_0^1\int_\Omega |\partial_j u(x + \theta h e_j)|^p dx d\theta\\
		(3)&= \|\partial_j u\|_p
	\end{align*}
	where (1) follows from fundamental theorem of integral and differentiation, (2) follows from Jensen, and (3) follows from the translation invariance of the norm.\\
	$\Leftarrow":$\\
	We would like to show that there exists a $w_j$, such that:
	\begin{align*}
		\forall \varphi \in C^\infty_C(\Omega): -\int_\Omega w_j \varphi dx = \int_\Omega u \partial_j \varphi dx
	\end{align*}
	Now we define:
	\begin{align*}
		I_k(u) := \int_\Omega \partial_j^{h_k} u \varphi dx
	\end{align*}
	Since the difference quotient is bounded, by Banach-Alaoglu we know $\exists w_j: I_k(u) \overset{weakly}{\rightharpoonup} \int_\Omega w_j \varphi dx,  \forall \varphi \in C_C^\infty$. Take $\frac{1}{h_k} \rightarrow 0$ we get:
	\begin{align*}
		I^k(u) &= \int_\Omega \frac{u(x + h_ke_j) - u(x)}{h_k} \varphi(x) dx \\
		&= \frac{1}{h_k} (\int_\Omega u(x + h_k e_j)\varphi(x)dx - \int_\Omega u(x) \varphi(x)dx)\\
		&= \frac{1}{h_k}(\int_\Omega u(\tilde{x}) \varphi(\tilde{x} - h_k e_j)dx - \int_\Omega u(\tilde{x})\varphi(\tilde{x})dx)\\
		&=  \int_\Omega u(\tilde{x}) \frac{\varphi(\tilde{x} - h_k e_j) - \varphi(\tilde{x})}{h_k} d\tilde{x}    \\
		&\overset{h_k \rightarrow 0}{\rightarrow} -\int_\Omega u(x) \partial_j \varphi(x)dx 
	\end{align*}
	we get:
	\begin{align*}
		\lim_k I_k(u) = \int_\Omega w_j \varphi dx = -\int_\Omega u \partial_j \varphi dx 
	\end{align*} 
\end{proof}

\begin{proposition}
We have the rule of calculation for the difference quotient, given $\Omega = \mathbb{R}^d$:
	
\end{proposition}

\begin{enumerate}
\item
	$\textbf{transformation rule} :
	  \int_\Omega \partial^h_j u v dx = -\int_\Omega u \partial_j^{-h}vdx$ 
\item $\textbf{product rule} :
	 \partial_j^h(au) = a \partial_j^h u  + \partial_j^h a T_j^h u$
\end{enumerate}	



\subsection{Bochner space}
\subsubsection{The Bochner Integral}
\begin{definition}[strongly measurable function]
	Given $X$ a Banach space, we define the \textbf{simple function} as the functions that take finite many value
	\begin{align*}
		f(t) = \sum_{i = 1}^j a_i \mathbbm{1}_{t \in B_i}(t), \quad a_i \in X, B_i \subset [a, b]
	\end{align*}
	Given $-\infty \le a \le b \le \infty$ and the function
	\begin{align*}
		f: (a, b) \rightarrow X
	\end{align*}
	we say the function is \textbf{strongly measurable} if it's the pointwise limit of the simple function, i.e., there exists $(u_n)$ sequence of simple functions such that
	\begin{align*}
		u_n(t) \rightarrow u(t) \quad \texttt{for } t \in [a, b], a.e.
	\end{align*}
\end{definition}

\begin{theorem}[Pettis]
	If the Banach space $X$ is separable, if
		\begin{align*}
		\forall \alpha \in \mathbb{R} \forall \lambda \in X^* : \{t \in \mathbb{R}: \langle u(t), \lambda \rangle \le \alpha \}  \texttt{ is measurable}
	\end{align*}
	Then the function $u$ is strongly measurable. We call the condition \textbf{weak measurable}.
\end{theorem}

\begin{definition}[Bochner Integral]
	Given a Banach space, the \textbf{Bochner-integral} is defined as 
	\begin{align*}
		\int_a^b f(t) dt = \int_0^T \sum_{i = 1}^N \mathbbm{1}_{B_i}(t) x_i dt = \sum_{i = 1}^N m(B_i) x_i \in X
	\end{align*}
	Accordingly we define the \textbf{integrability }by taking the norm of $X$ in the integrand:
	\begin{align*}
		f \in L^p((a, b), X) \Leftrightarrow
		\|f\|_{L^p((a, b), X)}^p 
		= \int_a^b \|f\|_X^p dt < \infty
	\end{align*}
\end{definition}

\begin{theorem}[Bochner]
	Given $X$ a Bochner space and $f: [0, T] \rightarrow X$ strongly measurable, then:
	\begin{enumerate}
		\item \textbf{integrability equivalent to integrability w.r.t. the norm function.}
		\begin{align*}
			f \in L^1((a, b), X) \Leftrightarrow \bigg(t \mapsto \|f(t)\|_X \bigg)\in L^1((a,b), \mathbb{R})
		\end{align*}
		there is a bijection between the $L^1$ function of Bochner space and $L^1$ function in the real space. 
		\item \textbf{linearity and sublinearity}:
		\begin{align*}
			\|\int_0^T f(t) dt \|_X  &\le \|f\|_{L^1((a, b), X)}	\\
			\int_0^T \langle \Lambda, u \rangle dt &= \langle \Lambda, \int_0^T u dt \rangle \quad \forall \Lambda \in X^*
		\end{align*}
	\end{enumerate}
\end{theorem}
\begin{proof}
\quad\\
\textbf{1).} Since the function is summable we can apply the dominated convergence theorem to show the result. 
\quad\\
\textbf{2).}
\begin{align*}
	\int_0^T \langle \Lambda, u \rangle dt \overset{Pettis}{\longleftarrow} \int_0^T \langle \Lambda, u_k \rangle dt \overset{linearity}{=} \langle \Lambda, \int_0^T u_k dt \rangle \overset{Fubini}{\longrightarrow} \langle \Lambda, \int_0^T u dt \rangle   
\end{align*}
\end{proof}

\begin{definition}[Bochner space]
For $1 \leq p \leq \infty$, the Bochner space is defined as:
\begin{align*}
	L^p(0, T; \Omega) = \{ u: [0, T]\rightarrow X : \|u\|_{L^p(0, T;X)} < \infty \}
\end{align*}
where the norm is defined as :
\begin{equation}
\|u\|_{L(0, T;X)}:=
\left\{
\begin{aligned}
	p &\neq \infty: (\int_0^T \|u(t)\|_X^p dx)^\frac{1}{p}\\
	p &= \infty : \esssup_{t \in [0, T]} \|u(t)\|_X
\end{aligned}
\right.	
\end{equation}
\end{definition}


\subsubsection{Properties}
\begin{theorem}[Properties of Bochner space]
	Given $X$ a reflexiv Banach space and $-\infty < a < b < \infty$, $\frac{1}{p} + \frac{1}{q} = 1$, $p, q < \infty$. 
	\begin{enumerate}
		\item \textbf{Mollification:} Define $\eta_\epsilon: \mathbb{R} \rightarrow \mathbb{R}$ the standard mollifier with $\eta_\epsilon \overset{\epsilon \rightarrow 0}{\rightarrow} \delta_0$. And $u \in L^p(0, T; X)$, we have:
		\begin{align*}
			C^\infty((0, T), X) \ni u * \eta_\epsilon \longrightarrow u \texttt{ in } L^p(0, T, X)
		\end{align*}
		where
		\begin{align*}
			u* \eta_\epsilon := \int_\mathbb{R} u(s) \eta( t - s ) ds
		\end{align*}
		\item \textbf{the Dual Space:}
		\begin{align*}
			L^p([0, T], X) = L^q([0, T], X^*)
		\end{align*}
	\end{enumerate}
\end{theorem}
\begin{proof}
	\quad\\
	\textbf{1).}\\
	\textbf{2).} 
\end{proof}

\begin{lemma}
	For $u \in W^{1, p}(0, T; X)$, we have the fundamental formula: $ 0\le t_1 < t_1 \le T:$
	\begin{align*}
		u(t_2) = u(t_1) + \int_{t_1}^{t_2} \partial_t u(s) ds\quad \texttt{for a.e.}t_1, t_2. 
	\end{align*}
\end{lemma}
\begin{remark}
$u(t_1), u(t_2)$ are only defined \textbf{a.e.}.	 This formula holds true only for a $t_1, t_2$-dependant null set.
\end{remark}

\begin{proof}\quad\\
	\textbf{1). }Observe the mollifier of $u$: $u_\epsilon$, we have
	\begin{align*}
		\partial_tu_\epsilon(t) dt 
	=  \int_{\mathbb{R}} \partial_t \eta_\epsilon (t - \tau) u(\tau) d\tau dt 
	= \int_{\mathbb{R}} \partial_s u (t -s) \eta_\epsilon (s)ds dt  = \int_{t_1}^{t_2} \partial_t u * \eta_\epsilon (t) dt 
	\end{align*}
	\textbf{2).}
	We verify the equality in Bochner space:
	\begin{align*}
	\textbf{claim: }  u_\epsilon(t_2) - u_\epsilon(t_1) = \int_{t_1}^{t_2} \partial_t u * \eta_\epsilon (t) dt
	\end{align*}
	By taking all $\lambda \in X^*$ and apply the theorem of Bochner to interchange the duality bracket and the integral, we verify the result.\\
	\textbf{3). }Taking the limit $\epsilon \rightarrow 0$, by $L^p$ convergence we have pointwise convergence entlang a subsequence.
	\begin{align*}
		u(t_2) - u(t_1) = \lim_{\epsilon \rightarrow 0} \int_{t_1}^{t_2}\int_{\mathbb{R}} \partial_t \eta_\epsilon (t - \tau) u(\tau) d\tau dt = \int_{t_1}^{t_2}\int_{\mathbb{R}} \lim_{\epsilon \rightarrow 0} \partial_t \eta_\epsilon (t - \tau) u(\tau) d\tau dt = \int_{t_1}^{t_2} \partial_t u(t)dt
	\end{align*}
\end{proof}

\begin{theorem}[The trace theorem in Bochner space]
	For $u \in W^{1, p}(0, T, X)$ we define
	\begin{align*}
		\texttt{spur}_{t_0}: = - \int_{t_0}^T \partial_t u \phi ds - \int_{t_0}^T u \phi_t ds \qquad \forall \phi \in C^\infty_c([0, T), X), \phi(t_0) = 1
	\end{align*}
	the trace mapping. Then we have
	\begin{align*}
		u(t) = \texttt{spur}_t(u) \quad a.e.
	\end{align*}
	Moreover, the mapping 
	\begin{align*}
		t \mapsto \textbf{spur}_t
	\end{align*}
	is also continuous.
\end{theorem}
\begin{remark}
The theorem above give a definition of $u(t_0)$ by the integral formula. In particular we may now define $u(T)$ by measurability(integrability). 	
\end{remark}
\begin{proof}
	\quad\\
	\textbf{1). } $\textbf{spur}_{t_0}$ is well defined by testing $\phi(t_0) = 0$.\\
	\textbf{2). } $\textbf{spur}_{t_0}$ is bounded linear.\\
	\textbf{3). } By mollification $u_\epsilon \rightarrow u$ in $L^p$, w.l.o.g. $u_\epsilon \rightarrow u$ a.e. (entlang a subsequence). Then the boundedness of the trace mapping 
	\begin{align*}
		\textbf{spur}_t(u_\epsilon) \rightarrow \textbf{spur}_t(u) \quad a.e.
	\end{align*}
\end{proof}

\begin{theorem}[Continuous Representation in Bochner space]
Given $ u \in L^1(0, T; X)$, $\partial_t u \in L^1(0, T; X)$, we have for $0 \le b \le a \le T$, up to a null set $N \subset [0, T]$ that:
\begin{align*}
	u(b) = u(a) + \int_a^b \partial_t u(s) ds , \quad u \in C^0([0, T], X)
\end{align*}
we have the continuous embedding
\begin{align*}
	W^{1, 1}(0, T, X) \hookrightarrow C^0([0, T], X)
\end{align*}
\end{theorem}
\begin{proof}
	
\end{proof}




\newpage
\section{The Representation Formulas}
\subsection{The Poisson Equation}
The object to be considered first is the Poisson equation:
\begin{equation}\label{PE}
\left\{
\begin{aligned}
	-\Delta u &= f  \quad x \in \Omega\\
	u(x) &= g(x) \quad x \in \Gamma_D\\
	A(x)\nabla u(x) &= h(x) \quad x \in \Gamma_N
\end{aligned}
\right.	
\end{equation}


We start with the simplest case where $f$ is a Dirac function(it does point landing):
\begin{equation*}
\delta_0: C^\infty_C(\Omega) \rightarrow \mathbb{R}, \varphi \mapsto \varphi(0)	
\end{equation*}
(By taking the integral of the fundamental solution $\Psi_d$ and an arbitrary smooth function $\varphi$, we would get a point value of the smooth function.)\\
Objective: Solve for the distribution,
\begin{align}
	 \int_{\mathbb{R}^d} u \Delta \varphi d\mu = \delta_0(\varphi) = \varphi(0) \quad \forall \varphi \in C_C^\infty(\Omega)
\end{align}
Therefore we only need $u \in L^1_{loc}(\Omega)$. In the Laplace equation where $0 \notin supp(\varphi)$. So instead we would like to solve:
\begin{align*}
	\int_{\mathbb{R}^n} u \Delta \varphi d\mu = 0 \quad \forall \varphi \in C^\infty_C(\Omega)
\end{align*} 

\subsubsection*{Find the distributional solution of $\Delta u = 0$}
The Laplace operator is invariant under rotations, in the sense that:
\begin{equation*}
	\forall \delta \textbf{ a ratation} : \Delta u(x) = \Delta  u(T_\delta x)
\end{equation*}
(\url{https://math.stackexchange.com/questions/970892/show-laplace-operator-is-rotationally-invariant/2617730#2617730})\\
Therefore we make our guess:
\begin{align*}
	u(x) = c \cdot |x|^\alpha
\end{align*}
Plug in $\Delta u = 0$ we get the distributional solution, which we call the fundamental solution of the Poisson equation.

\begin{remark}
Here we encounter the Dirac-function. The Dirac function is defined as(in the sense of distribution):
\begin{align*}
	\forall \varphi \in C_C^\infty(\Omega): \int_\Omega \delta \varphi d\mu = \varphi(0)
\end{align*}	
What we have calculated is a solution to $\Delta u = -\delta_0$ in the sense of distribution.
\end{remark}
By calculating $\Delta u = 0$ we get the fundamental \textbf{up to a constant}. Our objective is to get a weak solution for the Laplace equation, which satisfies:
\begin{align*}
		\forall \varphi \in C_C^\infty(\Omega): \int \Delta \Psi_d \varphi dx =   -\int \Psi_d \Delta \varphi dx = \varphi(0)	
\end{align*}
As the formula shows we have a \textbf{singularity} at $x = 0$, interchange integration and differentiation is not directly applicable. As the proof below shows, the usual technique to deal with the singularity is to cut a ball around it and approximate it by shrinking the ball. 

\begin{proposition}[The fundamental solution of Poisson equations]
 	For $d \in \mathbb{N}$, we define for $x \neq 0$: 
\begin{equation}\label{PE: fundamental solution}
	\Psi_d(x) = \left\{
	\begin{aligned}
		&-\frac{1}{2}|x| \quad &d = 1\\
		&-\frac{1}{2\pi} \log|x| \quad &d = 2\\
		& \frac{1}{d(d-2)\omega_d} |x|^{-(d-2)} \quad &d \geq3
	\end{aligned}
	\right.
\end{equation}
the fundamental solution. Then we have
\begin{align*}
	-\Delta \Psi_d = \delta_0 
\end{align*}
in the sense of distribution, i.e.,  
	\begin{align*}
		\forall \varphi \in C_C^\infty(\Omega): \int -\Delta \Psi_d \varphi dx = \varphi(0)
	\end{align*}
\end{proposition}
\begin{proof}
	W.l.o.g. assuming $\Omega = B_0(r)$. Since the integrand admits a singularity, we cannot apply bounded convergence theorem directly. Instead:
	\begin{align*}
		\int_{B_0(r)} \Psi_d \Delta \varphi dx = \lim_{\epsilon \rightarrow 0}\bigg(  \int_{B_0(r)\slash B_0(\epsilon)} \Psi_d \Delta \varphi dx + \int_{B_0(\epsilon)} \Psi_d \Delta \varphi dx\bigg)
	\end{align*}
	For the first part:
	\begin{align*}
		&\int_{B_0(r)\slash B_0(\epsilon) }  \Psi_d \Delta \varphi dx\\
		=& -  \underbrace{ \int_{B_0(r)\slash B_0(\epsilon)} \Delta \Psi_d \varphi dx}_{\rightarrow 0  \texttt{ since } \Delta \Psi_d \equiv 0}
		+ \underbrace{ \int_{\partial B_0(r)}  \Psi_d \nabla \varphi \cdot \nu dS}_{\rightarrow 0 \texttt{ since } \varphi \equiv 0}
		+ \underbrace{ \int_{\partial B_0(\epsilon)} \Psi_d \nabla \varphi \cdot \nu dS}_{:= I_\epsilon}
		- \underbrace{ \int_{\partial B_0(\epsilon)} \varphi \nabla \Psi_d \cdot \nu dS}_{:= J_\epsilon}
		- \underbrace{ \int_{\partial B_0(R)} \varphi \nabla \Psi_d \cdot \nu dS}_{\rightarrow 0}
		  \\
		 =& I_\epsilon - J_\epsilon
	\end{align*}

	Observe by the definition of fundamental solution:
	\begin{align*}
		I_\epsilon &\le  \int_{\partial B(0, \epsilon)} \hat{c} \epsilon^{2 - d} \|\varphi\|_{L^\infty}  \\
		&\leq \hat{c} \|\varphi\|_\infty \epsilon^{2-d} \int_{\partial B(0, \epsilon)} 1 dS\\
		&= \hat{c} \epsilon^{2 - d} (d \omega_{d} \epsilon^{d - 1})\\
		&= \hat{C} \epsilon \rightarrow 0
	\end{align*}
To calculate $J_\epsilon$ observe:
	\begin{align*}
		\nabla \Psi_d(x) \cdot \nu &= -\bigg( (d-2) \frac{1}{d(d-2)\omega_d} \bigg) |x|^{-(d-1)}  \frac{x}{|x|} \cdot \underbrace{ (-\frac{x}{\epsilon}) }_{:= \nu(x)}\\
		&= \frac{1}{d\omega_d |x|^d} \cdot x^2 \cdot \frac{1}{\epsilon}\\
		&= \frac{1}{d\omega_d \epsilon^{d - 1} }\\
		&= \frac{1}{vol(\partial B_0(\epsilon))}
	\end{align*}
	Therefore we have:
	\begin{align*}
		\varphi(0) = \varphi(0) \underbrace{ \int_{\partial B_0(\epsilon)} \nabla \Psi_d(x) \cdot \nu dS}_{ = 1}
	\end{align*}
	hence
	\begin{align*}
		J_\epsilon - \varphi(0) &= \int_{\partial B_0(\epsilon)} (\varphi(x) - \varphi(0)) \nabla \Psi_d \cdot \nu dx\\
		&\leq \int_{\partial B_0(\epsilon)} c |x|\Psi_d \cdot \nu dx\\
		&\leq c\epsilon \int_{\partial B_0(\epsilon)} \frac{1}{m(\partial B_0(\epsilon))}d S\\
		&\rightarrow 0
	\end{align*}
	For the second part we observe $\psi_d$ are at least integrable. Restricted on shrinking domain we can prove that it converges to 0 by dominated convergence theorem. 
\end{proof} 

The importance of the fundamental solution is that, for $\Delta u = f$, under certain regularity condition, the convolution gives point-wise evaluation:
\begin{align*}
	u(x) &= \int_{\mathbb{R}^n} \Psi_d(x - y) f(y) dy\\
	\Delta u(x) &= \int_{\mathbb{R}^n} \Delta \Psi_d (x - y) f(y) dy = f(x)
\end{align*}
The \textbf{Question} now remains is under which regularity condition

\begin{theorem}
For the Poisson equation:
\begin{itemize}
	\item For $f \in C_0^1(\mathbb{R}^d)$, the convolution $u = \Psi_d * f$ lies in $C^2(\mathbb{R}^d)$ and satisfies $-\Delta u = f$, i.e. we have found a classical solution
	\item For $f \in C_C^0(\mathbb{R}^d)$, the convolution $u = \Psi_d * f$ lies in $C^1(\mathbb{R}^d)$ and satisfies $-\Delta u= f$ weakly, i.e.:
	\begin{align*}
		\forall \varphi \in C^\infty_C(\mathbb{R}^d): \int \nabla u \nabla \varphi dx = -\int f \varphi dx
	\end{align*}
\end{itemize}
\end{theorem}
\begin{remark}
$\Delta \Psi_d$ is not in $L^2$ anymore.	
\end{remark}

\begin{proof}We proof the first and leave the second for exercise.
	\begin{itemize}
		\item Step 1: $f \in C^\infty_C(\mathbb{R}^d)$\par 
		By applying the differentiation rule:
		\begin{align*}
			-\Delta u(x) &\overset{(1)}{=} -\Delta (\Psi_d * f)(x)\\
			&= -\Delta_x \int \psi_d(z) f(x-z) dz\\
			&= -\int \psi_d(z) \Delta_x f(x - z) dz \\
			(\varphi(z):= f(x-z))
			&= -\int \Psi_d(z) \Delta \varphi(z) dz\\
			(1)&= \delta_0(\varphi)\\
			&= \varphi(0)\\
			&= f(x)
		\end{align*}
		\item Step 2: the first and the second derivative of $u$ could be approximated by $\|f\|_{C_1}$. In an other word
		\begin{align*}
			\forall R > 0 \exists C_R: 
			\|u\|_{C_2(B_{R}(0))} = \|\Psi_d*f\|_{C_2(B_{R}(0))}
			            \leq C_R \|f\|_{C_1}
		\end{align*}
		Notice:
		\begin{align*}
			|u(x)| &= \int_{B_R} \Psi_d(x-z) f(z) dz\\
			(1)&\leq \|f\|_{L_\infty} \int_{B_R} \Psi_d(x-z)dz\\
			(2)&\leq \|f\|_{L_\infty} \|\Psi\|_{L^1(B_{2R})} \\
		\end{align*}
		Analog:
		\begin{align*}
			|\partial_j u(x)| &= \int_{B_R}  \Psi_d(x-z) \partial_j f(z)dz \texttt{ (property of convolution)}\\
			&\leq \|\Psi_d\|_{L^1(B_{2R})} \|\nabla f\|_{L^\infty(\Omega)}  \\ 
		\end{align*}
		Finally:
		\begin{align*}
			\partial_i \partial_j u(x) &= \int_{B_R} \partial_i\Psi_d(x-z) \partial_j f(z) dz\\
			&\leq \|\nabla \Psi_d\|_{L^1(B_{2R})} \|\nabla f\|_{L^\infty(\Omega)}
		\end{align*}
		(For details of the proof in convolution see the section of \textbf{MOLLIFICATION})\par
		So we see that the $C^2$-norm of $u$ could be controlled by the $C^1$-norm of $f$.
		\item Step 3: Approximation and limit passage:\par 
		By Stone-Weierstrass we know every continuous function on a compact support could be uniformly approximated by the smooth function. Assume $\|f_n - f \|_{C^1(B_R)}$. We define
		\begin{align*}
			u_n:= \Psi_d * f_n
		\end{align*}
		By the result of STEP 2, a Cauchy sequence $f_n$ in $C^1$ gives a Cauchy sequence $u_n$ in $C^2$.
	\end{itemize}
\end{proof}


\subsubsection{Green's function}
The fundamental solution provides a method to solve the 
\begin{definition}[Green's identity]
	For $u, v \in C^2(\overline{\Omega})$, we have:
	\begin{equation}
		\begin{aligned}
			& \int_\Omega u \Delta v dx = \int_{\partial \Omega} u  \nabla v \cdot   \mathbf{n}  ds - \int_\Omega \nabla u \cdot  \nabla v dx \\
			& \int_\Omega u \Delta v dx = \int_{\partial \Omega}
			u \nabla v \cdot \mathbf{n} - v \nabla u \cdot \mathbf{n} ds 
			+ \int_\Omega v \Delta u dx\\
		\end{aligned}
	\end{equation}
\end{definition}

\begin{example}
If we replace the $\delta_0(\varphi)$ with $\delta_x(u)$, then we get:
\begin{align*}
	u(x) = \delta_x(u) &= \int_\Omega (- \Delta \Psi_d(x-z)) u(z) dz \\
	&= \int_\Omega \Delta u(z)(\Psi_d(x-z))dz + \int_{\partial \Omega} \psi_d(x-z)\nabla u(z)\cdot \nu dz - \int_{\partial \Omega}\nabla\Psi_d(x-z) u(z)dz
\end{align*}
\end{example}

In the following we would like to solve the Mixed Poisson Problem(MPP):
\begin{equation}
\textbf{(MPP)}\left\{
\begin{aligned}
	 	\Delta u &= f\\
	 	u &= g \quad \texttt{on} \quad \Gamma_{dir} \\
	 	\nabla u \cdot \nu &= h \quad \texttt{on} \quad \Gamma_{neu}
\end{aligned}
\right.
\end{equation}

\begin{definition}[Green's function]
	A function $G:\bar{\Omega} \times \bar{\Omega} \rightarrow \mathbb{R}$ is called Green's function for \textbf{(MPP)} if the function $ G(x, z) = H(x, z) + \Psi_d(x, z)$ satisfies:
	\begin{enumerate}
		\item $H \in C^2(\Omega \times \Omega) \cap C^1(\bar{\Omega} \times \bar{\Omega})$
		\item $\forall x \in \Omega : \Delta_z H(x, \cdot) = 0$
		\item $\forall y \in \Gamma_{dir}, x \in \Omega: G(x, y) = 0$
		\item $\forall y \in \Gamma_{neu}, x \in \Omega: \nabla_y G \cdot \nu = 0$
	\end{enumerate}
\end{definition}

\begin{theorem}[Representation via Green's function]
	If the Green's function exists, then the solution of MPP also exists and has the form:
	\begin{equation}\label{PE: Green's formula}
		u(x) = \int_\Omega G(x, y)f(y)dy - \int_{\Gamma_{dir}}\nabla_x G(x, y)\cdot \nu(y) \cdot  g(y) dy + \int_{\Gamma_{Neu}} G(x,y) h(y)dy
	\end{equation}
\end{theorem}
\begin{proof}
	By the calculation above:
	\begin{align*}
		u(x) &= \int_\Omega f(y) \Psi_d(x-z)dz 
		+ \int_{\partial \Omega} \psi_d(x-z)h(y) dz 
		- \int_{\partial \Omega}\nabla\Psi_d(x-z) \cdot \nu g(y) dz
	\end{align*}
	Apply (2) and the second Green's identity:
	\begin{align*}
		0 &= \int_{\Omega} -\Delta_y H(x, y) u(y) dy \\
		&= \int_{\Omega} H(x, y) f(y) dy 
		  + \int_{\partial \Omega} H(x, y) h(y) dy
		  - \int_{\partial \Omega} \nabla_y H(x, y) \cdot \nu(y) g(y) dy
	\end{align*}
	Add to $u(x)$ we get:
	\begin{align*}
		u(x) = \int_{\Omega} G(x, y) f(y)dy
				+\int_{\Gamma_{Neu}} G(x,y)h(y)dy
				 - \int_{\Gamma_{Dir}}\nabla_y G(x, y)\cdot \nu(y) g(y)dy
	\end{align*}
	Apply (3) and (4) we get the result.
\end{proof}

\begin{example}
	The function 
	\begin{align*}
		z \rightarrow \Psi_d(x - Rz)
	\end{align*}
	satisfies
	\begin{align*}
		-\Delta \Psi_d (x -Rz) = \delta_x = -\Delta \Psi_d(x -z)
	\end{align*}
	In the reflected space $R\Omega$($R$ is the reflection operation, the reflected space is defined as the symmetry of the original space with respect to the boundary axis. $R((x, y)) = (x, -y)$ ) we have:
	\begin{align*}
		-\Delta \Psi_d(x - z) = 0
	\end{align*}
	Now define
	\begin{align*}
		H(x, z) = -\Psi_d(x - Rz)
	\end{align*}
	we see that $\Delta_z \Psi_d(x - Rz) = 0$. and $G = H +\Psi$ satisfies the boundary condition of the Green's function. 
\end{example}

\subsubsection{The Laplace equation in the halfspace}



\subsection{The Representation via Fourier's series}
In this section we focus on the Dirichlet Problem(DP):
\begin{equation}\label{PE:DP(simple)}
(DP)\left\{
\begin{aligned}
	\Delta u &= f \\
	u_0 &= 0 \quad u \in \Gamma_{dir} \subset \partial\Omega
\end{aligned}
\right.	
\end{equation} 
We know for a \textbf{cuboid} $\Omega \subset \mathbb{R}^d$, we have an ONB of $L^2(\Omega)$ (denoted by $(0, l_1)\times \dots (0, l_d)) $:
\begin{equation}
	\varphi_k = c_\pi \prod_{i = 1}^d \sin(\frac{2\pi k_i}{l_i}x_i)	, \quad c_\pi = \sqrt{\frac{2}{vol(\Omega)}}
\end{equation} 
The basis function is nicely connected to the Laplace operator in the sense that:
\begin{equation*}
	-\Delta \varphi_k = -\sum_{i = 1}^d (\frac{2\pi k_i}{l_i})^2-\varphi_k = \lambda_k \varphi_k, \quad \lambda_k = \sum (\frac{2\pi k_i}{l_i})^2
\end{equation*}

\begin{theorem}[Strong solutions on cuboids]
	Let $\Omega \subset \mathbb{R}^d$ be a cuboid, then
	\begin{align*}
		\forall f\in L^2(\Omega) \exists!u \in H^{2}: \textbf{u satisfies (DP)}
	\end{align*}
\end{theorem}
\begin{proof}.
\begin{enumerate}
	\item \textbf{Approximation in finite dimension}
	\begin{align*}
		f_N = \sum^N_{k = 1} g_k \varphi_k, \quad g_k = \langle f, \varphi_k \rangle \quad (g \in l^2(\mathbb{Z}^d, \mathbb{R}))
	\end{align*}
	we claim that $u_N = \sum_{|k| \leq N} \frac{1}{\lambda_k} g_k \varphi_k $ solves the finite dimensional Poisson problem, indeed
	\begin{align*}
		\Delta u_N \overset{(1)}{=} \sum_{|k| \leq N} -\lambda_k \frac{1}{\lambda_k} g_k \varphi_k = \sum_{|k| \leq N} -g_k \varphi_k = -f_N
	\end{align*}
	where (1) follows from the property the cuboid $L^2$ space in $\mathbb{R}^d$.
	\item \textbf{$u_N$ converges in $H^2(\Omega)$}\\
	1)For $\|u_M - u_N\|_2^2 \rightarrow 0$ :
	\begin{align*}
		\|u_M - u_N\|^2_2 &= \| \sum_M^N \frac{g_k}{\lambda_k}\varphi_k\|^2_2\\
		&= \sum_M^N \frac{g_k}{\lambda_k}\\
		&\leq \frac{1}{\lambda_{\min}} \sum_M^N g_k
	\end{align*}
	
	By the Parval identity and the fact that $f = \sum g_k \varphi_k$ we know the sum goes to 0 as $N\rightarrow \infty $ \\
	
	2)For $\|\partial u_N - \partial u_N\|^2_s \rightarrow 0$:
	\begin{align*}
		\partial_i \varphi_k = c_\pi \prod_{i \neq j} \sin(\frac{2\pi k_i}{l_i} x_i) \frac{2\pi k_i}{l_i} \cos (\frac{2\pi k_j }{l_j} x_j)
	\end{align*}
	We can prove that $\partial \varphi_k$ also forms a OGS. By $\partial u_M = \sum^M_k \frac{g_k}{\lambda_k} \partial \varphi_k $ we get
	\begin{align*}
		\|\partial_i u_M - \partial_i u_N\|^2_2 &= \sum_M^N (\frac{g_k}{\lambda_k})^2 (\frac{2\pi k_i}{l_i})^2\\
		&\leq \sum^N_M g_k^2 \cdot C
	\end{align*}
	where $C$ is a constant. By $f\in L^2(\Omega)$ we get the norm goes to 0 as $M, N$ goes to infinity.\\
	3)For$\|\partial_i \partial_j u_N - \partial_i \partial_j u_M\|^2_2 \rightarrow 0$:\\
	\item \textbf{Show that $u$ indeed solves the Poisson equation}\\
	cONB
\end{enumerate}
\end{proof}
\begin{remark}
By solving the Poisson problem on a rectangle we may use Fourier series.
\end{remark}


\subsection{The Heat Equation on the  half-space}
First we investigate the homogeneous heat equation:
\begin{equation} 
\textbf{(HE)}\left\{
\begin{aligned}
	\partial_t u - \Delta_x u = 0 \quad &(t, x) \in (0, \infty) \times \mathbb{R}^d\\
	u(0, x) = u_0(x) \quad &x \in \mathbb{R}^d
\end{aligned}
\right.	
\end{equation}


\subsubsection*{Derivation of Heat Kernel}
We apply transform to the \textbf{(HE)}, such that the Differentiation could be transformed into multiplication of monoids:
\begin{align*}
	F[\partial_t u] &= \hat{\partial_t u} = \partial_t \hat{u} \\
	F[\Delta u] &= \sum_i^d F[\partial_{ii} u]\\
	&= \sum_i^d -\xi_j^2 F[u]\\
\end{align*}
Then we can transform \textbf{(HE)} into:
\begin{align*}
	\hat{u}_t(t, \xi) &= -|\xi|^2 \hat{u}(t, \xi) \\
	\hat{u}(0, \xi) &= \hat{u_0}(\xi)
\end{align*}
Solve the homogeneous ODE we get:
\begin{align*}
	\hat{u}(t, \xi) = \hat{u_0}(\xi) e^{-|\xi|^2t}
\end{align*}
Apply the inverse Fourier transformation:
\begin{align*}
	F^{-1}[\hat{u}] &= F^{-1}[\hat{u_0}(\xi) e^{-|\xi|^2t}]\\
	&= C F^{-1}[\hat{u_0}] * F^{-1}[e^{-|\xi|^2 t}]\\
	(1)&= \int u_0(y)w(t, x-y)dy
\end{align*}

\begin{definition}[Heat Kernel]	
The Heat kernel is defined as:
\begin{align*}
	H_d(t, x) := \frac{1}{(4\pi t)^{d/2}} e^{-\frac{|x|^2}{4t}}
\end{align*}
\end{definition}

\begin{lemma}[Properties of the heat kernel]
	The heat kernel satisfies:
	\begin{enumerate}
		\item $\forall t >0: \int_{\mathbb{R}^d} H_d(t,x) dx = 1 $
		, $H_d(t, x) > 0$
		\item $H_d(t, x) \in C^w((0, \infty), \mathbb{R}^d)$
		\item $\partial_t H_d = \Delta H_d$ in $(0, \infty) \times \mathbb{R}^d$.
		\item Duality in the $BC(\mathbb{R}^d)$
		\begin{align*}
			\forall \varphi \in BC(\mathbb{R}^d): \int \lim_{t \rightarrow 0} H_d \varphi dx = \varphi(0)
		\end{align*}
	\end{enumerate}
\end{lemma}
\begin{proof}
	For (4):
	\begin{align*}
		\int \lim_t H_d(t, x) \varphi(x) dx - \varphi(0) & \overset{(1)}{=}
		\int \lim_t H_d(t, x) (\varphi(x) - \varphi(0)) dx\\
		&= \int_{B_R(0)} \lim_t H_d(t, x) (\varphi(x) - \varphi(0) dx 
		+ \int_{\mathbb{R}^d/B_R(0)} \lim_t H_d(t, x) (\varphi(x) - \varphi(0)) dx\\
	 (2)&\leq  \frac{\epsilon}{2} + \|\varphi\|_\infty \int_{\mathbb{R}^d/B_R(0)} H_d(t, x)dx    \\
	 (3)&\leq  \frac{\epsilon}{2} + \frac{\epsilon}{2}\\
	 &= \epsilon
	\end{align*}
\end{proof}

\begin{theorem}[Representation formula of classical solution]
	Given the data $u_0 \in BC(\mathbb{R}^d)$, then
\begin{align*}
	u(t, x) = \int_{\mathbb{R}^d} H_d(t, x - y)u_0(y)dy
\end{align*}
is a classical(even an analytical) solution to the homogeneous heat equation \ref{HE:homogeneous}.
\end{theorem}
\begin{proof}
\begin{align*}
	\partial_t (H_d*u_0) - \Delta_x (H_d*u_0) &\overset{(1)}{=} \int \partial_t H_d(t, x -y) u_0(y) - \Delta_x H_d(t, x -y)u_0(y) dy\\
	&= \int u_0(y) (\partial_t - \Delta_x) H_d(t, x -y)dy\\
	&= 0
\end{align*}
\end{proof}
\begin{remark}.
\begin{enumerate}
	\item The homogeneous heat equations have infinitely many solutions.
	\item models violate the theory of relativity.
\end{enumerate}	
\end{remark}

\subsubsection*{The inhomogeneous case}
In the inhomogeneous case, the equation becomes:
\begin{equation} \label{HE:homogeneous}
\textbf{(HE)}\left\{
\begin{aligned}
	\partial_t u - \Delta_x u = f(t, x) \quad &(t, x) \in (0, \infty) \times \mathbb{R}^d\\
	u(0, x) = u_0(x) \quad &x \in \mathbb{R}^d
\end{aligned}
\right.	
\end{equation}
the Fourier transform become:
\begin{align*}
	\hat{u}_t(t, \xi) &= -|\xi|^2 \hat{u}(t, \xi) + \hat{f}(t, \xi) \\
	\hat{u}(0, \xi) &= \hat{u_0}(\xi)
\end{align*}
Use the method of solving linear ODE we get:
\begin{align*}
	\hat{u_p} = \int_0^t e^{(t - s)|\xi|^2} \hat{f}(s, \xi)ds
\end{align*}
Therefore
\begin{align*}
	\hat{u} = \hat{u}_f + \hat{u}_p
\end{align*}
Apply the inverse Fourier transform:
\begin{align*}
	F^{-1}[\hat{u}](x) &= F^{-1}[\hat{u_f}] + F^{-1}[\hat{u_p}]\\
	&= H_d*u_0 + F^{-1}[\int_0^t e^{(t-s)|\xi|^2} \hat{f}(s, \xi)ds]\\
 (1)&= \dots + \int_0^t F^{-1} [e^{(t-s)|\xi|^2} \hat{f}(s, \xi)]\\
 &= H_d(t, x)*u_0 + \int_0^t H_d(t-s, x)*f(s, x) ds
\end{align*}
For (1) the DUIS to work we need $f$ to be reasonable integrable.

\begin{definition}[mild solution]
The solution we get above:
\begin{equation}
	H_d(t, x)*u_0 + \int_0^t H_d(t-s, x)*f ds
\end{equation}
is called mild solution of the inhomogeneous heat equation.
\end{definition}
\begin{remark}
The construction of $u$ is related to Duhamel's principle. Notice that $H_d(t-s, x)*f$ solves the homogeneous heat equation with starting time $s$ instead of $0$.	
\end{remark}

\begin{theorem}
	Assume $u_0 \in BC(\mathbb{R}^d)$, the $f, \partial_t f, \nabla f, \nabla^2 f \in BC[[0, \infty) \times \mathbb{R}^d)$. Then the mild solution is a classical solution.
\end{theorem}
Apply parametric integral
\begin{align*}
	\partial_t u - \Delta u &= (\partial_t - \Delta_x) \int_0^t \int_{\mathbb{R}^d} H_d(t - s, x - y) f(s, y) dy ds\\
	&= (\partial_t - \Delta_x) \int_0^t \int_{\mathbb{R}^d}
	   H_d(s, y) f(t - s, x -y)dy ds\\
	&= \int_0^t \int_{\mathbb{R}^d} (\partial_t - \Delta_x)H_d(s, y)f(t-s, x-y)dyds + \int_{\mathbb{R}_d}H_d(t, y)f(0, x-y)dy 
\end{align*}
\begin{remark}
\end{remark}


\subsubsection{Smoothing Properties}
From the last section we learned:
\begin{align*}
	u(t, x) = H_d * u_0,\quad \forall t >0: u(t, \cdot) \in C^\infty(\mathbb{R}^d)
\end{align*}
The solution is bounded, for $t \rightarrow \infty$, $u$ would become flat. The smoothing properties is the study of the decreasing rate.
\begin{align*}
	\nabla_x H(x, t) &= \frac{C}{t^{d/2}} e^{\frac{-|x|^2}{4t}} \nabla_x (\frac{-|x|^2}{4t})\\
	&= \frac{C}{t^{d/2}} e^{\frac{-|x|^2}{4t}}(-\frac{1}{4t}2x)\\
	&= \frac{C_1}{t^{(d/2)+ 1}} x e^{\frac{-|x|^2}{4t}}\\
	\nabla^2_x H(x, t)
	&= \frac{C_1}{t^{d/2}}(\frac{I_d}{t} + \frac{C_2}{t^2} x \otimes x) e^{\frac{-|x|^2}{4t}} = \frac{C_1}{t^{d/2 + 2}}(I_d t + C_2 x \otimes x) e^{\frac{-|x|^2}{4t}} \\
\end{align*}
By induction we get:


\begin{theorem}[Smoothing Properties]
	For each $\alpha \in \mathbb{N}_0^d$, there exists a $\tilde{C}_\alpha = C_\alpha M_\alpha > 0:$
	\begin{enumerate}
		\item
		\begin{align*}
			\|\partial_x^\alpha u(u, \cdot)\|_\infty \leq
			\frac{\tilde{C}_\alpha}{t^{|\alpha|/2}} \|u_0\|_\infty
		\end{align*}
	\end{enumerate}
\end{theorem}

\begin{align*}
	\partial_j(|x|) &= \partial_j (\sqrt{x_1^2 + \dots x_d^2})\\
	&=\frac{1}{2} (x_1^d + \dots + x_d^2)^{-\frac{1}{2}} \partial_j (\sum^d x_i^2)\\
	&= x_j |x|^{-1}
\end{align*}
\begin{align*}
	\partial_j(\sum x_i^2) = 2x_j
\end{align*}












\newpage
\section{The Elliptic Problems via Energy Method}
Observe the second type of Dirichlet Problem:
\begin{equation}\label{PE:DP(2)}
\textbf{(DP)} \left\{
\begin{aligned}
	-\div(A(x) \nabla u) = f \quad x &\in \Omega\\
	u = g \quad x &\in \Gamma_{dir}
\end{aligned}
\right.	
\end{equation}
In this section we study the solution of the \textbf{Weak Dirichlet Problem} (WDP):
\begin{equation}
\begin{aligned}
	\int_\Omega \nabla \varphi A(x) \nabla u dx = \int_\Omega \varphi f dx
	\quad \forall \varphi &\in H^1_0(\Omega)\\
	 u - g \in H^1_0(\Omega) \quad x &\in \Gamma_{dir}
\end{aligned}	
\end{equation}
In order for $g$ to be well-defined, we first need the trace theorem. 


\subsection{The Trace Theorem}
We know $\forall u \in C(\Omega)$, the trace mapping:
\begin{equation*}
\gamma\left\{
\begin{aligned}
	u &\mapsto u_{|\partial \Omega}\\
	C(\bar{\Omega}) &\rightarrow C(\partial \Omega)
\end{aligned}
\right.	
\end{equation*}
is well-defined. But for $u \in L^2(\Omega)$, since $u$ is only defined almost everywhere, it's impossible to define the evaluation norm or the trace mapping of the function $u$. In order to have a proper boundary value, we need to know the trace mapping is well-defined.

\begin{theorem}[Existence of traces]
	Given a d-dimensional domain $\Omega$ with Lipschitz boundary $\partial \Omega$. $\forall p \in [1, \infty] \forall u \in W^{1, p}:$ we have a unique mapping $\gamma$
	\begin{equation*}
	\left\{
	\begin{aligned}
		W^{1, p}(\bar{\Omega}) &\rightarrow L^p\\
		u &\mapsto \gamma(u) 
	\end{aligned}
	\right.	
	\end{equation*}
	such that $\gamma$ extends the classical trace mapping from $C^1_c(\bar{\Omega})$ to $C(\partial\Omega)$.
\end{theorem}
\begin{proof}
	Here we only give a special case where $p = 2$. We want to show
	\begin{align*}
		\forall u \in H^1(\bar{\Omega})\exists \gamma(u) \in H^1(\partial \Omega): \text{$\gamma(u)$ extends the trace mapping of the dense subset $C^1(\bar{\Omega})$}
	\end{align*}
	i.e. for a sequence $(f_k)$ in the dense subset, the image of the trace mapping would converge in low-dimensional $L^p$ norm
\begin{itemize}
\item \textbf{Step 1}: Give a cylinder domain $\Omega = (0, l) \times \Sigma$. For $u \in C^1(\bar{\Omega})$ we know form the exercise:
	\begin{align*}
		\|u(x_1, \cdot) - u(\bar{x_1}, \cdot)\|_{L^2(\Sigma)}^2 \leq 
		\|\nabla u\|_{L^2(\Omega)} |x_1 - \bar{x_1}|^{\frac{1}{2}}
	\end{align*}
    Therefore $u \in C^1(\Omega) \Rightarrow u \in C^{\frac{1}{2}}((0, l), L^2(\Sigma))$. By the density of $C^1(\Omega)$ in $H^1(\Omega)$ we obtain the same inequality in $H^1(\Omega)$.
    \begin{align*}
    	\|\gamma(u)\|_{L^2} = \|\lim_{x_1 \rightarrow 0} u(x_1, \dot)\|_{L^2(\Sigma)}
    \end{align*}
    Notice the $(u(x_1, \cdot))_{x_1 \rightarrow 0}$ is a Cauchy sequence in $L^2(\Sigma)$ therefore admits an unique limit in $L^2(\Sigma)$
\end{itemize}
\end{proof}


\subsection{Symmetric Elliptic Problem}
\subsubsection{Dirichlet Problem}
\begin{lemma}[Poincare lemma]
Given $\Omega \subset \mathbb{R}^n$ open and bounded, $p \in [1, \infty)$, and there exists some $\nu \in \mathbb{R}^n$, such that 
\begin{align*}
	\Omega \subset \{x \in \mathbb{R}^n| |x \cdot \nu| \le R\},
\end{align*}
then there exists some constant $C_0$, such that
\begin{align*}
	\forall u \in W_0^{1,p}: \|u\|_{L^p} \leq C_0(p, diam(\Omega)) \|\nabla u\|_{L^p},
\end{align*}
\end{lemma}
\begin{remark}
This result can be extended to $\Omega$ bounded Lipschitz. 	
\end{remark}


\begin{proposition}[equivalence of the $\|\|_{H_1}$ and $\|\cdot\|_{A}$]
	Given a bounded domain $\Omega$ in $R^d$ and function system $A : \Omega \rightarrow \mathbb{R}^{d \times d}_{sym}$ satisfying:
	\begin{enumerate}
		\item \textbf{uniformly bounded: }$A \in L^\infty(\Omega, \mathbb{R}^{d \times d}_{sym})$
		\item $A$ \textbf{uniformly elliptic}(positive definite), i.e., there exits some $C > 0$ such that
		\begin{align*}
			\forall_{a.e.} x \in \Omega \forall \xi \in \mathbb{R}^d \exists \alpha >0 : 
			\xi A(x)\xi \geq C|\xi|^2
		\end{align*}
	\end{enumerate}
	Define a new norm with respect to the scalar product:
	\begin{align*}
		\langle \cdot , \cdot \rangle_A:= \int_\Omega \nabla u A \nabla u d\mu.
	\end{align*}
	Then there exists a constant $C > 0$:
	\begin{align*}
		\frac{1}{C} \|u\|_{H^1} \leq \|u\|_A^2 \leq C \|u\|_{H^1} 
	\end{align*}
\end{proposition}
\begin{remark}
Remember the $\langle .\rangle_A$ is defined using $\nabla u$	
\end{remark}

\begin{proof}
	\begin{align*}
		\|u\|_A^2 &:= \int_\Omega \nabla u A(x) \nabla u dx\\ 
		&\leq \int_\Omega |\nabla u|^2 A(x)dx\\
		&\leq \|A\|_\infty \|\nabla u\|_2^2 \\
		&\leq \|A\|_\infty \|u\|_{H_1}
	\end{align*}
	On the other hand:
	\begin{align*}
		\|u\|_{H_1}^2 &= \|u\|_2^2 + \|\nabla u\|^2_2\\
		(1)&\leq (1 + R^2) \|\nabla u\|^2_2\\
		(2)&\leq (1 + R^2) \frac{1}{c} \|\nabla u\|_A^2
	\end{align*}
	where in (1) we use the Poincare lemma and (2) the uniform ellipticity.
\end{proof}

\begin{theorem}[Existence and Uniqueness of (WDP)]
	 Given a bounded domain $\Omega$ and an uniformly SPD function $A(x)$. Suppose the data in the \textbf{(WDP)} satisfies 
	 \begin{enumerate}
	 	\item the source $f \in L^2(\Omega)$,
	 	\item the boundary $g \in H^1(\Omega)$ (not necessarily compact supported). 
	 \end{enumerate}
	  Then the \textbf{(WDP)}  admits a unique solution $u \in H^1(\Omega)$.
\end{theorem}
\begin{proof}.
\begin{enumerate}
	\item Step 1: $g = 0$: $\forall \varphi \in H^1_0(\Omega)$, we have by the Cauchy-Schwartz:
		\begin{align*}
		 	\int_\Omega f\varphi dx \leq \|f\|_{L^2} \|\varphi\|_{L^2} 
		 \end{align*}
		 Furthermore by the Poincare inequality:
		 \begin{align*}
		 	\|f\|_{L^2} \|\varphi\|_{L^2} \leq C_P \|f\|_{L^2} \|\varphi\|_{H_1}
		\end{align*}
		Finally by equivalence of the norm we have
		\begin{align*}
			\|f\|_{L^2} \|\varphi\|_{L^2} \leq C_P \|f\|_{L^2} \|\varphi\|_{A}
		\end{align*}
		The linear form 
		\begin{align*}
			F(\varphi) = \int_\Omega f \varphi dx,
		\end{align*}
		is therefore bounded in $(H^1_0(\Omega), \langle,\rangle_A)$, by Riesz the linear form can be represented by the scalar product, i.e., there exits $u \in H_0^1(\Omega)$, such that
		\begin{align*} 
		\langle u, \varphi \rangle_A = F(\varphi) \qquad \forall \varphi \in H^1_0(\Omega).	
		\end{align*}
		The boundary condition $g = 0$ implies $u \in H^1_0(\Omega)$. So the boundary condition is satisfied.
	\item Step 2: $g \neq 0$: \par 
	Define $w := u - g$. Notice $w \in H^1_0(\Omega)$. 
	\begin{align*}
		F(\varphi) &= \langle g,  \varphi \rangle_A + \langle w, \varphi \rangle_A
	\end{align*}
	\begin{align*}
		\langle w, \varphi \rangle_A &= F(\varphi) - \langle g , \varphi \rangle_A
	\end{align*}
	By the same argument above we have:
	\begin{align*}
		F(\varphi) - \langle g,  \varphi \rangle_A 
		&\leq \|f\|_{L^2} \|\varphi\|_{L^2} - \|g\|_{H_1} \|\varphi\|_{H_1} \\
		&\leq C(g, f, C_P)  \|\varphi\|_A
	\end{align*}
	Therefore $F(\varphi) - \langle g,  \varphi \rangle_A$ defines a bounded linear form in $(H^1_0(\Omega), \langle,\rangle_A)$. By Riesz we get an unique $w \in H^1_0(\Omega)$ with:
	\begin{align*}
		\langle w, \cdot \rangle_A = F(\cdot) - \langle g,  \cdot \rangle_A
	\end{align*} and therefore a unique $u = w + g \in H^1(\Omega) $.
\end{enumerate}
\end{proof}
\begin{remark}
$u, g$ are not necessarily in $H^1_0(\Omega)$	
\end{remark}







\subsubsection{Mixed Elliptic Problem}
We now return to the mixed elliptic problem: \textbf{(MP)}
\begin{equation}
\begin{aligned} \qquad
	-\div(A(x) \nabla u) &= f\\
	u &= g \qquad u \in \Gamma_{dir}\\
	(A(x)\nabla u)\cdot \nu + \alpha u &= h \qquad  u \in \Gamma_{neu} = \partial \Omega \slash \Gamma_{dir} 
\end{aligned}
\end{equation}

By the 1.st Green's identity we obtain:
\begin{align*}
	\int_\Omega f\varphi = \int_\Omega \nabla \varphi( A(x)\nabla u )dx - \int_{\Gamma_{neu}} (h - \alpha u)\varphi ds - 
	\int_{\Gamma_{dir}} (A(x)\nabla u \cdot \nu)\varphi ds 
\end{align*}
To obtain an scalar product on the righthand side, we need the last term to be 0, so $\varphi \in H^1_{\Gamma_{dir}}$. Now we rewrite the \textbf{(MP)} in the weak formulation: 
\begin{equation}
\begin{aligned}
	\int_{\Omega}f\varphi dx + \int_{\Gamma_{neu}} h\varphi dx &= \int_\Omega(A(x) \nabla u)\nabla\varphi dx + \int_{\Gamma_{neu}} \alpha u \varphi dx
	\quad &\forall \varphi \in H^1_{\Gamma_{dir}}(\Omega)\\
	\gamma(u - g) &= 0 \qquad  \texttt{on} \quad \Gamma_{dir}
\end{aligned}
\end{equation}
Here we introduced the new space:
\begin{align*}
	H^1_{\Gamma_D}:= \{u \in H^1(\Omega) |\quad \gamma(u)|_{\Gamma_D} = 0\}
\end{align*}
Again we would like to solve the \textbf{(WMP) }using Riesz. In order to apply Riesz, we need to defne a proper scalar product for the \textbf{(WMP)}.
\begin{align*}
	( u| \varphi )_{A, \alpha}:= \langle u, \varphi \rangle_{A, \Omega} + \alpha \langle u, \varphi \rangle_{\Gamma_{neu}}
\end{align*}
corresponding to the righthand side of the \textbf{(WMP).}

\begin{definition}[$H^1$ coecivity]
Given $\Omega \subset \mathbb{R}^d$ and $\partial \Omega = \Gamma_{dir} \overset{\circ}{\cup} \Gamma_{neu}$. $A$ is a matrix function satisfying all properties in proposition 5.2 other than uniform ellipticity.\par 
We say $\alpha \in L^\infty(\Gamma_N)$ satisfies the \textbf{$H^1$ coercivity condition} if:
\begin{align*}
	\exists c_0 > 0 \forall u \in H^1_{\Gamma_D}: \|u\|_{A, \alpha} \geq c_0 \|u\|_{H^1(\Omega)}
\end{align*}
\end{definition}

\begin{theorem}[Existence and Uniqueness of (WMP)]
	Given a Lipschitz domain $\Omega$. The data of the \textbf{(WMP)} satisfies the condition: $(f, g, h) \in L^2(\Omega) \times H^1(\Omega) \times L^2(\Gamma_N)$. Then the \textbf{(WMP)} admits an unique solution $u \in H^1(\Omega)$.
\end{theorem}
\begin{proof}.
\begin{itemize}
	\item \textbf{Step 1:} $(H^1_{\Gamma_D}, ( \cdot| \cdot )_{A, \alpha})$ is a Hilbert space:
	\begin{align*}
		\langle u, u \rangle_{A, \alpha} &\leq \|A\|_\infty \|\nabla u\|^2_2 + \|\alpha\|_\infty \|u\|_2^2\\
		&\leq \max\{\|A\|_\infty, \|\alpha\|_\infty\} \|u\|^2_{H^1}
	\end{align*} 
	And by coercivity we have:
	\begin{align*}
		\|u\|_{A, \alpha}^2 \geq c_0 \|u\|_{H_1}
	\end{align*}
	The norms are therefore equivalent. By the trace extension of $H^1$ we know $(H^1_{\Gamma_D}, \langle, \rangle_{A, \alpha})$ is a closed subspace of $H^1(\Omega)$ and therefore defines a Hilbert space.
	\item \textbf{Step 2:} $g = 0$. \par 
	We show $F(\varphi) = \int_\Omega f\varphi dx + \int_{\Gamma_{new}} h \varphi dx$ indeed defines bounded linear form in $(H^1_D, \langle, \rangle_{A, \alpha})$.
	\begin{align*}
		F(\varphi) &\leq \|f\|_{2, \Omega} \|\varphi\|_{2, \Omega} + \|h\|_{2, \Gamma_N} \|\varphi\|_{2, \Gamma_N}
	\end{align*}
	therefore $F$ is bounded.
	\item \textbf{Step 3: $g \neq 0$}
\end{itemize}	
\end{proof}


\subsubsection{Pure Neumann Problem}
Now we observe the pure Neumann problem:
\begin{equation}
\textbf{(NP)}\left\{
\begin{aligned}
	-\div(A(x)\nabla u) &= f\\
	A(x)\nabla u \cdot \nu &= h \quad u \in \partial\Omega 
\end{aligned}	
\right.
\end{equation}
There is a non-trivial solution, namely $u \equiv const$ to the homogenous case $(f, g) \equiv (0, 0)$. By superposition principle we know in the inhomogeneous case a solution plus a constant would also solve the \textbf{(NP)}. We would like "quotient out" the space such that the solution is unique. Observe:
\begin{align*}
	&\int_\Omega u_1 dx = 0, u_2 = u_1 + \alpha\\
	\Rightarrow &\int_\Omega u_2 dx = \int_\Omega u_1 + \alpha dx\\
	\Rightarrow &\alpha = 0
\end{align*}
Intuitively, we would get an unique solution in the quotient space $H^1(\Omega)/ \{\int_\Omega u dx = 0\}$. \\
Integrating by parts yield:
\begin{equation}
\textbf{(WNP)}\left\{
\begin{aligned}
	&\int_\Omega (A(x) \nabla u) \nabla\varphi dx = \int_\Omega f\varphi dx + \int_{\Gamma_{neu}}h \varphi ds \quad \forall \varphi \in H^1(\Omega) \\
	& u \in H_{av}^1(\Omega)
\end{aligned}	
\right.
\end{equation}

\begin{proposition}[Poincare-Wirtinger inequality]
	Given a bounded Lipschitz domain, we have: 
	\begin{align*}
		\exists C_{PW} \forall u \in H^1(\Omega)\text{ with } \int_\Omega u dx = 0: \|u\|_{L^2}^2 \leq C_{PW} \|\nabla u\|_{L^2}^2
	\end{align*}
\end{proposition}
\begin{proof}
	
\end{proof}


\begin{theorem}[Existence and Uniqueness of (WNP)]
	Given $\Omega$ a bounded, simply connected domain with Lipschitz-boundary. The \textbf{(WNP)} admits a solution $u \in H^1(\Omega)$,  that is unique up to a constant if and only if:
	\begin{align*}
		\int_\Omega fdx + \int_{\partial \Omega} hdx = 0
	\end{align*}
\end{theorem}
\begin{proof}.
\begin{itemize}
\item \textbf{Step 1: Introduce $H^1_{av}(\Omega)$} 
Introduce the adapted space:
\begin{equation}
	H^1_{av}(\Omega) := \{ u \in H^1(\Omega)| \int_\Omega u dx = 0\}
\end{equation}
$H^1_{av}(\Omega)$ is a closed subspace since it's the pre-image of the linear form:
\begin{align*}
	F(u):= ( \mathbbm{1}_{\Omega}| u )_{H_1}
\end{align*}
since $\Omega$ is bounded $\mathbbm{1}_{\Omega} \in H^1(\Omega)$ :
\begin{align*}
	&F(u) = \int_\Omega u \mathbbm{1}_{\Omega} dx + \int_\Omega \nabla u  \nabla \mathbbm{1}_{\Omega} dx \\
	\Rightarrow &F^{-1}(0) = \{u |\int_\Omega u dx = 0 \}
\end{align*}
We've established $H_{av}^1(\Omega)$ is a closed subspace of $H^1(\Omega)$ and in the proof of (WDP) we know $\langle \rangle_{H_1}$ and $\langle \rangle_{A}$ are equivalent norms in $H^1$.\par
Next we prove the righthand side defines a bounded linear form, indeed
\begin{align*}
	F(\varphi)&:= \int_\Omega f\varphi dx + \int_{\Gamma_N} h\varphi dx\\
	&\leq (\|f\|_2 + \|h\|_2) \|\varphi\|_2 \\
	&\leq C(\|\varphi\|_2 + \frac{1}{C_{PW}}\|\varphi\|_2)\\
 (1)&\leq \tilde{C} \|\varphi\|_{H_1}
\end{align*}
where (1) we apply the Poincare-Wirtinger inequality.\\
Now apply Riesz we know the problem:
\begin{equation*}
(\tilde{\textbf{WNP}})\left\{
\begin{aligned}
	&\int_\Omega \nabla u A(x) \nabla \varphi = \int_\Omega f\varphi dx + \int_{\Gamma_N} h\varphi dx \quad \forall \varphi \in H^1_{av}(\Omega)\\
	&u \in H^1_{av}(\Omega)
\end{aligned}
\right.
\end{equation*}
admits an unique solution.
\item \textbf{Step 2: Existence}\\
We know show:
\begin{align*}
	\int_\Omega f dx + \int_{\partial \Omega} h dx = 0 \Leftrightarrow
	\text{There exists a solution for \textbf{(WNP)}}
\end{align*}
"$\Rightarrow$":\\
$\forall \varphi \in H_{av}^1(\Omega) \exists! u \in H_{av}^1(\Omega) : F(\varphi) = \langle u , \varphi \rangle_A$ \\
We still need to check with the test function $\varphi \equiv 1$
\begin{align*}
	F(\varphi) &= 0
\end{align*}
therefore $u$ solves \textbf{(WNP)}\\
"$\Leftarrow$": \\
Suppose $\int_\Omega f dx + \int_{\partial \Omega} h dx \neq 0$: Inserting $\varphi \equiv 1$, knowing that there can not be a solution.
\item 
\textbf{Step 3: Uniqueness}\\  
Assuming there exists $u_1, u_2 \in H^1(\Omega)$ that don't differ by a constant, solves \textbf{(WNP)}. Then there exists 2 different $\tilde{u_1}, \tilde{u_2} \in H^1_{av}(\Omega)$ solving $\tilde{(\textbf{WNP})}$. Contradiction.
\end{itemize}
\end{proof}




\subsection{Asymmetric Elliptic Problem and Lax-Milgram Theorem}
\begin{theorem}[Lax-Milgram theorem]
	Given a bounded coercive bilinear form $B(u, v)$ in a Hilbert space $H$, then $\forall F(\varphi) \in H^* \exists!u \in H:F(\varphi) = B(u, \varphi)$. To be concrete:\par 
	Given a bilinear form in $B:H \times H \rightarrow \mathbb{R}$ satisfying:
	\begin{enumerate}
		\item \textbf{boundedness:}$\exists C_1>0\forall u, v\in H: B(u, v) \leq C_1 \|u\|\|v\|$
		\item \textbf{coercive:}(ellipticity) $\exists c_1 > 0 \forall u \in H: B(u, u) \geq c_0 \|u\|^2$
	\end{enumerate}
	Then we have:
	\begin{align*}
		\forall F \in H^* \exists u \in H: F(\varphi) = B(u, \varphi)\quad \forall \varphi \in H
	\end{align*}
\end{theorem}
\begin{proof} We \textbf{refrain ourself to the separable Hilbert space}, where we have a cONB: $\{ \psi
_m | m \in \mathbb{N} \}$
\begin{itemize}
\item \textbf{Step1(finite representation)}: 
For the subspace spanned by $V_N = \spa \{ \psi_m | m \in \{1, \dots, N\} \}$, find $u_N$, such that 
\begin{align*}
	\textbf{(PN)} \quad \forall F \in H^* \exists u_N \in H: F(\varphi) = B(u_N, \varphi) \quad \forall \varphi \in V_N
\end{align*}
For $u_N = \sum^N_i \alpha_i \psi_i$, denote $f_i = F(\psi_i) $. Write $A_N$ as $A_{i, k} = B(\psi_i, \psi_k ) \in \mathbb{R} $, \textbf{(PN)} could be written as a linear system:
\begin{align*}
	A^N 
	\begin{bmatrix}
		\alpha^1 \\
		\dots\\
		\alpha^N
	\end{bmatrix}
	= 
	\begin{bmatrix}
		f_1 \\
		\dots\\
		f_N
	\end{bmatrix}
\end{align*}
Since $\ker(A_N) = \{0\}$ by the \textbf{coercivity}, $A_N$ is invertable. We have an unique vectot $\alpha_N$.
\item \textbf{Step 2: (a priori estimate)}
We want to apply Banach-Alaoglu. Therefore we have to show the boundedness of the sequence $u_N$. Indeed by coercivity of the Bilinear form $B(u, v)$ we have:
\begin{align*}
	c_0 \|u_N\|_H^2 &\leq B(u_N, u_N) = F(u_N) \leq \|F\|_{H^*}\|u_N\|_H\\
	\Rightarrow 
	 \|u_N\|_H &\leq \frac{1}{c_0}  \|F\|_{H^*}
\end{align*}
\item \textbf{Step 3: (convergent subsequence)}
By Banach-Alaoglu we know every bounded sequence in the separable Hilbert space admits a weakly convergent subsequence. By passing through a subsequence we know
\begin{align*}
	\exists u^* \in H: \quad u_N \overset{weak}{\rightharpoonup} u^*
\end{align*}
\item \textbf{Step 4:($u^*$ solves the P in infinite dimension):} We know
\begin{align*}
	B(u_N, \varphi_M) = F(\varphi_M) \quad \forall \varphi_M \in V_M \subset V_N
\end{align*}
to show:
\begin{align*}
	B(u^*, \varphi) = F(\varphi) \qquad \forall \varphi \in H
\end{align*}
observe for all $\varphi_M$ arbitrary, since $B(\varphi_M, \cdot)$ defines a bounded functional on H, we have by weak convergence
\begin{align*}
	F(\varphi_M) = B(u_N, \varphi_M) \longrightarrow B(u^*, \varphi_M),
\end{align*}
and by the density argument:
\begin{align*}
	B(u^*, \varphi) = \underbrace{ B(u^*, \varphi_M)}_{\rightarrow B(u^*, \varphi_M) } + \underbrace{ B(u^*, \varphi - \varphi_M)}_{\rightarrow 0}
\end{align*}
\end{itemize}
\end{proof}
\begin{remark}
Since only the reflexivity is essential here for the limit passage, we may weaken the condition such that for the bilinear form $B: X \times Y \rightarrow \mathbb{R}$, $X$ is Hilbert and $Y$ is Banach.\\
There is an alternative proof using Riesz and Banach fixed point. 	
\end{remark}



With Lax-Milgram we are able to solve the asymmetrical problems. Now we don't need $B(u, v) = B(v, u)$. Observe the elliptic problem
\begin{equation}\label{EP}
(EP)\left\{
\begin{aligned}
	div(A(x)\nabla u) + b\cdot \nabla u + cu &= f \\
	u &= g \qquad u \in \Gamma_{dir}\\
	A\nabla u \cdot \nu + \alpha u &= h \qquad u \in \Gamma_{neu}
\end{aligned}
\right.
\end{equation}
We need $A \in L^\infty(\Omega, \mathbb{R}^d \times \mathbb{R}^d)$ but not symmetric. Other coefficients $b, c, \alpha$ also need to be in $L^\infty$, and $A$ needs to satisfy the ellipticity condition:
\begin{align*}
	\exists c_0 > 0 \forall \xi \in \mathbb{R}^d \forall_{a.e} x \in \Omega: \xi A(x) \xi \geq c_0 |\xi|^2
\end{align*} 
Integrating by parts in the first term and we get the weak form of the (EP):
\begin{equation}\label{EP:WEP}
\textbf{(WEP)}\left\{
\begin{aligned}
	&\int_\Omega f\varphi dx + \int_{\Gamma_{neu}} h \varphi dx = \int_\Omega \nabla\varphi A(x)\nabla u dx + \int_\Omega b\nabla u \varphi dx + \int_\Omega cu \varphi dx \qquad \forall \varphi \in H^1_{\Gamma_D}(\Omega) \\
	&u - g \in H^1_{\Gamma_{D}}(\Omega)
\end{aligned}
\right.	
\end{equation}
(The asymmetry arises from the odd order derivative of $u$ and to get the \textbf{(WEP)} from \textbf{(EP)} we assumed $\alpha = 0$) 


\begin{theorem}[Existence and Uniqueness of the (WEP)]
	  Given a bounded domain $\Omega$ with pure Dirichlet condition, and:
	  \begin{enumerate}
	  	\item $A \in L^\infty$ and $A$ is uniformly elliptic
	  	\item $b \in W^{1, \infty}$, $c \in L^\infty$ and satisfies:
	  	\begin{align*}
	  		c(x) - \frac{1}{2}div(b)(x) \geq 0 \qquad \text{a.e. in $\Omega$}
	  	\end{align*}
	  \end{enumerate}
	  then 
	  \begin{align*}
	  	\forall f \in L^2(\Omega), g \in H^1(\Omega)\exists!u \in H^1(\Omega): \textbf{u solves the (WEP)}
	  \end{align*}
\end{theorem}
\begin{proof}
Boundedness follows immediately by Cauchy-Schwartz. As for the coercivity:
\begin{align*}
	B(u, u) &= \int_\Omega \nabla u A \nabla u dx + \int_\Omega u  b(x) \cdot  \nabla u dx + \int_\Omega c(x) u^2 dx \\
	&\geq \int_\Omega a_0 |\nabla u|^2 + b \cdot \nabla (\frac{1}{2} u^2) + cu^2 dx\\
	&= \int_\Omega a_0 |\nabla u|^2 + (c -\frac{1}{2}div(b))u^2 dx \\
	&\geq a_0 \|\nabla u\|^2_2 \\
	&\overset{(1)}{\geq} c_0 \|u\|_{H_1} 
\end{align*}
\end{proof}

\subsubsection*{Important Elliptic System}
The weak form of the elasticity problem:
\begin{equation}
\left\{
\begin{aligned}
	&B(u, \varphi) := \int_\Omega \lambda (\nabla \cdot u)(\nabla \cdot \varphi) + \mu (\nabla u + \nabla u^T) :(\nabla \varphi + \nabla \varphi^T)dx = \int_\Omega f \varphi dx \quad \forall \varphi \in H^1_0(\Omega) \\
	&u - g \in H^1_0(\Omega)
\end{aligned}	
\right.
\end{equation}
Notation: $\nabla \cdot u$ : divergence; $A:B := \sum_{i,j} A_{ij}B_{ij}$



\subsection{Regularity of Solutions}
In this section we would like to explore, under which condition does \textbf{(WEP)} possess a more regular solution. \par 
Consider the most simple case of the Poisson equation. By linear Elliptic theory we know the solution operator is the inverse Laplacian 
\begin{align*}
	(-\Delta)^{-1}: H^{-1} \rightarrow H_0^1(\Omega)
\end{align*}

We would like to require a more regular $ f \in L^2(\Omega)$. We know that the weak derivative of $f$ exists and $\partial_j f \in H^{-1}$. Define $v = \partial_j u$, we have the elliptic equation:
\begin{align*}
	\Delta v = \partial_j f
\end{align*} 

If $v = 0$ on $\partial\Omega$ we would have $v \in H_0^1$ hence $u \in H^2_0$. But $v$ clearly doesn't necessarily satisfies the Dirichlet boundary condition. We now explore step by step the $L^2$-regularity with different boundary condition assumptions.\par

\subsubsection{Regularity on a cube}
\begin{theorem}
	Given a cube $ \Omega = [0, l_1] \times \dots [0, l_n]$ on $\mathbb{R}^n$, if $f \in L^2(\Omega)$, then we have
	\begin{align*}
		\|u\|_{H^2} \le C \|f\|_{L^2}
	\end{align*}
	where the constant $C$ independent of $u$ and $f$. 
\end{theorem}
\begin{proof}(via spectral theorem)\par 
We know $L^2(\Omega)$ has the orthonormal basis consisting of countable eigenfunction. 
\begin{align*}
	\varphi_k(x) = \sin(\frac{\pi k_1}{l_1}) \dots \sin (\frac{\pi k_n}{l_n}) \qquad k \in \mathbb{N}^n
\end{align*}
For $f \in L^2(\Omega)$ we have 
\begin{align*}
	f = \sum_{k \in \mathbb{N}^n} a_k \varphi_k  \quad (a_k)_k \in l^2(\mathbb{R}^n) 
\end{align*}
By the PDE
\begin{align*}
	\Delta u = f
\end{align*}
we can calculate $u$ explicitly:
\begin{align*}
	u = \sum_{k \in \mathbb{N}^n} b_k \varphi_k
\end{align*}
with 
\begin{align*}
	b_k = \frac{a_k}{\sum_{i = 1}^n (\frac{\pi k_i}{l_i})^2}
\end{align*}
we have the estimate 
\begin{align*}
	|k|^2|b_k| \le C |a_k|
\end{align*}
where the constant $C$ depends on $l$. Notice $(b_k)_k \in l^2(\mathbb{R}^n$ as well. We calculate the second derivative of $u$:
\begin{align*}
	\partial_1 \partial_2 u = \sum_{k \in \mathbb{N}^n} b_k \frac{\pi k_1 \pi k_2}{l_1 l_2} \cos(\frac{\pi k_1}{l_1}) \cos(\frac{\pi k_2}{l_2}) \dots
\end{align*} 
Observe
\begin{align*}
	\sum_{k \in \mathbb{N}^n}  |b_k \frac{\pi k_1 \pi k_2}{l_1 l_2}|^2 < \infty
\end{align*}
we have 
\begin{align*}
	\|u\|_{H^2} \le C \|f\|_{L^2}
\end{align*}
\end{proof}


\begin{theorem}
 	Given $k \in \mathbb{N}_0$, $\Omega \subset \mathbb{R}^d$ with $\partial \Omega \in C^{k+2}(\Omega)$. Moreover, assuming we have the data:
 	\begin{itemize}
 		\item $A \in C^{k+1}(\bar{\Omega}, \mathbb{R}^{d \times d})$, $A$ uniformly elliptic(i.e. $\exists a_0 \geq 0 \forall \xi \in \mathbb{R}^d: |\xi| A(x) |\xi| \geq a_0 \|\xi\|^2$)
 		\item $b \in C^k(\bar{\Omega}, \mathbb{R}^d)$
 		\item $c \in C^k(\bar{\Omega})$
 		\item $f \in H^k(\Omega)$
 		\item $g \in H^{k+2}(\Omega)$
 	\end{itemize}
 	Then the weak solution $u \in H^1(\Omega)$ of the pure Dirichlet problem satisfies 
 	\begin{align*}
 		u \in H^{k+2}(\Omega)
 	\end{align*}
 \end{theorem}
\begin{remark}
Here we prove the simple case where $k = 0$ and $\Omega = \mathbb{R}^d$, for the general case see Schweizer.
\begin{proof}
 	By the proposition, in order to prove $u \in H^2(\Omega)$ we only need to prove the difference quotient of $\nabla u$ is bounded.
 	Define the bilinear form:
 	\begin{align*}
 		\hat{B}(u, \varphi) := \int_\Omega \nabla u A \nabla \varphi dx + \int_\Omega a_0 u \varphi dx\\
 	\end{align*}
 	By the uniform ellipticity we get:
 	\begin{align*}
 		\hat{B}(u, u) \geq a_0\|u\|_{H^1}^2
 	\end{align*}
	By the weak form of \textbf{(WEP)} we get:
 	\begin{align*}
 		\hat{B}(u, v) &= \int_\Omega ((a_0 - c)u - b\nabla u)\varphi dx     + \int_\Omega f \varphi dx\\
 		&= \int_\Omega \underline{ [(a_0 -c)u - b\nabla u + f)]}\varphi dx\\
 		&:= \int_\Omega w \varphi dx\\
 		(1)&< \infty
 	\end{align*}
 	(1)Define the underlined part as $w$, since $u \in H^1, f \in L^2, \dots$ we get the  boundedness of $\hat{B}$ by Cauchy-Schwartz.
 	Inserting the difference quotient in the bilinear form:
 \begin{align*}
 	B(u, \partial_j^h u)
 \end{align*}
\begin{align*}
	a_0\|\partial^h_j u \|_{H^1} &\leq \hat{B}(\partial_h^j u, \partial_h^j u)\\
	&= \int_\Omega (a_0 - c) \partial^h_j u 
\end{align*}
\end{proof}	
\end{remark}



\subsection{Spectral Theory for elliptic operator}
This section we focus on the bilinear form:
\begin{align*}
	H^1_{\Gamma_D} \times H^1_{\Gamma_D} &\rightarrow \mathbb{R} \\
	B(u, v) &= \int_\Omega u A v + buv dx + \int_{\Gamma_N} c uv dx
\end{align*}
This weak formulation relates to the eigenvalue problem(EWP)
\begin{align*}
	\div(A \nabla\psi) + au &= \lambda \rho \psi \\
	\psi &= 0 \qquad  \psi \in \Gamma_D\\
	A \nabla \psi \cdot \nu + \alpha \psi &= 0 \qquad \psi \in \Gamma_N
\end{align*}

\begin{lemma}[Rellich's compact embedding]
	The embdding:
	\begin{align*}
		E: H^1_{\Gamma_D}(\Omega) &\hookrightarrow L^2(\Omega)\\
			u &\mapsto u
	\end{align*}
\end{lemma}
\begin{proof}
\end{proof}

\begin{lemma}[Abstract result on symmetric and compact operator]
	Given a Hilbert space $(H, \langle , \rangle)$ and \textbf{compact} and symmetric(self-adjoint) ($\langle Kf, g \rangle = \langle f, Kg   \rangle$) operator $K: H \rightarrow H$. Then there exists an ONS $\{ \varphi_k | k \in A, A \subset \mathbb{N}\}$ and eigenvalue $\{\lambda_k | \lambda_k \overset{k \rightarrow \infty}{\rightarrow} 0\}$, such that $K$ admits the spectral representation:
	\begin{align*}
		\forall u \in H: Ku = \sum_{k} \lambda_k \langle u, \varphi_k\rangle \varphi_k
	\end{align*}
\end{lemma}
\begin{proof}
	\texttt{cf.linear functional Analysis}
\end{proof}

\begin{theorem}[The elliptic spectral theorem]
	Given a bounded domain $\Omega \subset \mathbb{R}^d$ with Lipschitz boundary $\partial \Omega$. And a \textbf{nice}(symmetric, bounded, coercive) bilinear form $B$, \textbf{a general scalar product} $S$. Then there exists a cONS $\{\psi_k | k \in \mathbb{N}\}$, and $\{(\lambda_k)_{k \in \mathbb{N}}: \lambda_k \overset{ k \rightarrow \infty}{\rightarrow} \infty, \lambda_k > 0 \}$(i.e the eigenpair $(\psi_k, \lambda_k)$), such that:
\begin{align*}
		\forall \varphi \in H^1_{\Gamma_D}(\Omega): B(\psi_k, \varphi) =\lambda_k S(\psi_k, \varphi) 
\end{align*}
Moreover, in the case of a trivial kernel the cONB is complete. 
\end{theorem}
\begin{remark}
The general scalar product is defined as:
\begin{align*}
	S: L^2(\Omega) \times L^2(\Omega) \rightarrow \mathbb{R}, 
	(u, v) \mapsto \int_\Omega \rho(x)u(x)v(x)dx , \quad \rho \in L^\infty(\Omega), \rho(x) \geq \rho_0 > 0
\end{align*}	
\end{remark}
\begin{proof}
We define the operator $K$ that finds the solution for the given data $f$:
\begin{align*}
	K: L^2(\Omega) &\rightarrow H^1_{\Gamma_N}(\Omega) \\
	Kf &= u :\forall \varphi\in H^1_{\Gamma_N}(\Omega):  B(u, \varphi) = S(f, \varphi) \
\end{align*}
by Lax-Milgram we know the linear operator is well defined.
\begin{itemize}
	\item \textbf{Compactness}
	By Rellich's we know the the embedding
	\begin{align*}
		E: H^1_{\Gamma_N} \rightarrow L^2
	\end{align*}
	is compact. 
	The composition:
	\begin{align*}
		E \circ K: L^2(\Omega) \overset{K}{\hookrightarrow} H^1_{\Gamma_N}(\Omega) \overset{E}{\hookrightarrow} L^2(\Omega)\\
		f \mapsto u = Kf  \mapsto u= Eu = Kf
	\end{align*}
	is also compact by the theorem that one of the operator is compact is sufficient for the composition to be compact.\par
	Without loss of generality we can define a compact operator:
	\begin{align*}
		K: L^2(\Omega) &\rightarrow L^2(\Omega)\\
		Kf &= u: \forall \varphi \in H^1_{\Gamma_N}(\Omega): B(u, \varphi) = S(f, \varphi)
	\end{align*}
	\item \textbf{Symmetry}
	To show the symmetry of $K$ in $L^2(\Omega)$:
	\begin{align*}
		\forall f, g \in L^2(\Omega):  S(Kf, g) = S(f, Kg)
	\end{align*}
	By the definition:
	\begin{align*}
		S(Kf, g) &= S(g, Kf) \\
		&= B(Kg, Kf)\\
		&= B(Kf, Kg)\\
		&= S(f, Kg)
	\end{align*}
	follows from the symmetry of $B$ and $S$. 
	\item \textbf{Completeness}
	We follow the arguments:\\
	The linear operator has trivial kernel\\
	$\Rightarrow$ the linear operator $K$ is bijective\\
	$\Rightarrow$ $K$ has a cONB in $H^1_{\Gamma_N}$\\
	$\Rightarrow$ $K$ has a cONB with $K \psi_k = \lambda_k \psi_k$ in $L^2$ by density.
\end{itemize}
By the \textbf{abstract elliptic spectral theory} we have:
\begin{align*}
	\exists \gamma_k \rightarrow 0: K\psi_k = \gamma_k \psi_k
\end{align*}
Plug the compact operator into the bilinear form, we obtained a cONS satisfying that for all $\varphi \in H^1_{\Gamma_N}(\Omega)$:
\begin{align*}
	B(\psi_k, \varphi) &= \frac{1}{\gamma_k} B(K\psi_k, \varphi)\\
	(1) &= \lambda_k S(\psi_k, \varphi)
\end{align*}
where in (1) we define $\lambda_k := \frac{1}{\gamma_k}$. With $\gamma_k \rightarrow 0$ we get $\lambda_k \rightarrow \infty$.\\
For the positivity of $\lambda_k$ we observe:
\begin{align*}
	\lambda_k &= \frac{1}{\gamma_k}\\
	\gamma_k &= \gamma_k S(\psi_k, \psi_k) \\
	&= S(\psi_k, \gamma_k \psi_k)\\
	&= B(K\psi_k, \gamma_k \psi_k)\\
	&= B(K \psi_k, K \psi_k)\\
	(1)&>0 
\end{align*}
where (1) follows from the coercivity of $B$.
\end{proof}


\subsubsection*{Solve the Elliptic Problem with Spectral Representation}
Given $f \in L^2(\Omega)$, we have
\begin{align*}
	\forall \varphi \in H^1_{\Gamma_D}(\Omega): B(u, \varphi) = \int_\Omega f \varphi dx
\end{align*}
The last theorem provides a system of eigenpairs in the Hilbert space $\{(\psi_k, \lambda_k)\}_{k \in \mathbb{N}}$ , our goal is to find the spectral representation of $u$:
\begin{align*}
	u = \sum u_n \psi_n
\end{align*}
Plug in the spectral representation of the bilinear form and u:
\begin{align*}
	\int_\Omega f\psi_k dx =B(u, \psi_k) &= B(\sum u_n \psi_n, \psi_k)\\
	&= \sum u_n B(\psi_n, \psi_k)\\
	&= \sum u_n \lambda_n S(\psi_n, \psi_k) \\
	&= u_k \lambda_k
\end{align*}
We get the coefficient:
\begin{align*}
	u_k = \frac{\int_\Omega f \psi_k dx}{\lambda_k}, \quad 
	u= \sum \frac{1}{\lambda_k}( \int_\Omega f \psi dx ) \psi_k
\end{align*}
We still need to show that the sum converges in $H^1_{\Gamma_D}$.

\subsubsection*{Connection to the Green's function}
\subsubsection*{Connection to the Poincare's Friedrich's inequality}
The Poincare-Friedrich's inequality states:
\begin{align*}
	\exists C_{PW}>0: \| f\|_{L^2} \le C_{PW} \|\nabla f\|_{L^2}
\end{align*}




\newpage
\section{The Parabolic Equation}
\subsection{The Gronwall's Inequality}
\begin{theorem}
	
\end{theorem}


\subsection{Interpolation in Bochner Space}
 \subsubsection*{The Gelfand triple}
 The "$\hookrightarrow$" means that there is a \textbf{continuous inclusion} from the left side to the right side. Given a Hilbert space $H$, and $V$ its dense, reflexive subset, we say $(V, H, V^*)$ is a \textbf{Gelfand triple} if
 \begin{enumerate}
 	\item $V \subset H \subset V^*$. 
 	\item There exists a \textbf{consistent structure}.
 	\item we have the compact embedding $V \overset{c}{\hookrightarrow} H \overset{c}{\hookrightarrow} V^*$. 
 \end{enumerate}
 where the consistent structure means:
 \begin{align*}
 	\forall v \in V \forall h \in H: \langle u, v \rangle_{V^*, V} = \langle u, v \rangle_{H}
 \end{align*}

 in general, the second inclusion map is not injective, but the \textbf{Gelfand triple} states that if $K$ is a compact subset of $H$, then this inclusion is also injective. \par 
 
 From now we observe the Gelfand's triple
 \begin{align*}
 	V = H_0^1(\Omega) \quad H = L^2(\Omega) \quad V^* = H^{-1}(\Omega)
 \end{align*}
 
 \begin{lemma}[Embedding in to $C^0(0, T; H)$]
 	Given a Gelfand triple above and $u \in L^2(0, T; H)$.  Then we have:
 	\begin{enumerate}
 		\item $u \in L^2(0, T; V)$ and $\partial_t u \in L^2(0, T; V^*)$
 		\item \textbf{Continuity to $H$: }$ u \in C^0([0, T], H)$
 		\item \textbf{Measurability: }$t \mapsto \|u(t)\|_H^2 \in W^{1,1}([0, T])$
 		\item \textbf{time-derivative of the norm: }$\frac{d}{dt} \|u(t)\|_H^2 = 2 \langle E \partial_t u(t), u(t)\rangle_{V^*, V} $
 	\end{enumerate}
 \end{lemma}
 \begin{proof}\quad\\
 \textbf{1).} trivial by the definition of $H^1_0$. \\
 \textbf{2).} we define $u^\delta$, $u^\eta$ 2 mollifier of $u$, observe:
 \begin{align*}
 	\frac{d}{dt}\|u^\eta - u^\delta\|_H = 2 \langle \partial_t u^\eta - \partial_t u^\delta, u^\eta - u^\delta \rangle_{H} = \langle \partial_tu^\delta - \partial_t u^\eta, u^\delta - u^\eta \rangle_{V^*, V}
 \end{align*}
 By the fundamental lemma of calculus in Bochner space:
 \begin{align*}
 	\|u^\eta(t) - u^\delta(t)\|_H^2 &= \|u^\eta(s) - u^\delta(s)\|_H^2 + \int_s^t 2 \langle u_t^\eta(\tau) - u_t^\delta(\tau), u^\eta(\tau) - u^\delta(\tau) \rangle_{V^*, V} d\tau\\
 	&\le \|u^\eta(2) - u^\delta(s)\|_H^2 + \|\partial_t u^\eta - \partial_t u^\delta \|_{V^*} \|u^\eta - u^\delta\|_V
 \end{align*}
 Since $u^\eta, u^\delta \in C^\infty(0, T, H)$, and the Frechet derivative exists by continuous differentiability, therefore the second term vanishes. We have
 \begin{align*}
 	\sup_{|t - s| \le \delta}\|u(t) - u(s)\|_H & = \sup_{|t - s| \le \delta} \lim_{\eta, \delta \rightarrow 0}\|u^\eta(t) - u^\delta(s)\|_H  \\
 	& = \lim_{\eta, \delta \rightarrow 0} \sup_{|t -s| \le \delta}
 \end{align*}
 \textbf{3):} Follows from \textbf{3)} and $u \in L^2(0, T; V)$.\\
 \textbf{4):} Observe 
 \begin{align*}
 	\underbrace{\|u_\epsilon(t_2)\|_H^2}_{\rightarrow \|u(t_2)\|^2_H} - \|u_\epsilon(t_1)\|_H^2 
 	&\overset{DUIS}{=} \int_{t_1}^{t_2} 2\langle \partial_t u_\epsilon(t), u_\epsilon(t) \rangle_H dt\\
 	& ^= \int_{t_1}^{t_2} 2 \langle \partial_t u_\epsilon(t), u_\epsilon(t) \rangle_{V^*,V} dt 
 \end{align*}
 since $u_\epsilon(t) \overset{V}{\rightarrow} u(t)$,  $\partial_t u_\epsilon(t) \overset{V^*}{\rightarrow}\partial_t u(t)$ a.e. up to a subsequence. By the boundedness of $\|\partial_t u(t)\|_{V^*}$ we get the Frechet derivative.  
 \end{proof}
 \begin{remark}
 Here we used the Poincare's lemma explicitly, 	
 \end{remark}

 
 \begin{lemma}[Complete ONS in $V, H, V^*$]
 	Let $(V, H, V^*)$ be a Gelfand's triple. Then there exists 
 	\begin{align*}
 		\{(\lambda_k, \varphi_k)\}_{k \in \mathbb{N}} \subset \mathbb{R} \times V
 	\end{align*}
 	with
 	\begin{align*}
 		\langle \varphi_k, \varphi_l \rangle_V = \delta_{kl}, \quad
 		\langle \varphi_k, \varphi_l \rangle_H = \lambda_k \delta_{kl}, \quad
 		\langle E\varphi_k, E\varphi_l\rangle_{V^*} = \lambda_k^2 \delta_{kl}
 	\end{align*}
 	We get cONB in each of the triple:
 	\begin{align*}
 		\{\varphi_k\}_{k \in \mathbb{N}} \subset V, \quad 
 		\{\frac{1}{\sqrt{\lambda_k}} \varphi_k\}_{k \in \mathbb{N}} \subset H, \quad
 		\{\frac{1}{\lambda_k}E\varphi_k\}_{k \in \mathbb{N}} \subset V^*
 	\end{align*} 
 \end{lemma}
 \begin{proof}
 	Use the Gelfand triple to define a compact and symmetric operator $K: V^* \rightarrow V^*$ and apply \textbf{spectral theorem for compact and symmetric operator}. \quad\\
 	Define:
 	\begin{align*}
 		K: &V^* \rightarrow V^*\\
 		&V^* \overset{J_v}{\approxeq}V \overset{I:\texttt{cpt dense}}{ \hookrightarrow} H \overset{E}{\approxeq} H^*  \overset{I^*}{\hookrightarrow} V^* 
 	\end{align*}
 	1. show $K$ is compact and symmetric, get $\{\xi_k, k \in \mathbb{N}\}$ cONB of $V^*$   \\
 	2. Observe $\varphi_k = J_v \xi_k$, we have :
 	\begin{align*}
 		\langle \varphi_k, \varphi_l \rangle &= \delta_{kl}\\
 		\langle I \varphi_k, Iv \rangle_H &= \langle K \xi_k, v \rangle_{V^*, V} = \langle \lambda_k \xi_k, v \rangle_{V^*, V} = \lambda_k \langle \varphi_k, v \rangle_V\\
 		\Rightarrow 
 		\langle \varphi_k, \varphi_k \rangle &= \lambda_k \delta_{lk}
 	\end{align*}
 \end{proof}
 
 \begin{definition}[weak solution of PaE]
 	$u: [0, T] \rightarrow V$ is a weak solution to the parabolic PDE if:
 	\begin{enumerate}
 		\item $u \in L^2(0, T; V)$
 		\item $E \partial_t u \in L^2(0, T; V^*)$
 		\item $\forall v \in V: \langle E \partial_t u + L u , v \rangle_{V^*, V} = \langle f(t), v \rangle_{V^*, V}$ in $L^2([0, T]; \mathbb{R})$
 		\item $u(0) = u_0 \in H$. 
 	\end{enumerate}
 \end{definition}
 \begin{remark}
 \textbf{3)} is the \textbf{operator formulation}, we have the equivalent \textbf{weak formulation:}
 \begin{align*}
 	\int_0^T \langle E\partial_t u + Lu, v\rangle_{V^*, V}  = \int_0^T \langle f, v \rangle_{V^*, V} dt
 \end{align*}	
 \end{remark}

\subsection{Energy Method of the Parabolic Equation }
Given a connected domain with Lipschitz boundary $\Omega$, $u: [0, T] \times \Omega \rightarrow \mathbb{K}$,  we consider equations of the form: \textbf{(PaE)}
 \begin{equation}  
 \begin{aligned}
 	\alpha(x) \partial_t u - \div(A(x)\nabla u) + b \cdot \nabla u + c u = f(t, x)   &\texttt{ on }(0, T) \times \Omega \\
 	u(0, \cdot) = u_0(x) &\texttt{ on } \{0\} \times \Omega \\
 	u(t, x) = 0 &\texttt{ on }(0, T) \times \Gamma_{dir}\\
 	A(x) u\cdot \nu = 0 &\texttt{ on } (0, T) \times \Gamma_{neu}
 \end{aligned}	
 \end{equation}
 We would like to transform it into the weak formulation. After integrating by parts w.r.t. $x$ we get the a bilinear form: $B: V \times V \rightarrow \mathbb{K}$ \textbf{not symmetric}:
 \begin{align*}
 	&B(u, v) = \int_\Omega \nabla u A(x) \nabla v + u b \cdot \nabla v + ucv dx 
 \end{align*}
 where $V = H_0^1(\Omega) \overset{dense}{\subset} L^2(\Omega) = H$. And the Lax-Milgram theorem gives a bounded inverse in $V$. \\
 \textbf{Remark:} $H^1_0$ is a Hilbert space.
 \begin{align*}
 	\texttt{(Lax) }
 	&L :  \langle Lu , v \rangle_{V^*, V} = B(u, v)
 \end{align*}
 Define the map:
 \begin{equation*}
 E: = \left\{
 \begin{aligned}
 	H &\rightarrow H^*\\
 	u &\mapsto (\alpha u|\cdot)_H
 \end{aligned}
 \right.
 \end{equation*}
 with the weight norm
 \begin{align*}
 	 (u|v)_H := \langle Eu, v \rangle_{H^*, H}
 \end{align*}
 and $L$ the Lax-Milgram map
 \begin{align*}
 	L: V \rightarrow V^*, B(u, v) = \langle Lu, v \rangle_{V^*, V}
 \end{align*}
 we define:
 \begin{align*}
 	H= L^2(\Omega), \quad V:= H_{\Gamma_{dir}}^1 (\Omega), \quad V^* = H_{\Gamma_{dir}}^1 (\Omega)^*
 \end{align*}
 Observe by the norm equivalence by Poincare.
 \begin{align*}
 	B(u, u) 
 	&\ge \alpha_{min}\|\nabla u \|_{H}^2 - \underbrace{ b_{max}\|u\|_H \|\nabla u \|_H}_{\le c^1\|u\|_H^2 + c^2 \|\nabla u\|^2_H } - c_{max}\|u\|_H^2 \\
 	&\ge C_1 \| \nabla u\|_H^2  - C_2 \|u\|_H^2\\
 	&\ge C_1 \|u\|_V^2 - C_2 \|u\|_H^2
 \end{align*}
 
 


 
 \begin{theorem}[Existence of solutions]
 		Let $H, V, E, L$ defined as above, then the problem:
 		\begin{align*}
 			E \partial_t u + Lu &= f(t) \in V^* \\
 			u(0) &= u_0 \in H
 		\end{align*}
 		admits an unique weak solution. 
 \end{theorem}
 \begin{remark}
 The solution mapping:
 \begin{align*}
 	\mathcal{K}: H \times L^2(0, T; V^*) &\rightarrow L^2(0,T; V) \cap H^1(0, T; V^*)\\
 	(u_0, f ) &\mapsto u
 \end{align*}	
 is a topological isomorphism. i.e. it admits a bounded inverse. 
 \end{remark}
 \begin{proof}\quad\\
 \textbf{1). Finite dimensional approximation:}
 Let $V_n = \spa \{\varphi_1, \dots, \varphi_n \} \subset V $ with $\{\varphi_k\}_{k \in \mathbb{N}}$ the cONB defined above. $u_n:= \sum_{k = 1}^n \alpha_k(t) \varphi_k \in V_n.$ with $\alpha(t) \in H^1(0, T, \mathbb{R}^m)$, by the PDE we have
 \begin{align*}(*) \quad
 	\forall v_n \in V_n: \langle E \partial_t u_n(t), v_n \rangle_{V^*, V} + B(u_n(t), v_n) = \langle f(t), v_n \rangle_{V^*, V} 
 \end{align*}
 For $\alpha := (\alpha_1, \dots, \alpha_n)$, we have the ODE:
 \begin{align*}
 	&\dot{\alpha}(t) = B^{(m)} \alpha(t) + f^{(m)}(t) \quad  \forall 1 \le m \le  n \\
 	&\alpha(0) = (\alpha_1(0),  \dots, \alpha_n(0))
 \end{align*}
 the ODE has an unique solution by the linearity of $B$ (Lipschitz continuous). We hence proved the existence of finite dimensional solution in $L^2(0, T, V) $
 \\
 \textbf{2). A priori estimate:}
 claim:
 \begin{align*}
 	(1):\|u_n\|_{L^2(0, T; V)}< \infty,\quad
 	(2):\|E \partial_t u_n\|_{L^2(0, T; V^*)} < \infty
 \end{align*}
 For \textbf{2):} we have
 \begin{align*}
 	\| \partial_t u_n\|_{V^*}  &= \|Lu_n - f_n\|_{V^*}\\
 	&\le \|Lu_n\|_{V^*} + \|f_n\|_{V^*}\\
 	&\le  \sup_{\|v\|_V = 1}B(u_n, v)  + \|f\|_{V^*}\\
 	&\le C \|u_n\|_V + \|f\|_{V^*}\\
 	&\le C \|u\|_V + \|f\|_{V^*}
 \end{align*}
 therefore by the integral in Bochner space
 \begin{align*}
 	\|\partial_t u_n(t)\|_{L^2(0, T; V^*)}^2  \le  \int_0^T \|\partial_t u(t)\|_{V^*}^2 dt < \infty
 \end{align*}
 \textbf{For 1).}
 we plug in $v_n = u_n$ and get 
 \begin{align*}
 	\underbrace{\langle \partial_t u_n, u_n \rangle_{V^*, V}}_{ = \frac{1}{2}\frac{d}{dt}\|u_n\|_H^2} + \underbrace{B(u_n, u_n)}_{\le C_1 \|u_n\|_V - C_2 \|u_n\|_H} &\le \|f_n(t)\|_{V^*}\|u_n\|_V\\
 	\frac{1}{2}\frac{d}{dt}\|u_n\|_H^2 + C_1 \|u_n\|_V^2 &\le C_2 \|u_n\|_H^2 + \|f\|_{V^*}\|u_n\|_V\\
 	\frac{1}{2}\frac{d}{dt} \|u_n\|_H^2 + C_1 \|u_n\|_V^2 &\le C_2 \|u_n\|_H^2 + b_1 \|f\|_{V^*}^2 + b_2 \|u_n\|_V^2
 \end{align*}
 define $\|u_n(t)\|_H^2 := g(t)$,  first consider
 \begin{align*}
 	 \frac{d}{dt} \|u_n\|_H^2 &\le C_2 \|u_n\|_H^2 + h(t)  \\ 
 	\dot{g}(t)  &\le C_2g(t) + \|f(t)\|_{V^*}
 \end{align*}
 By the Gronwall we have:
 \begin{align*}
 	\|u_n(t)\|_H^2 &\le e^{2C_2 t} \|u_0\|_H^2 + \int_0^t e^{2 C_2(t - s)}h(s)ds := \tilde{C}\\
 	\Rightarrow \|u_n\|_{C(0, T; H)} &\le \tilde{C}T
 \end{align*}
 up to a constant:
 \begin{align*}
 	\int_0^T \|u_n\|_V &\le \int_0^T \frac{d}{dt}\|u_n\|_H^2 + \|u_n\|_H^2 + \|f\|_{V^*} dt \\
 	&\le \|u_n(T)\|_H^2 - \|u_0\|_H^2 + \tilde{C}T + T\|f\|_{V^*}\\
 	&< \infty
 \end{align*}
 \\
 \textbf{3). Extraction of a subsequence}
 \\
 \textbf{4). Limit passage:}
 By reflexivity of $L^2(0, T;V)$ and $L^2(0, T; V^*)$ we extract the weak convergent subsequences:
 \begin{align*}
 \partial_t u_n \rightharpoonup w^*\\
 	u_n \rightharpoonup u^*
 \end{align*}
 For $\varphi \in C_C^\infty(0, T;V)$:
 \begin{align*}
 	\int_0^T \langle w^*, \varphi \rangle dt 
 	\leftarrow
 	\int_0^T \langle E\partial_t u_n, \varphi \rangle dt = -\int_0^T \langle E\partial_t \varphi, u_n \rangle dt \rightarrow -\int_0^T \langle E\partial_t\varphi, u^* \rangle dt
 \end{align*}
 therefore:
 \begin{align*}
 	E\partial_t u_* = w^* \quad a.e. 
 \end{align*}
 \textbf{5):Compatibility w.r.t. the initial conditon: } Since we have a continuous embedding by the last lemma, we have:
 \begin{align*}
 	u \in C^0([0, T], V)
 \end{align*}
 \quad\\
 \textbf{6):Uniqueness}:	
 \end{proof}
 
  

\section{The Hyperbolic Equation}
\subsection{Separation of variables}
\subsubsection{Transport Equation with separation of variables}
To simplify the problem we first consider the homogeneous case. Observe the 1-dimensional case
\begin{equation*}
\left\{
\begin{aligned}
	&u_t = a u_x \quad \texttt{on } [0, T] \times \mathbb{R} \\
	&u(0, x) = u_0(x)
\end{aligned}
\right.
\end{equation*}
Try the ansatz(seperation of variables) 
\begin{equation*}
	u(t, x) = \varphi(t) \psi(x)
\end{equation*}
then
\begin{align*}
	\varphi'(t) \psi(x) &= a\varphi(t) \psi'(x)\\
	\Rightarrow \frac{\varphi'(t)}{\varphi(t)} &= a\frac{\psi'(x)}{\psi(x)} := \lambda (\textbf{Given } \varphi, \psi \neq 0)
\end{align*}
and have reduced the PDE into 2 ODEs:
\begin{equation*}
	\left\{
	\begin{aligned}
		&\varphi'(t) = \lambda \varphi(t)\\
		&\psi'(x) = \frac{\lambda}{a} \psi(x) 
	\end{aligned}
	\right.
\end{equation*}
The ODEs have the general solutions 
\begin{equation*}
	\begin{aligned}
		\varphi(t) = c e^{\lambda t}\\
		\psi(x) = \tilde{c} e^{\frac{\lambda}{a} x}
	\end{aligned}
\end{equation*}
By superposition we get the solution has the form:
\begin{equation*}
	u(t, x) = \sum_{i = 1}^{k^*} c_i e^{\lambda_i(\frac{x}{a} + t)}
\end{equation*}
\textbf{Question:} Is it possible to match the initial condition:
\begin{align*}
	u(0, x) = u_0(x) = \sum_{i = 1}^\infty c_i e^{\lambda_i(\frac{x}{a})}
\end{align*}
It seems to be difficult. We make the \textbf{ansatz}:
\begin{equation*}x
	u(t, x) = G(x + at)
\end{equation*}
It satisfies 
\begin{align*}
	G_t = a G_x
\end{align*}
check the initial condition:
\begin{equation*}
	G(0, x) = u_0(x) = u_0(x + a\cdot 0)
\end{equation*}
We guess $u_0(x + at)$ is a solution, (check if it satisfies the PDE):
\begin{equation*}
	G(t,x) = u_0(x + at)
\end{equation*}
We have a solution provided $u_0 \in C^1$
\\
Yet we haven't showed the \textbf{uniqueness}. For the inhomogeneous 1-dimensional transport equation
\begin{equation}\label{TE}
\left\{
	\begin{aligned}
		&u_t = a \cdot u_x + f(t, x)\\
		&u(0, x) = u_0(x)
	\end{aligned}
\right.
\end{equation}
we have the solution of the form(by parametric integral)
\begin{equation} \label{TE: solution}
	u(t, x) = u_0(x + at) + \int_0^t f(s, x + a(t -s))ds
\end{equation}

\subsection{The Energy Estimate}

\subsection{well-posedness}
\begin{definition}
	We say the linear system is well-posed in $L^2(\mathbb{R}^d, \mathbb{R}^n)$ if
	\begin{enumerate}
		\item $\forall u_0 \in L^2(\mathbb{R}) \exists ! u \in C([0, T], L^2(\mathbb{R}^d, \mathbb{R}^n))$ s.t. $\forall \varphi \in S(\mathbb{R} \times \mathbb{R}^d ; \mathbb{R}^n))$ with $\varphi(T, \cdot) = 0$:
		\begin{align*}
			\int_{\mathbb{R}^d}\int_0^T u \partial_t \varphi   +  u \sum_{j = 1}^d A_j \partial_{x_j} \varphi dx dt + \int_{\mathbb{R}^d} u_0 \varphi(0, \cdot) dx = 0
		\end{align*}
		\item The linear solution map:
		\begin{align*}
			u_0 \mapsto u, \quad L^2(\mathbb{R}^d, \mathbb{R}^n) \rightarrow C([0, T], L^2)
		\end{align*}
		is continuous, i.e., 
		\begin{align*}
			\|u\|_{C([0, T], L^2)} \le C_T \|u_0\|_{L^2}
		\end{align*}
	\end{enumerate}
\end{definition} 


\subsection{Evolution process with semigroup}

\section{The Calculus of Variation}

\subsection{The Energy functional}




Given $g \in Lip(\mathbb{R}^n)$ and particular $H \in C^2(\mathbb{R}^n)$. We aim at global solution to the \textbf{Hamilton-Jacobi(the scalar conservation law)}:
\begin{align*}
	&H(Du, x) + u_t = 0 \texttt{ in } \mathbb{R}^n \times \mathbb{R}(\texttt{space time cylinder)}\\
	&u = g \qquad \qquad \quad \texttt{ on } \mathbb{R}^n \times \{0\} = \partial\Omega
\end{align*}
\quad\\
In the calculus of variation problem, we try to solve:
\begin{align*}
	\int_\Omega L(Du, u, x) dx
\end{align*}
where $L$ is the \textbf{Lagrangian.} and is the  \textbf{convex conjugate} of the \textbf{Hamiltonian}:
\subsection{The Euler Lagrange Equation}
\begin{definition}[the COV problem]
	Fix time interval $[0, T]$ for $0 < T < \infty$ and BC at $t = 0$ and $t = \Gamma$ by $a, b \in \mathbb{R}^n$. Define the admissible functions:
	\begin{align*}
		A = \bigg\{ w \in C^2([0, T], \mathbb{R}^n): w(0) = a, w(T) = b \bigg\}
	\end{align*}
	and minimize:
	\begin{align*}
		I(w):= \int_0^T L(\dot{w}(t), w(t)) dt \quad \texttt{for } w \in A
	\end{align*}
\end{definition}

Notice the Lagrangian in this context admits arguments of the trajectory and the speed. $u$ doesn't appear in this case. 

\begin{theorem}[Euler-Lagrange equation of energy $I$]
	Suppose the trajactory $\gamma \in A$ minimises $I$ in $A$, then 
	\begin{align*}
		I(\gamma) = \inf I(A) = \min (A) \in \mathbb{R}
	\end{align*}
	then we have for $x = x(t)$:
	\begin{align*}
		-\frac{d}{d t}\bigg(L_{\dot{x}}(x, \dot{x})\bigg) + L_x(x, \dot{x}) = 0
	\end{align*}
	for all $0 < t< T$. 
\end{theorem}
\begin{remark}
in the absence of further conditions we don't know that there exists a minimizer. 
\end{remark}
\begin{proof}
	For $\varphi \in D(0, T; \mathbb{R}^n) = D(0, T)^n$ with
	\begin{align*}
		D(0, T)^n:=   C^\infty_c([0, T])
		\quad(\texttt{not ness.analytic}).
	\end{align*} 
	Since function in $D(0, T)^n$ admits the vanishing boundary we have
	\begin{align*}
		\forall s \in \mathbb{R}: \gamma + s \varphi \in A.
	\end{align*}
	And $I(\gamma + s \varphi)$ is smooth $s \in \mathbb{R}$ with a global minimum $s = 0$. i.e.
	\begin{align*}
		0&= \frac{\partial}{\partial s }I( \gamma + s \varphi)|_{s = 0}\\
		&= \int_0^T \frac{\partial}{\partial s }L (\dot{\gamma} + s \dot{\varphi}, \gamma + s \varphi) dt\\
		&= \int_0^T L_{\dot{y}}()\cdot  \dot{\varphi} + L_y \cdot  \varphi dt\\
		\texttt{(integration by parts) }
		&= \int_0^T -  \frac{d}{dt} L_{\dot{y}} \varphi + L_y \varphi dt + L_{\dot{y}} \varphi |_a^b\\
	\end{align*}
	by the fundamental lemma of calculus of variation we obtain the result. 
\end{proof}
\begin{remark}
The Euler-Lagrange equation 
\begin{equation}
	\frac{d}{dt}L_{y'} = L_y
\end{equation}
is similar to the Newton's 2nd law
\begin{equation}
	\frac{d}{dt} p = F
\end{equation}
\end{remark}


\subsubsection{The Hamiltonian Formalism}
\begin{definition}[Hamiltonian]
Suppose $\forall x, p \in \mathbb{R}^n$, 
\begin{align*}
	p= L_q(x, q) \texttt{ (generalised momentum)}
\end{align*}
can by uniquely solved for $q$ with a smooth function 
\begin{align*}
	q = \textbf{q}(p, x) \texttt{ (generalised velocity)}
\end{align*}
then we can define the \textbf{Hamiltonian} associated with Lagrangian $L$:
	\begin{align*}
		H(p, x) = p \cdot q - L(q(p, x) , x)
	\end{align*}
\end{definition} 

\begin{example}[point mechanics]
we denote $p$ the generalised momentum and $q$ the generalised velocity. $\phi$ the potential energy.
\begin{align*}
	L(q, x) = \frac{m}{2}|q|^2 - \phi(x)
\end{align*}	
with $m > 0$ mass, $\phi \in C^2(\mathbb{R}^n)$ potential with induced field $f(x):= -\nabla \phi(x)$. We check the assumption:
\begin{align*}
	p &= L_q(q, x) = mq \Leftrightarrow q = \frac{p}{m}
\end{align*}
so $Q(p, x) = \frac{p}{m}$. we have the Hamiltonian
\begin{align*}
	H(p, q, x) &= p \cdot q(x, p) - L(q(p, x), x)\\
	&=  \frac{|p|^2}{2 m }  - (\frac{m}{2}) (\frac{p}{m})^2 + \phi(x)\\
	&= \frac{m}{2} |q|^2 + \phi(x)
\end{align*}
In another form, we see that the Hamiltonian is the sum of the kinetic energy and the potential energy:
\begin{align*}
	H = \frac{m}{2} |q|^2 + \phi(x)
\end{align*}
For the Lagrangian of the classic mechanics, we get the corresponding \textbf{Eular-Lagrange equation}:
\begin{align*}
	&\frac{d}{dt}(mq) = D \phi(x)\\
	&ma = F
\end{align*}
which corresponds to the Newton's law.  \qed
\end{example}



The Euler-Lagrange equation can be solved by the Hamiltonian ODE. 

\begin{theorem}[Hamiltonian ODE]
Define the 
\begin{align*}
	p:= L_q(\dot{x}, x) \in C^1(0, T, \mathbb{R}^n)
\end{align*}
then we have \textbf{Hamiltonian ODE} satisfying Euler-Lagrange equation
\begin{align*}
\textbf{(a)}: \dot{p}(s) &= - H_x(p, x) \\
\textbf{(c)}: \dot{x}(s) &= H_p(p, x)	
\end{align*}
Moreover
\begin{align*}
	s \mapsto H(p(s), x(s)))
\end{align*}
is constant in time.
\end{theorem}
\begin{proof}\quad\\
\textbf{1):(c)}
\begin{align*}
	H_{p_i}(p, x) &= \frac{d}{dp_i}(p \cdot q(p, x) - L(q, x))\\
	&= \langle e_i, q \rangle + \langle p, \nabla_{p_i} q \rangle - \langle L_q, \nabla_{p_i} q \rangle \\
	&= q_i \\
	\texttt{ (by definition of $q$)}
	&= \dot{x}^i 
\end{align*}
\textbf{2): (a)}
\begin{align*}
	 H_{x_i}(p, x) &= \frac{d}{dx_i}(p \cdot q(p, x) - L(q, x))\\
	&=  \langle p,  \nabla_{x_i}q \rangle - (\langle L_q, \nabla_{x_i}q \rangle + L_{x_i}) \\
	&= \langle p,  \nabla_{x_i}q \rangle - \langle p, \nabla_{x_i} q \rangle - L_{x_i} \\
	\texttt{(Euler-Langrange)}&= \frac{d}{ds}(L_{q_i})  \\
	 \texttt{(duality) }&= \dot{p_i} 
\end{align*}
\textbf{3):} $H(p(s), x(s))$ is constant in time:
\begin{align*}
	H_s(p(s), x(s)) &:= \frac{d}{ds}(p(s) \cdot q(p, s) - L(q(p(s), s), x(s))\\
	&= \langle \nabla_s p , q \rangle  + \langle p, \nabla_s q \rangle 
	- \langle L_q , \nabla_s q \rangle + \langle L_x, \nabla_s x \rangle)\\
	&= \langle \nabla_s p , q \rangle  + \langle p, \nabla_s q \rangle
	-\langle L_q , \nabla_s q \rangle -  \langle \nabla_s L_q , q \rangle\\
	&= \langle \nabla_s p , q \rangle  + \langle p, \nabla_s q \rangle
	-\langle p , \nabla_s q \rangle - \langle \nabla_s p , q \rangle \\
	&= 0
\end{align*}

\end{proof}

\begin{remark}
Due to the general difficulties with the characteristics ODEs we circument in the examples below with arguments form the calculus of variations to deduce solutions of the Hamiliton ODEs. 	
\end{remark}


\subsection{Legendre transform, Hopf-Lax formula}
In order to satisfy the duality assumption, where the generalised velocity $q$ can be represented by the generalised momentum $p$ and the position $x$, we introduce the \textit{Legendre Transformation.} With the help of \textit{convex conjugate} we may get rid of $\dot{x}$ in $H(t, x, \dot{x}, p)$ and write the Hamiltonian as $H(x, p, q)$.   \\
\textbf{Assumption 1:}\\
\textbf{1)}$L$ doesn't depend on $x$ ($L_x \equiv 0)$. $L$ is convex\\
\textbf{2)}$L$ superlinear growth:
\begin{align*}
	\lim_{|q| \rightarrow \infty} \frac{L(q)}{|q|} = \infty
\end{align*}
\textbf{Duality Assumption of the Lagrangian}: $L \in C^2(\mathbb{R}^n \times \mathbb{R}^n)$ satisfy $\forall x , p, q \in \mathbb{R}^n$ with
\begin{align*}
	p &= L_q(q, x)\texttt{ (the generated momentum)}.
\end{align*}
$q$ can be represented by $p, x$, i.e., 
\begin{align*}
	q &= \textbf{q}(p, x).
\end{align*}
\begin{remark}\quad\\
1. \textbf{convex functions are continuous and has second derivative a.e.} by the \textbf{theorem of Alexandrov}.	\\
2. The convexity ensures the existence of the convex dual by applying the 1st variational equality. \\
3. The superlinearity:
\end{remark}
\quad\\
we henceforth have the modified Hamilton-Jacobi Equation:
\begin{align*}
	&H(Du) + u_t = 0 \texttt{ in } \mathbb{R}^n \times \mathbb{R}\\
	&u = g \qquad \qquad \quad \texttt{ on } \mathbb{R}^n \times \{0\} = \partial\Omega
\end{align*}

\begin{definition}[Legendre transformation]
	$\forall p \in \mathbb{R}^n$:
	\begin{align*}
		L^*(p):= \sup_{q \in \mathbb{R}^n} (p\cdot q - L(q))
	\end{align*}
\end{definition}

\begin{theorem}[Convex duality]
	$L$ convex, superlinear groweh, then
	\begin{align*}
		H&= L^* \texttt{ convex, superlinear growth}\\
		H^* &= L
	\end{align*}
\end{theorem}
\begin{proof}\quad\\
\textbf{1) $L^*$ is convex}: observe $p^*: q \rightarrow \langle p, q \rangle$ defines a functional on $\mathbb{R}^n$. 
\begin{align*}
	L^*(p^*) = \sup_{p \in \mathbb{R}^n} p^*(q) - L(q)
\end{align*}
	defines a conjugate function $L^*$ with argument of $p^*$, therefore convex. \\
\textbf{2): $L^*$ superlinear growth:} we normalize the argument, observe:
\begin{align*}
	L^*(p) &:= \sup_{q \in \mathbb{R}^n} \langle q, p \rangle - L(q)\\
	&\ge \lambda p - L(\frac{\lambda p}{|p|})   \texttt{ (choose $q = \frac{\lambda p}{|p|}$)}\\
	&\ge \lambda p - \max_{x \in B(0, \lambda)} L(x) 
\end{align*}
since we can choose any $\lambda$, the function is superlinear.\\
\textbf{3): The convex duality:} 
\begin{align*}
	\frac{d}{dp}\bigg( p \cdot q - L(q, x)\bigg) &\overset{!}{=} 0\\
	p &= L_q
\end{align*}
which is the definition of the Hamiltonian. 
\end{proof}

\begin{definition}[Hopf-Lax Formula]
	$g: \mathbb{R}^n \rightarrow \mathbb{R}$ Lipschitz initial condition 
	$\forall x \in \mathbb{R}^n \forall t >  0$:
	\begin{align*}
		u(x, t):= \min_{y \in \mathbb{R}^n} \bigg(t L(\frac{x - y}{t}) + g(y)\bigg)
	\end{align*}
\end{definition}

\begin{lemma}
	Given the $u$ defined above, and $L$ satisfies the \textbf{Assumption 1}, set
\begin{align*}
	u(t, x) := \inf_{w \in C^1} \bigg\{ \int_0^t L(\dot{w}(s) ) ds + g(w(0)) \bigg| \quad w(0) = y, w(t) = x\bigg\}
\end{align*}
we have:
\begin{align*}
		u(x, t):= \min_{y \in \mathbb{R}^n} \bigg(t L(\frac{x - y}{t}) + g(y)\bigg)
\end{align*}
\end{lemma}
\begin{proof}
First we prove the infimum:\\
"$\leq$": set $w(s) = y + \frac{s}{t}(x - y)$, we have
\begin{align*}
	u(t, x) \le \int_0^t L(\frac{x - y}{t})ds + g(y) = t L(\frac{x - y}{t}) + g(y)
\end{align*}
"$\geq$": By the averaging Jensen:
\begin{align*}
	 &L\bigg(\int_0^t \frac{1}{t} \dot{w}(s) ds \bigg) \overset{Jensen}{\le} \frac{1}{t} \int_0^t L(\dot{w}(s))ds  \\
	 \Rightarrow
	 & t L(\frac{x - y}{t}) + g(y) \le \int_0^t L(\dot{w}(s))ds + g(y) = u(t, x)
\end{align*}
\end{proof}

\begin{lemma}[A functional identity]
	\begin{align*}
		u(x, t) = \min_{y \in \mathbb{R}^n} ((t-s)L(\frac{x - y}{t - s}) + u(x, s)) 
	\end{align*}
	which means in order to calculate $u( \cdot, t)$ we can use $u( \cdot, s)$ as the initial condition on the remaining interval $[s, t]$. 
\end{lemma}
\begin{proof}
	convex tranformation
\end{proof}
	

\begin{theorem} Given the assumption above:($L$ convex and superlinear-gorowth, $g$ Lipschitz). We have:\\
 \textbf{1)}. The function $u$ defined by the \textbf{Hopf-Lax} formula is Lipschitz in $ Q =\mathbb{R}^n \times \mathbb{R}^+$ and therefore a.e. differentiable. \\
\textbf{2).} At the point $(x, t)$ where $u$ is differentiable, $u$ satisfies the HJE:
		\begin{align*}
			&u_t + H(Du) = 0\\
			&u(x, 0) = g(x)
		\end{align*}

\end{theorem}
\begin{proof}\quad\\
	\textbf{1): Lipschitz continuity(differentiability):} 
	\begin{align*}
		u(\hat{x}, t) - u(x, t) &= \min_z \{ t L(\frac{\hat{x} - z}{t}) + g(z)\} - \min_y\{tL(\frac{x - y}{t}) + g(y)\}\\
		(\texttt{choose } y = x - \hat{x} + z)
		&\le g(z) - g(z + \hat{x} - x)\\
		&\le C_g \|\hat{x} - x \|
	\end{align*}
	therefore $u$ is Lipschitz continuous with the Lipschitz constant of $g$ $C_g$.  \\
	For Lipschitz in $t$, first notice $u(x, 0) = g(x)$, we have:
	\begin{align*}
		u(x, t) &= \min_{y}\bigg\{tL(\frac{x - y}{t}) + g(y) - g(x) + g(x) \bigg\}\\
		&\ge  g(x) +  \min_y \bigg\{ tL(\frac{x - y}{t})  -C_g|x - y|\bigg\}\\
		(z = \frac{x - y}{t})
		&= g(x) - t\max_{z \sim (x, t)} \bigg\{ C_g |z| - L(z)\bigg\}\\
		&= g(x) - t\max_{w \in B(0, C_g)} \max_z \{ w \cdot z - L(z)\}\\
		&= g(x) - t H_{w \in B_{C_g}(0)}(w)
	\end{align*}
	Therefore we have a lower bound for $u(x, t) - g(x)$:
	\begin{align*}
		u(x, t) - g(x) \ge -t H_{w \in B_{c_g}(0)} (w)
	\end{align*}
	Moreover choosing $y= x$ we get an upper bound
	\begin{align*}
		u(x, t) &\le  tL(0) + g(x)\\
		\Rightarrow u(x, t) - g(x) &\le t L(0)
	\end{align*}
	Altogether
	\begin{align*}
		|u(x, t) - u(x, 0)| \le \max \bigg\{\max_{\|w\| \le C_g} H(w), L(0)\bigg\}t 
	\end{align*}
	The global Lipschitz continuity could be proved by the triangular inequality:
	\begin{align*}
		|u(x, \hat{t}) - u(x, t)| &=| u(x, \hat{t}) - u(x, 0) + u(x, 0) - u(x, t) |\\
		&\le \hat{t} \hat{C} - \hat{C} t\\
		&\le \hat{C}|\hat{t} - t|  
	\end{align*}
	By Radamacher theorem the function $u(x, t)$ is than also on $X \times \mathbb{R}$ a.e. differentiable. \\
	\textbf{2): Satisfies the Hamilton-Jacobi PDE:} \\
	"$u_t + H(D_x u) \le 0$:": \\
	Observe
	\begin{align*}
		u(x + hq, t + h)&= \min_y \biggl\{ h L(\frac{x + hq - y}{h}) + u(y, t)\biggr\} \\
		\texttt{(set $y = x$)} \quad 
		&\le h L(q) + u(x, t) \\
		\Rightarrow 
		\frac{u(x+ hq, t + h) - u(x, t)}{h} &\le L(q)\\
		\overset{h \rightarrow 0}{\Rightarrow} u_t + \langle D_x u, q \rangle &\le L(q) \quad \texttt{ (by differentiability})\\ 
		\\\Rightarrow
		u_t + H(D_x u) = u_t + L^*(D_x u) &=  u_t + \max_q \{ D_x u \cdot q - L(q)\}\le 0
	\end{align*}
	"$u_t + H(D_x u) \ge 0$".\\
	Define $h := t - s$. set $z := arg \min_y \{ t L (\frac{x - y}{t}) + g(y) \}$
	\begin{align*}
		u(x, t) - u(y, s) &=  tL(\frac{x - z}{t}) - g(z) - \min_{q \in \mathbb{R}^n} \{ s L(\frac{y - q}{s}) + g(q) \} \\
		\texttt{(choose $z = q$)}
		&\ge tL(\frac{x - z}{t}) - g(z) - sL(\frac{y - z}{t}) - g(z)\\
	(\texttt{set $z$ s.t. $\frac{x - z}{t} = \frac{y - z}{s}$})\\
		u(x, t) - u((1-\frac{h}{t}) x + \frac{h}{t}z  , t - h) &\geq h  L(\frac{x - z}{t})
	\end{align*}
	divide both sides by $h$ then taking the derivative we have
	\begin{align*}
		D_x u \cdot \frac{x - z}{t} + u_t \geq L(\frac{ x - z}{t})
	\end{align*}
	we have
	\begin{align*}
		u_t + H(D_x u) &= u_t + \max_q\{ D_x u \cdot q - L(q)\}\\
		&\ge 0 \quad (\texttt{choose }q = \frac{x - z }{t})
	\end{align*}
	For the boundary condition we get:
	\begin{align*}
		u(x, 0) = g(x)
	\end{align*}
	by the Hopf-Lax formula.
\end{proof}

Now we see that the solution to the Hamilton-Jacobi-Equation could be solved by a calculus of variation problem:
\begin{align*}
	u(t, x) := \inf_{w \in C^1} \bigg\{ \int_0^t L(\dot{w}(s) ) ds + g(w(0)) \bigg| \quad w(0) = y, w(t) = x\bigg\}
\end{align*}

We have a Lipschitz-continuous solution. Such solutions are in general not unique. For the solution to be \textbf{an unique weak solution}, we need to investigate the \textbf{semi-concavity}.

\subsubsection*{Uniqueness of the weak solution and semi-concavity}

\begin{definition}[weak solution of the Hamilton-Jacobi equation]
	A \underline{\textbf{Lipschitz continuous}} function $u : \mathbb{R}^n \times [0, \infty) \rightarrow \mathbb{R}$ is a weak solution of the initial-value problem if
	\begin{enumerate}
		\item $u(0, x) = g(x)$ 
		\item $u_t(x, t) + H(Du(x, t)) = 0$ a.e.
		\item $u(x + z, t) - 2u(x, t) + u(x - z, t) \le C(1 + \frac{1}{t})|z|^2 $
	\end{enumerate}
\end{definition}

\begin{theorem}
	If $H \in C^2(\Omega)$ and the underlying $L$ satisfies the assumptions and $g$ is Lipschitz. If either $H$ is uniformly convex, i.e. $D^2 H \geq \theta >0$ in $\mathbb{R}^n$ , or $g$ is semi-concave in the sense that :
		\begin{align*}
			\exists c_1 > 0 \forall x, y \in \mathbb{R}^n:
			g(x + y) - 2g(x) + g( x - y) \le c_1 |y|^2
		\end{align*}
		then $u$ is semi-concave, i.e.  $\exists c_2> 0 \forall x, y \in \mathbb{R}^n \forall t > 0$:
		\begin{align*}
			u(x+ y, t) - 2u(x, t) + u(x-y, t) \le c_2 (1 + \frac{1}{t} )|y|^2
		\end{align*}
		and the Hopf-Lax formula is an \textbf{unique} weak solution to the Hamilton-Jacobi problem. 
\end{theorem}
\begin{proof}\quad
Assuming $u$ and $\tilde{u}$ are 2 solutions. Define $w = u - \tilde{u}$
\begin{enumerate}
	\item First we prove that the \textbf{semi-concavity} implies the \textbf{uniqueness of weak solution.}
	\begin{align*}
	u_t(y, s) - \tilde{u}_t(y, s) 
		&= -H(Du(y, s)) + H(D\tilde{u}(y, s))\\
		\texttt{($H$ differentiable ) }
		&= -\int_0^1 \frac{d}{dr} H(rDu) +  \frac{d}{dr}H( (1-r)D\tilde{u})dr  \\
		\texttt{(chain rule) }
		&= \underbrace{ -\int_\Omega DH\bigg( rDu + (1 - r) D\tilde{u}\bigg)dr}_{:= b(y, s)} \underbrace{ (Du(y, s) - D\tilde{u}(y, s))}_{:= Dw(y, s)}
	\end{align*}
	\begin{align*}
		w_t + \textbf{b} \cdot Dw = 0 \quad a.e.
	\end{align*}
	\item
	\begin{align*}
		v_t + \textbf{b}_\epsilon \cdot Dv = 0 \quad a.e.
	\end{align*}
	\item Define $u^\epsilon$, $\tilde{u}^\epsilon$ the mollification of $u$ and $\tilde{u}$. Since the solution given by Hopf-Lax formula is Lipschitz we have:
	\begin{align*}
		\|D u^\epsilon\|_{op} \le Lip(u) \quad \|D\tilde{u}^\epsilon\|_{op} \le Lip(\tilde{u})
	\end{align*}
	and by the property of mollification:
	\begin{align*}
		Du^\epsilon \rightarrow Du \quad D\tilde{u}^\epsilon \rightarrow D \tilde{u}
	\end{align*}
	\item
	\begin{align*}
		v_t + div(v \textbf{b}_\epsilon) &= v_t + div(\textbf{b}_\epsilon)v + \textbf{b}_\epsilon \cdot D v \\
		&= div(\textbf{b}_\epsilon) v + (\textbf{b}_\epsilon - \textbf{b}) \cdot Dv \quad a.e.
	\end{align*}
	\item 
	\begin{align*}
		div(\textbf{b}_\epsilon)&= 
		\int_\Omega \Delta H(ru - (1 - r) \tilde{u}) dr \\
		(\texttt{semi-concavity of } u)\quad
		&\le C(1 + \frac{1}{s})
	\end{align*}
	\item Defne $B:= B(x_0, R(t_0 - t))$
	\begin{align*}
		e(t) = \int_B v(x, t) dx 
	\end{align*}
	\begin{align*}
		\dot{e}(t) \le C (1 + \frac{1}{t})e(t) \quad a.e. , 0 \le t \le t_0
	\end{align*}
	\item Choose $\phi(z)$ such that 
	\begin{align*}
		\phi(z) |_{|z| \le \epsilon [Lip(u) + Lip(\tilde{u})]} = 0
	\end{align*} 
	By Gronwall's lemma we have:
	\begin{align*}
		e(r) = 0 \quad 0< \epsilon < r < t
	\end{align*}
\end{enumerate}
\end{proof}

\subsection{Scalar conservation laws}
we observe the \textbf{IVP} of the \textbf{scalar conservation law}
\begin{equation*}
\left\{
\begin{aligned}
	u_t + F_x(u) = 0 &\texttt{ on } \Omega \times (0, \infty) \\
	u(x) = g(x) &\texttt{ on } \Omega \times \{t = 0\}
\end{aligned}
\right.
\end{equation*}
\subsubsection*{The traffic model}
Street: $\mathbb{R}$, $n = 1$, $t \geq 0$. Density $\rho(x, t):= \texttt{number of cars per length}$. 
\begin{align*}
	\int_a^b \rho(t, x) dx &= \texttt{ total number of cars at $t$. }\\
	v(x, t) &= \texttt{ velocity of cars at } (x, t)
\end{align*}
\textbf{1. Principle:} conservation of cars:
\begin{align*}
	\frac{d}{dt} \int_a^b \rho(x, t)dx &= \rho(a, b) v(a, t) - \rho(b, t) v(b, t)\\
	&= - \int_a^b (\rho v)_x dx 
\end{align*}
we have the simplified setting:
\begin{align*}
	v(\rho):= 1 - \rho
\end{align*}
In the traffic model we have the initial condition
\begin{align*}
	g(x) = & 1 \quad - l < x < 0\\
	&0 \quad else
\end{align*}
Finally we formulate the \textbf{Traffic PDE}
\begin{align*}
	&\rho_t + (\rho - \rho^2)_x = 0\\
	&\rho(0, x)= g(x)
\end{align*}
we have a shock(discontinuity) in the boundary condition. Achange the notation
\begin{align*}
	F(p, z, x) = p_2 - p_1(1 - 2 z)
\end{align*}
calculate the characteristic ODE:
\begin{align*}
	\dot{\gamma}(s) = F_p = 
	\begin{bmatrix}
		1 - 2z\\
		1
	\end{bmatrix}
\end{align*}
there are some region with jumps and some where the characteristic curve don't cross. 

\subsubsection*{Rakine-Hugonoite-Condition}
Observe the PDE of scalar conservation law. Observe $G: (t, x) \mapsto (u_t(t, \cdot), F(u(\cdot, x))$.  Apply the divergence theorem we have $\forall \varphi \in \mathcal{D}([0, \infty) \times \mathbb{R} \rightarrow \mathbb{R})$ we have:
\begin{align*}
	\int_0^\infty\int_\Omega div(G) \varphi dx ds
	&= -\int_0^\infty\int_\Omega u \varphi_t + F(u) \varphi_x dx ds + \int_\Omega (G \cdot \nu) \varphi dx|_{t = 0} \\
	(\nu = -1)
	&= -\int_0^\infty\int_\Omega u \varphi_t + F(u) \varphi_x dxds - \int_\Omega u(0, x) \varphi dx\\
	(\varphi \texttt{ compact supported})
	&= -\int_0^\infty\int_\Omega u \varphi_t + F(u) \varphi_x dxds 
\end{align*}

\begin{definition}[jump]
	At $x \in \Gamma$ 
	\begin{align*}
		[u]_\Gamma(x)&:= \lim_{\Omega_2 \ni x_2 \rightarrow x} u(x_2) - \lim_{\Omega_1 \ni x_1 \rightarrow x} u(x_1)\\
		[F(u)]_\Gamma(x) &:= \lim_{\Omega_2 \ni x_2 \rightarrow x}F(u(x_2)) - \lim_{\Omega_1 \ni x_1 \rightarrow x} F(u(x_1))
	\end{align*}
\end{definition}

\begin{proposition}[Rakine-Hugonoit condition]
	Recall $g \in L^\infty(\mathbb{R}), F \in C^1(\mathbb{R}), u \in C^1(\bar{\Omega})\cup C^1(\Omega_2)$, then $u$ is a weak solution to the IVP iff
	\begin{itemize}
		\item $u$ is a weak solution on the respective domains
		\item The Rakine-Hugonoit condition holds
		\begin{align*}
			\nu \perp ([F(u)]_\Gamma, [u]_\Gamma) \quad a.e. \quad x \in \Gamma
		\end{align*}
		(the condition is independent of the sign of the normal). 
	\end{itemize}
\end{proposition}
\begin{proof}
	\quad\\
	$\Rightarrow$:\\
	1)$u$ is a weak solution in $\Omega$, then $u$ is also a weak solution in any smaller open set in $\Omega$. By the regularisation we know $u|_{\Omega_i}$ is a classical solution of 
	\begin{align*}
		u_t + F(u)_x = 0
	\end{align*}
	2) Given $x_0 \in \Gamma, r_0$ s.t. $B(x_0, 2 r_0)$ and $0 < \epsilon < r_0$, $\varphi:= \eta_\epsilon * \chi_B(x_0, r_0) \in D(B(x_0, r_0 + \epsilon)) \subset D(\Omega)$. Since $u$ is a weak solution we have:
	\begin{align*}
		0 &= \int_0^\infty \int_{Q \cap \Omega} (u \varphi_t + F(u) \varphi_x) dQ\\
		&=  \sum_{j = 1}^2 \int_0^\infty \int_{Q \cap \Omega_j} (u\varphi_t + F(u) \varphi_x) dQ  \\
		&=  \sum_{j = 1}^2 \underbrace{ - \int_{\Omega_j \cap Q}  (u_t + F(u)_x)\varphi dQ}_{vanishes} + 
		\int_{\Gamma \cap \overline{\Omega_j}} (\nu_j \cdot (F(u), u)) \varphi dQ\\
		& =  \int_{\Gamma_j } [\nu_1(F(u_+ ) - F(u_-))+ \nu_2(u_+ - u_-) ] \varphi dQ
	\end{align*}
	all terms vanishes except for the last. Notice that $\nu_1, \nu_2$ admit the opposite signs. We therefore have at a neighborhoods of $x_0$:
	\begin{align*}
		\int_{B(x_0, r_0)}[ \nu_1 F(u_1) + \nu_2u_1 - (\nu_1 F(u_2) + \nu_2 u_2)] \varphi dx = 0
	\end{align*}
	By the fundamental lemma of calculus of variation we conclude the result.
	\quad\\
	$\Leftarrow$: trivial.
\end{proof}

\begin{example}[continue with the traffic model]
suppose the arc $\Gamma_{12}$ is parametrised time $t$ by 
\begin{align*}
	\gamma(t) = (s(t), t) \in \Gamma_{12} \quad t \ge 0
\end{align*}
The tangent vector along $\gamma$ is
\begin{align*}
	\dot{\gamma}(t) = (s'(t), 1) \Rightarrow 
	\tau(x):= \frac{(s'(t), 1)}{|\dot{\gamma}|}
\end{align*}
and the normal vector
\begin{align*}
	\nu(x) = \frac{(-1, s'(t))}{|\nu(x)|}
\end{align*}
By the RHC we have
\begin{align*}
	[F(u)]_\Gamma = s'(t) [u]_\Gamma
\end{align*}
In the traffic model we have $\Gamma_{12} = \{l\} \times [0, l], s(t) := -l$
\begin{align*}
	\sigma = 0
\end{align*}
The RHC condition is satisfied and we have a global weak solution. 
\end{example}


\subsubsection*{Entropy solution}
The intuition behind the entropy condition is that there cannot be characteristics coming from the discontinuity. We now derive a sufficient condition to rule out the non-physical solutions.
\quad\\
\textbf{The uniformly convex case}: Recall \textbf{(IVP)} of the conservation law with $F \in C^2(\mathbb{R})$ uniformly convex (i.e. $D^2f(x)$ strictly positive with a global lower bound $m > 0$). Function like this has the property
\begin{align*}
	(F^*)' = (F')^{-1}
\end{align*}
\begin{proof}
\begin{align*}
	F^*(p)&:= \sup_{q \in \mathbb{R}}(p q - F(q))\\
	\texttt{(max attained)}
	&= p \cdot (F')^{-1}(p) - F((F')^{-1}(p))\\
\end{align*}
differentiating both sides w.r.t. $p$ and the uniform convexity ensures that $(F'^{-1})'$ exists. 
\end{proof}

By the non-decreasing entropy we get the condition:
\begin{align*}
	F'(u_2) < \sigma < F'(u_1) \qquad 
\end{align*}

\begin{definition}[Entropy condition]
	A weak solution is an entropy solution if  $\exists c > 0 \forall t > 0 \forall x \in \mathbb{R} \forall z \in \mathbb{R}_+:$
	\begin{align*}
		u(x + z, t) - u(x, t) \le C \frac{z}{t} \texttt{ (one-sided jump)}
	\end{align*}
\end{definition}
\begin{remark}
entropy condition relaxes the obeservation that $u(, t)$ is non-decresing in a vague analogy to the physical entropy. 	
\end{remark}

\begin{example}
Given $u_l, u_r \in \mathbb{R}, u_l \neq u_r$ and let
\begin{align*}
	g(x)&:= u_l \quad x < 0\\
	&= u_r \quad x > 0
\end{align*}	
1.\textbf{case: shock wave}: $u_l > u_r$: $\sigma = \frac{F(u_l) - F(u_r)}{u_l - u_e}$
\begin{align*}
	u(x, t) &= u_l \qquad x < \sigma t\\
	&= u_r \qquad x > \sigma t
\end{align*}
\textbf{2. case: rarefaction wave $u_l < u_r$} , $\sigma = (F')^{-1}$
\begin{align*}
	u(x, t) &= u_l \qquad x < \sigma_l(t) \\
	&= G(\frac{x}{t}) \quad else\\
	&= u_r \qquad x > \sigma_r(t)
\end{align*}
\end{example}

\begin{theorem}
	Assuming $F$ is uniformly convex and smooth, then there exists at most 1 entropy solution. 
\end{theorem}
\begin{proof}
	\texttt{cf. Evans p.151}
\end{proof}

\begin{theorem}[Lax-Oleinik formula]
For $F: \mathbb{R} \rightarrow \mathbb{R}$ smooth, uniformly convex, $ g \in L^\infty(\mathbb{R})$: 
	\begin{itemize}
		\item $\forall t > 0 \exists N(t) \subset \mathbb{R}$ countable $\forall x \in \mathbb{R} \slash N(t) \exists !y(x, t) \in \mathbb{R}$ :
			\begin{align*}
				\min_{y \in \mathbb{R}} \bigg\{ t L(\frac{x - y}{t}) + h(y) \bigg\} = t L\bigg(\frac{x - y(x, t)}{t}\bigg) + h(y(x, t))
			\end{align*}
			where
			\begin{align*}
				h(x) = \int_0^x g(y) dy
			\end{align*}
		\item $\forall t > t$,  $y(, t)$ is nondecreasing $\forall x \in \mathbb{R} \slash N(t)$ define
			\begin{align*}
				u(x, t) &:= (F')^{-1}(\frac{x - y(x, t)}{t})\\
				&= (F^*)'(\frac{x - y(x, t)}{t})
			\end{align*}
	\end{itemize}
	The solution determined by the formula above is the unique entropy solution. 
\end{theorem}
\begin{remark}
We have now an alternative method to solve the scalar conservation law. First we solve the HJE and use the convex conjugate to solve the stated IVP. Note that $F$ uniformly convex is crucial here. 	
\end{remark}



\subsection{Advection-reaction equation \& Graph space}
\subsubsection{advection-reaction equation}

the \textbf{steady}(time-independent) advection reaction equation is the linear, scalar steady first order equation:
\begin{align*}
	\beta \cdot \nabla u + \gamma  u = f &\texttt{ in }\Omega \\
	u = 0  &\texttt{ on } \Gamma_-
\end{align*}
where $\Gamma_-:= \{x \in \partial \Omega | \beta(x) \cdot \nu(x) < 0\}$ for bounded Lipschitz domain in $\mathbb{R}^n$ with negative outflow along $\partial \Omega$. 

\subsubsection*{Assumptions}
\textbf{1.} $\beta \in Lip(\Omega, \mathbb{R}^n), \gamma \in L^\infty(\Omega)$, $f \in L^2(\Omega)$\\
\textbf{2.}$\exists \alpha > 0: \alpha \le \gamma - \frac{1}{2} div (\beta)$ a.e. in $\Omega$. \\
\textbf{3. Advection: } $dist(\Gamma_-, \Gamma_+ )> 0$ seperation of the inflow and outflow of the boundary. \\


\quad\\
\textbf{New aspect:} $\Omega$ general, not half space, non-constant coefficient, particular boundary condition. \\
\textbf{Aim:} design weak solution. \\
The induced linear operator and nonsymmetric bilinear form:
\begin{align*}
		&a(v,w) := \int_\Omega (\gamma v w + (\beta \cdot \nabla v)w dx + \int_{\partial\Omega}(\beta \cdot v)_- (v w)   ds\\
		&F(w) := \int_\Omega fw dx  \quad v, w \in V \\
	&\mathcal{V} := \{v \in H^1(u) : v|_{\Gamma_-} = 0\}\\
	&\mathcal{D}(\Omega) := \mathcal{D}(\Omega, \mathbb{R}^n) = C^\infty_c(\Omega, \mathbb{R}^n) 
\end{align*}
where we define $(x)_+:= \max(0, x)$ and $(x)_:= \frac{1}{2}(|\cdot| - \cdot)$.  $V$ is not the final space of our solution, however it's dense in it.

\begin{definition}[the graph space]
	we define the \textbf{graph space}:
	\begin{align*}
		V:= \{ u \in L^2| \beta \cdot \nabla u \in L^2(\Omega)\}
	\end{align*} 
\end{definition}
\begin{remark}
The graph space has the nature form:
\begin{align*}
	V:= \{ u \in L^2| \nabla( u \beta) \in L^2 \}
\end{align*}	
the nature form and the conservative form is equivalent since:
\begin{align*}
	\int_{\Omega} \nabla (u \beta) \varphi dx &= \int_{\Omega}  (\beta \cdot \nabla u)\varphi dx + \int_\Omega (u \cdot \nabla \beta) \varphi dx
\end{align*}
since $\beta \in W^{1, \infty}$ we know the second term is $L^2$ bounded. The 2 definitions are therefore equivalent.
\end{remark}

\begin{lemma}[Hilbert structure]
		the graph space $V$ equipped with the scalar product:
	\begin{align*}
		(u, v):= (u, v)_2 + (\beta \cdot \nabla u, \beta \cdot \nabla v )_2
	\end{align*}
	is a Hilbert space
\end{lemma}
\begin{proof}\quad\\
\textbf{step 1)}: $(u, v)$ is a scalar product
\\
\textbf{step 2):} $V$ is complete. to show given a Cauchy sequence $v_n \in G \subset L^2(\Omega)$, we have $v_n \overset{\|\cdot \|_v}{\rightarrow} v \in G$. Indeed, by the fact that $b \nabla v_n, v_n$ are Cauchy sequences in $L^2(\Omega)$ we have the weak form: $\forall \varphi \in \mathcal{D}(\Omega)$:
\begin{align*}
	\int_\Omega \lim_n(\beta \cdot \nabla v_n) \varphi dx 
	&= \lim_n  \int_\Omega (\beta \cdot \nabla  v_n ) \varphi dx \qquad (\beta \cdot v_n \in L^\infty(\Omega))  \\
	&= -\lim_n  \int_\Omega v_n \nabla(\beta \varphi) dx \\
	&= -\int_\Omega v \nabla \cdot   (\beta \varphi) dx \\
	&= \int_\Omega w \varphi dx 
\end{align*}
this implies $\lim v_n = v \in V$ with $\beta \cdot \nabla v = w$. 
\end{proof}

\begin{theorem}[Properties of the graph space]\quad\\
	\textbf{a).} $C^\infty(\bar{\Omega})$ is dense in $V$. \\
	\textbf{b).} there exists a unique trace map:
	\begin{align*}
		\gamma: L(V, L^2(|b \cdot 
		\nu|, \partial \Omega))
	\end{align*}
	such that :
	\begin{align*}
		f|_{\partial\Omega} = \gamma(f) \quad \forall f \in \bar{\Omega}) \cap V
	\end{align*}
	Here $L^2(|b \cdot \nu|, \partial \Omega)$ is the weighted $L^2$ space \\
	\textbf{c):} the integration by parts formula holds
	\begin{align*}
		\int_\Omega (b \nabla v ) w + (b \cdot \nabla w) v + vw div( b) dx 
		&= \int_{\partial \Omega} | b \cdot v | \gamma (v) \gamma(w) ds
	\end{align*}	
\end{theorem}
\begin{proof}
\textbf{a):}\\
For $v \in V$ it satisfies $\beta \cdot \nabla v \in L^2(\Omega)$, and $\eta_\epsilon$ the standard mollifier, define for $G:= \beta v$:
\begin{align*}
	&v_\epsilon := v * \eta_\epsilon \in \mathcal{D}(\mathbb{R}) \\
	&G_\epsilon := G * \eta_\epsilon \in \mathcal{D}(\mathbb{R})
\end{align*}
\textbf{step 1). the boundedness}
\begin{align*}
	\|\beta \cdot \nabla v_\epsilon \|_2 \le \|\beta \cdot \nabla v_\epsilon - div(G_\epsilon)\|_2 + \|div(G_\epsilon)\|_2
\end{align*}
\textbf{step 2). weak convergence:}\\ By the reflexivity of $V$ we get a weak convergent subsequence 
\begin{align*}
	\beta \cdot \nabla v_\epsilon \overset{V}{\rightharpoonup} \beta \cdot \nabla v
\end{align*}
\textbf{step 3). density}\\ By Mazur lemma we find $v_j \in conv\{v_\epsilon\}$ such that 
\begin{align*}
	\beta \cdot \nabla v_j \overset{L^2}{\longrightarrow} \beta \cdot \nabla v
\end{align*}
we have:
\begin{align*}
	\|v_j - v\|_V 
	&= \int_\Omega |v - v_j|^2 dx + \int_\Omega |\beta \cdot \nabla v_j - \beta \cdot \nabla v|^2 dx \\
	&\rightarrow 0
\end{align*}
the first converges to 0 by mollification and the second by \textbf{ii).}\\
\quad\\
\textbf{b).} For $\varphi \in \mathcal{V}$:
\begin{align*}
	\gamma(\varphi^2) 
	&= \int_{\partial\Omega} |\beta \cdot \nu| \varphi^2 dS \\
	&= \int_{\Gamma_+} \chi_+ \varphi^2 (\beta \cdot \nu) dS - \int_{\Gamma_-}\chi_- \varphi^2 (\beta \cdot \nu)dS\\
	&= \int_{\partial\Omega}(\chi_+ - \chi_-) (\varphi^2 \beta) \cdot \nu dS\\
	&= \int_\Omega div\bigg((\chi_+ - \chi_-)\varphi^2 \beta\bigg) dx\\
	&= \int_\Omega \langle \varphi^2, \nabla \cdot [(\chi_+ - \chi_-)\beta]\rangle + \langle(\chi_+ - \chi_-)\beta, 2\varphi\nabla \varphi\rangle dx \\
	\|\gamma\|_{L^2(|b \cdot \nu, \partial\Omega)}^2
	&\le \|\nabla \cdot \beta\|_{L^\infty}\|\varphi\|^2_{L^2(\Omega)} +
	\underbrace{ \|\varphi\|_{L^2(\Omega)} \|\beta \cdot \nabla \varphi\|_{L^2(\Omega)}}_{ab \le a^2 + b^2}\\
	&\le C \|\varphi\|_V^2
\end{align*}
there exists a bounded trace map from a dense subset of $V$ to $L^2(|b\cdot \nu|,\partial\Omega)$. Hence there exists an unique completion of the bounded linear map.\\
\quad\\
\textbf{c).} trivial	
\end{proof}



\subsubsection*{the well-posedness}
\begin{lemma}[The well posedness in $V$]
Observe the bilinear form:
\begin{align*}
		a(v, w) &= \int_\Omega \bigg(\gamma v w + (b \cdot \nabla v) w \bigg) dx 
		+ \int_{\partial\Omega} (b \cdot  \nu) \gamma(v) \gamma(w) ds
\end{align*}
\textbf{1).}$a(\cdot, \cdot)$ is bounded\\
\textbf{2).}$a(\cdot, \cdot)$ is $L^2$ coercive, in the sense that
\begin{align*}
	a(v, v) \geq \mu_0 \|v\|^2_{L^2(\Omega)} + \int_{\partial \Omega }|\beta \cdot \nu | v^2 ds
\end{align*}
\textbf{3).}
define $V_0 := \{ v \in V:  v_{|\partial \Omega^- } = 0 \}$	
$a_0(\cdot, \cdot)$ is the corresponding bilinear form in $V_0$. Then $a_0$ satisfies the $\inf$-$\sup$ condition in the sense that:
\begin{align*}
	\inf_{v \in V_0, \|v\|_V = 1} \sup_{w \in L^2(\Omega), \|w\|_{L^2(\Omega)} = 1} a_0(v, w) >0
\end{align*}
\textbf{4.}
$a_0(v, w)$ is non-degeneracy in the second term ( in $w$).
\end{lemma}
\begin{proof}
\quad\\
\textbf{1): the weak formulation is well-defined bounded bilinear form $\mathcal{V} \times \mathcal{V} \rightarrow \mathbb{R}$. }
\begin{align*}
	a(u, \varphi) &= \int_\Omega (\beta \cdot  \nabla u) \varphi + \gamma u \varphi dx + \int_{\partial\Omega} (\beta \cdot \nu)_- u\varphi dx \\
	&\le \|\beta \nabla u\|_2\|\varphi\|_2 + \|\gamma\|_\infty \|u\|_2\|\varphi\|_2 + \|\beta \cdot \nu \|_\infty \|u\|_V \|\varphi\|_V
\end{align*}
since $\|\beta \cdot \nabla u\|_2, \|u\|_2$ are bounded by $\|u\|_V$, we have the desired result. \\
\textbf{2). coercivity}
\begin{align*}
	a(v, v) &= \int_\Omega (\beta \cdot \nabla v) v + \gamma v^2 dx + \int_{\partial\Omega} |b \cdot \nu| \gamma(v) \gamma(w)dS \\
	(\texttt{last theorem}) \quad
	&= \int_\Omega \underbrace{ \bigg(-\frac{1}{2} div(\beta) +  \gamma\bigg)}_{> 0 } v^2 dx + \frac{1}{2}\int_{\partial\Omega} |b \cdot \nu| \gamma(v) \gamma(w)dS \\
\end{align*}
\quad\\
\textbf{3). inf-sup condition}:\\
\textbf{4).Non-degeneracy}


\end{proof}
\quad\\

Next we are going to prove the well-posedness of the advection equation. Recall for $\varphi \in \mathbbm{1}_{D}(\mathbb{R}), b \in Lip(\Omega, \mathbb{R}^n)$, we have
\begin{align*}
	\int_\Omega div(b \varphi ) dx = \int_\Omega \varphi div(b) dx + \int_\Omega b \cdot \nabla \varphi dx 
\end{align*}




For the abstract linear problem: For a bounded bilinear form $a \in L(X \times Y, \mathbb{R}), f \in Y'$
\begin{align*}
	&\texttt{Find } u \in X\\
	&s.t.\quad a(u, w) = \langle f, w \rangle_{Y', Y} \quad \forall w \in Y 
\end{align*}
The problem is called \textbf{well-posed} if there exists an unique $u$ to the problem. 
The tool we are going to use is the so-called \textbf{Banach-Necas-Babuska theorem}

\begin{theorem}[Banach-Necas-Babuska]
	Let $X$ be a Banach space and $Y$ a reflexive Banach space. The problem is well-posed if and only if:
	\begin{enumerate}
		\item the \texttt{inf-sup} condition holds, i.e.
		\begin{align*}
			\inf_{\|x\| = 1} \sup_{\|y\|= 1} a(x, y) \geq C_{sta}
		\end{align*}
		\item \textbf{non-degeneracy:}  $\forall y \in Y$:
		\begin{align*}
			\bigg(\forall x \in X:a(x, y) = 0 \bigg)\Rightarrow (y = 0)
		\end{align*}
	\end{enumerate}
\end{theorem}

\begin{theorem}
	The problem is well-posed in $V$, i.e. 
	\begin{align*}
		\forall f \in L^2(\Omega) \exists! u \forall v \in V: a(u, v) = \int_\Omega fv dx
	\end{align*}
\end{theorem}
\begin{remark}
\quad
\begin{itemize}
	\item $\gamma$ is surjective
	\item If the weak form is satisfied then we have $\gamma(u) = g$ a.e. on $\Gamma$. 
\end{itemize}
\end{remark}

\begin{proof}
\quad\\
\textbf{2):the weak form solves the PDE}
by the assumption we have
\begin{align*}
	\int_\Omega (\beta \nabla u + \gamma u - f) \varphi dx  = 0 \qquad \forall \varphi \in C_c^\infty(\Omega)
\end{align*}
by the fundamental lemma of calculus we know $u \in V$ solves the PDE. Since $u \in \mathcal{V}$ the boundary condition holds trivially. \\
\textbf{3).}
assuming $u_1, u_2$ are 2 weak solutions, define
\begin{align*}
	e:= u_1 - u_2 \in \mathcal{V}	
\end{align*}
By coercivity we have:
\begin{align*}
	a(e, e) = a(u_1, 2) - a(u_2, e) = f(e) - f(e) = 0
\end{align*}
by $L^2$ coercivity we know $\|e\|_2$ = 0. Hence $e$ is 0 a.e. 
\end{proof}



\section{Excercise}
\subsubsection*{Integration by parts}
Introduction to Sobolev space:
Token: Hamilton-Jacobi-2022
\begin{definition}
	$\Omega \subset \mathbb{R}^N$ open, for $u \in C^1(\Omega)$ and $\varphi \in C_C^\infty(\Omega) : = \{ \varphi \in C^\infty(\Omega): supp(\varphi) \subset\subset \Omega\}$ ($\subset\subset$ compactly contained(there exists a compact set $K$ such that $A \subset K \subset \Omega$)
\end{definition}
Integration by parts: $\forall \varphi \in C_c^\infty$:
\begin{align*}
	\int_{\Omega} u \partial_j \varphi dx = - \int_{\Omega} \partial_j u \varphi dx + \int_{\partial \Omega} u \varphi \nu_j dx \quad j = 1, \dots, n
\end{align*}
A less smooth function $g_j \partial_j \in L_{loc}^1(\Omega)$ is called weak derivative of $u \in L^1_{loc}$ of 
\begin{align*}
	\int_{\Omega} u \partial_j \varphi u = \int_{\Omega} g_j \varphi dx\qquad \varphi \in C_c^\infty(\Omega)
\end{align*}

\begin{lemma}
	Weak derivatives are unique a.e. in $\Omega$. 
\end{lemma}
\begin{proof}
	Let $g_j, \tilde{g_j}$ satisfying (WD), we have
	\begin{align*}
		0 = \int (g_j - \tilde{g_j}) \varphi dx \quad \forall \varphi \in C_c^\infty(\Omega)
	\end{align*}
	The fundamental lemma of calculus of variations concludes the proof. 
\end{proof}

The formula (IBP) shows that (WD) holds for $C^1$ functions. i.e. classical derivatives are weak derivatives. 


\begin{example}[1-D example]
procedure for piecewise smooth functions: $u: \mathbb{R} \rightarrow \mathbb{R}, u(x) = |x|$. For $\varphi \in C_c^\infty(\mathbb{R})$ with $supp(\varphi) \subset [-M, M]$. We have
\begin{align*}
	\int_\mathbb{R} u \varphi' dx = \int_{-M}^0 -x \varphi'dx + \int_0^M x \varphi'dx \\
	= \int_{-M}^0 -1 \varphi dx + 0 \varphi(0) - M \varphi(-M) - \int_0^M 1 \varphi dx + M \varphi(M) - 0\varphi(0)
\end{align*}	
this vanishes iff it's continuous, i.e. $\varphi(0) = \varphi(0)$ from both sides.
\end{example}

\begin{example}[Heaviside functions]
$H: \mathbb{R} \rightarrow \mathbb{R}, H = \chi_{[0, \infty)}$. The same computation yields:
\begin{align*}
	\int_\mathbb{R} H \varphi'dx = \int_0^M \varphi' dx = \varphi(M) - \varphi(0) = - \varphi(0)
\end{align*}	
The map $\delta_0: \varphi \rightarrow \varphi(0)$ belongs to the dual space of $C_c^\infty(\mathbb{R})$ such bounded linear functions are called distributions. But there exists no functions $g \in L_{loc}^1(\mathbb{R})$ such that $\varphi(0) = g \varphi dx$ for alll $\varphi$. This means the generalised distributional derivative $\delta_0$ is no regular distribution. In particular, $H$ is not weakly differentiable. 
\end{example}

\subsubsection*{Lebesgue and Sobolev spaces}
\begin{definition}[Sobolev space]
	The Sobolev space is defined as 
	\begin{align*}
		W^{k, p}(\Omega) := \{ v \in L^p(\Omega), \forall \alpha \in \mathbb{N}_0^\alpha, |\alpha| \le k, D^\alpha v \textbf{ exists in weak sence}\}
	\end{align*}
	For $p = 2$, we have a Hilbert space with the scalar product:
	\begin{align*}
		\langle u, v \rangle = \sum_{\alpha} \langle D^\alpha u, D^\alpha v \rangle
	\end{align*}
	equivalent definition for $n = 1$: 
	\begin{align*}
		W^{1, p}(0, 1) = \{ f \in C^0[0, 1]: \exists g \in L^p(0, 1) \exists c \in \mathbb{R} \forall x \in [0, 1]: f(x) = c + \int_0^x g(y)dy \}
	\end{align*}
	this would relate to Bochner spaces later. 
\end{definition}

\subsubsection*{Mollification}
1. \textbf{Standard mollifier}: Define $\varphi \in C_c^\infty(\mathbb{R}^n)$ by 
\begin{align*}
	\varphi(x) := C\exp(\frac{1}{|x|^2 - 1}) \quad |x| < 1
\end{align*}
and 0 everywhere else. Choose $C>0$ such that $\int_\mathbb{R} \varphi dx = 1$. For $\epsilon >0$, we attain $\varphi_\epsilon(x) = \frac{1}{\epsilon^n} \varphi (\frac{x}{\epsilon})$. As we see the support is shrinking and \\
\textbf{2. Convolution:} Given $f \in L^p(\mathbb{R}^n)$, define 
\begin{align*}
	f_\epsilon := \varphi_\epsilon * f(x) = \int_{B_0(\epsilon)} \varphi_\epsilon (x - y) f(y) dy = \int \varphi_\epsilon(z)f(x - z)dz
\end{align*} 




\chapter{Non-linear PDE}
\section{Selected Examples}


\section{The Fixed Point Theorem}
\begin{theorem}
Given $X$ a Banach space and $M \subset X$ \textbf{closed} and \textbf{convex}. If $T : M \rightarrow M$ is continuous, then $T$ admits at least one fixed point. 	
\end{theorem}


\subsection{Calculus of Variations}
\begin{definition}[Null Lagrangian]
	The Lagrangian is called \textbf{null-Lagrangian} if the corresponding \textbf{Euler-Lagrange function} is satisfied by all $C^2$-functions.
\end{definition}


\subsection{The Brouwer fixed point theorem}
The first tool in the toolbox of fixed point theorem is the Brouwer fixed point theorem. It states that in the finite dimensional space, continuous map on a closed ball admits at least 1 fixed point. Although we aren't able to apply it directly in most of cases since the functional space is usually infinite dimensional, it provides a off-hand tool in constructing finite dimensional existence. 

\begin{theorem}[Brouwer fixed point theorem(finite dimensional)]
	Given a continuous function defined on a closed ball in the finite dimensional real space, the function admits \textbf{at least} one fixed point. \par 
	Given $T: \bar{B_R(0)} \rightarrow \bar{B_R(0)}$ continuous with $B_R(0) \subset \mathbb{R}^m$, $T$ has at least one fixed point. 
\end{theorem}
\begin{remark}
The result is wrong in the infinite dimensional case. 	
\end{remark}

\begin{lemma}[no-retraction principle]
	There exists no continuous map $ w: B \rightarrow \partial B$, such that $w(x) = x$ for all $x \in \partial B$. 
\end{lemma}



\subsection{The Schrauder Fixed Point Theorem}
When we prove existence with fixed point theorem in the infinite dimensional case, Brouwer is no longer applicable. 

\subsubsection{Nonlinear compact operator }
\begin{proposition}[Approximation of compact operator by finite dimensional operator]
Given $M$ bounded $T: M \rightarrow Y$, then the followings are equivalent
\begin{enumerate}
	\item T is compact
	\item $\forall n \in \mathbb{N} \exists T_n : M \rightarrow Y: T_n(M) \subset Y_n$ with $\dim (Y_n) < \infty$, $T_n \overset{uniform}{\rightarrow} T$. 
\end{enumerate}
\end{proposition}
\begin{proof}\quad\\
\textbf{2)$\Rightarrow$ 1)}: \\
\textbf{i).} $T$ is compact. By uniform convergence of the operator we have
\begin{align*}
	\forall \epsilon > 0: \sup_{x \in M} \|Tx - T_{n_0}x \| \le \frac{\epsilon}{3} \quad \texttt{for $n_0$ sufficiently large}
\end{align*}
By the compactness of $T_{n_0}$ we have
\begin{align*}
	\forall \epsilon  \exists N \in \mathbb{N}: T_{n_0}(M) \subset \bigcup_{i = 1}^N B_{\frac{\epsilon}{3}} (x_i)
\end{align*}
therefore $\forall y \in M$
\begin{align*}
	\|T( y) - T(x_i)\|_Y \le \|T(y) - T_{n_0} (y)\|_Y + \|T_{n_0}(y) - T_{n_0}(x_i)\|_Y + \|T_{n_0}(x_i) - T(x_i)\|_Y
\end{align*}
on the RHS the first and the third are $\le \frac{\epsilon}{3}$ by uniform convergence and the second by the compactness of $T_{n_0}$. \\
\textbf{2)$ \Rightarrow$ 1)} 
\end{proof}



\begin{theorem}[Schrauder Version I]
	Given $X$ a Banach space, $M \subset X$ convex and compact, $T: M \rightarrow M$ continuous. Then $T$ admits a fixed point. 
\end{theorem}

\begin{lemma}[Mazur's lemma]
	Convex hull of precompact set is compact. 
\end{lemma}

\begin{theorem}[Schrauder Version II]
	Given $X$ a Banach space, $M \subset X$ convex, closed and bounded. The map $T: M \rightarrow M$ is compact. Then $T$ admits a fixed point. 
\end{theorem}

\subsection{An application to ODEs}
\begin{theorem}[Peano]
Given $(t_0, y_0) \in \mathbb{R} \times \mathbb{R}^n$, $Q:= \{(t, y): |t - t_0| \le a, |y - y_0| \le b\}$, let $f: Q \rightarrow \mathbb{R}^n$ be continuous and uniformly bounded, i.e., $|f(t, y)| \le K$ for some $K \in \mathbb{R}$. Then the integral equation
\begin{align*}
	y(t):= y_0 + \int_{t_0}^t f(s, y(s)) ds 
\end{align*}
admit a local solution.	
\end{theorem}

\begin{corollary}[Peano existence theorem in ODE]
	
\end{corollary}


\subsection{Application of Schrauder}



\section{Nonlinear Poisson Problems}
\begin{lemma}[Estimate of the inverse Dirichlet Laplacian]
Consider the bounded linear map on $\Omega$ with $\partial\Omega \in C^{1, 1}$
\begin{align*}
	\Delta_D: H^1_0(\Omega) \cap H^2(\Omega) &\rightarrow L^2\\
	u &\mapsto -\Delta u
\end{align*}
The operator is an isomorphism therefore we may define the inverse
\begin{align*}
	A_D: L^2 &\rightarrow H^1_0 \cap H^2\\
	g &\mapsto (-\Delta_D)^{-1}(g)
\end{align*}
and we have the estimate
\begin{align*}
	\|A_D\|_{H^2} \le \|g\|_{L^2}\\
	\|A_D\|_{H^1_0} \le \|g\|_{L^2}
\end{align*}
\end{lemma}


\subsection{Solvability with Schauder}
The PDE under consideration is the \textbf{nonlinear Poisson equation}.       
\begin{align*} 
\textbf{(NP)}\qquad
		- \Delta u = f(u) &\texttt{ on } \Omega\\
		u  = 0 &\texttt{ on } \partial \Omega 
\end{align*}           
\begin{theorem}[Well-posedness of the nonlinear Poisson Problem]
	Observe the non-linear elliptic system:
	\begin{align*}
		- \Delta u = f(u) &\texttt{ on } \Omega\\
		u  = 0 &\texttt{ on } \partial \Omega 
	\end{align*}
	If $f \in L^\infty(\Omega) \cap C^0(\Omega)$, then the system admits a weak solution $u \in H^1(\Omega)$. Moreover, if $\partial\Omega \in C^{1, 1}$, $u \in H^2(\Omega)$.
\end{theorem}
\begin{proof}
Define the map:
\begin{align*}
	T: L^2(\Omega) \rightarrow L^2(\Omega), u \mapsto A_D(f(u))
\end{align*}
we would like to apply the Schrauder to see that it admits a fixed point. 
	\quad\\
	\textbf{Step1:  $T$ is well-defined}: By the continuity of $f$ we know
	\begin{align*}
		\forall u \in L^2: f(u) \in L^\infty \in L^2
	\end{align*}
	Moreover we know the inverse Laplace operator satisfies
	\begin{align*}
		 \|A_D(f(u))\|_{L^2} \le \|A_D(f(u))\|_{H^1_0} \le \|f(u)\|_{L^2} < \infty
	\end{align*}
	since $f \in L^\infty$. Therefore $T$ is well-defined.\\
	\textbf{Step2: $\exists R > 0: T(B_R(0)) \subset \overline{B_R(0)}$, moreover, $B_R(0)$ is precaompact in $H^1_0(\Omega)$. } Observe $\forall u \in L^2(\Omega):$
	\begin{align*}
		\|T(f(u))\|_{L^2} \le \|A_D(\|f\|_\infty)\|_{L^2} 
		\le \|f\|_\infty \|A_D(1)\|_{L^2} \le \|f\|_\infty C_p |\Omega|^{\frac{1}{2}}: = R
	\end{align*}
	Therefore we have 
	\begin{align*}
		T(B_R(0)) \subset B_R(0)
	\end{align*}
	By Rellich we know $B_R(0)$ relative compact in $H^1_0(\Omega)$. 
	therefore we know $T(B_R(0))$ is precompact.
	\\
	\textbf{Step3: $T$ is continuous.} Given $u_n \rightarrow u$ in $L^2$, we know up to a subsequence 
	\begin{align*}
		u_k  \overset{a.s.}{\rightarrow} u
	\end{align*}
	observe
	\begin{align*}
		\|T(u_k) - T(u)\|_{L^2} 
		\overset{\texttt{linearity of $A_D$}}{=} \|A_D(f(u) - f(u_k))\|_{L^2} 
		\overset{lemma}{\le} \|f(u_k) - f(u)\|_{L^2} 
		\overset{\texttt{DCT and $f \in C^0$}}{\rightarrow} 0
	\end{align*}
	\\
	In conclusion, $T$ is compact on a convex and compact ball $B_R(0)$ in $H^1_0$, therefore $T$ admits a fixed point and the elliptic system admits a solution.
\end{proof}


\subsection{A priori estimate with weak maximum principle}
\begin{definition}
	for $u \in W^{1, p}(\Omega)$, we say $u \ge 0$ on $\partial\Omega$ if $u_ \in W^{1, p}_0(\Omega)$.
\end{definition}

\begin{proposition}
	Let $u \in H^1(\Omega)$ and assume
	\begin{enumerate}
		\item $u \ge 0$ on $\partial\Omega$
		\item $-\Delta u \ge 0$ in the weak sense, i.e., 
		\begin{align*}
			\int_\Omega \nabla u \cdot \nabla v dx \ge 0 \qquad \forall v \in H^1_0(\Omega)
		\end{align*}
		Then we have $u \ge 0$ a.e.
	\end{enumerate}
\end{proposition}

\begin{example}[Exponential diffusion]
Observe the system
\begin{align*}
	-\Delta u = e^{-u} &\texttt{ in } \Omega\\
	u = 0 &\texttt{ on } \partial \Omega
\end{align*}
Notice we have on $\partial\Omega$ $u \ge 0$ and the weak form 
\begin{align*}
	 \int_\Omega \nabla u \nabla \varphi = - \int_\Omega \Delta u \varphi  \ge 0 \quad \forall \varphi \in H^1_0(\Omega) \ge 0 
\end{align*}
By the weak maximum principle we have
\begin{align*}
	u \ge 0 \texttt{ a.e. in  } \Omega 
\end{align*}
we may truncate the diffusion term $f$ by
\begin{align*}
	\tilde{f} &= f \quad u \ge 0\\
	&= 1 \quad u \le 0 
\end{align*}
we have $f \in L^\infty \cap C^0(\mathbb{R})$. By the theorem in the last chapter we get a solution
\begin{align*}
	u \in H^1_0(\Omega)
\end{align*}
Moreover, the solution is unique since $f$ is non-increasing. Assuming $u_1, u_2 \in H^1_0(\Omega)$ are 2 solution of the diffusion equation, by linearity $u_1 - u_2$ is also a solution, we have
\begin{align*}
	\int_\Omega |\nabla(u_1 - u_2)|^2 dx = \int_\Omega (f(u_1) - f(u_2))(u_1 - u_2) \le 0 
\end{align*}
\end{example}





\section{Variational Inequalities \& Monotone Operators}
Differentiation in Banach space. 
\subsection{Variational Inequality}
\textbf{Motivation}:
\begin{align*}
	M:= \{v \in H^1_0(0, 1), v(0) = a, v(1) = b \}
\end{align*}
define the functional 
\begin{align*}
	I(v) = \int_I \sqrt{1 + v'} dx 
\end{align*}
want to find
\begin{align*}
	\min_{ v \in M} I(v) 
\end{align*}
Let $\varphi \in H^1_0(0, 1)$, we know $u$ is a minimiser of $I(v)$ if 
\begin{align*}
	\frac{d}{d\epsilon} I(u + \epsilon \varphi)|_{ \epsilon = 0}   = \int_I \bigg(\sqrt{1 + |u'|^2} \bigg)^{-1} u' \varphi' dx := \langle F(u), \varphi \rangle  
\end{align*}

\begin{remark}
If $u \in H^2(0, 1)$, then we have
\begin{align*}
	-\int_I (\frac{u'}{\sqrt{1 + |u'|^2}})' \varphi = 0 \quad \varphi \varphi \in H_0^1(\Omega)
\end{align*}
therefore we know $u'$ is constant and the minimal curve is a straight line. 	
\end{remark}
The \textbf{Ellipticity} is clear in this case.  

\quad\\
\textbf{The obstacle problem:} Let $h \in C^1(I, \mathbb{R})$ we define
\begin{align*}
	M:= \{v \in H^1_0(0, 1), v(0) = a, v(1) = b, v \ge h \}
\end{align*}
The Gateaux derivative cannot be expected anymore, on the other hand we observe the variational inequality:
\begin{align*}
	\langle F(u), u - v \rangle  \le 0 \qquad \forall v \in M 
\end{align*}
on a closed and convex set $M$. 



\quad\\
\textbf{General setting}
\begin{enumerate}
	\item $X$ is reflexive, separable Banach space with the dual space $X^*$. 
	\item $M \subset X$ closed, non-empty, convex. 
\end{enumerate}

consider the variational problem:
\begin{align*}
	I(v) = \int_a^b L( \nabla u, u, x) dt 
\end{align*}

The most simple example would be the shortest curve between 2 points
\begin{align*}
	I(v) := \int_a^b \sqrt{1 + u'(x)} dx  
\end{align*}
and consider the minimisation problem:
\begin{align*}
	\min_{v \in M} I(v) \qquad M := \{ v \in H^1([0, 1]) | v(0) = a, v(1) = b \}
\end{align*}
$M$ is well-defined by the classical embedding result.

\quad\\
\textbf{Assumptions on $F$:}\\
\textbf{F1:} $F$ maps bounded set of in $X$ to bounded set in $X^*$.\\
\textbf{F2:} Coercivity: There exists $u_0 \in M: \langle F(u), u - u_0 \rangle / \|u - u_0\| \rightarrow \infty$ as $\|u \| \rightarrow \infty$. \\
\textbf{F3:} Given $u_n \rightharpoonup u, F(u_n) \rightharpoonup u^*$, we have
\begin{align*}
	\langle u^*, u \rangle \le \liminf_{n \rightarrow } \langle F(u_n), u_n \rangle 
\end{align*}
Furthermore, if $\langle u^*, u \rangle = \liminf \langle F(u_n), u_n \rangle$, it holds
\begin{align*}
	\langle F(u) - u^*, u - v \rangle \le 0  \quad \forall v \in M
\end{align*}
\begin{remark}
\quad\\
\textbf{1).} The coercivity condition endowed $u$ with boundary growth condition.\\
\textbf{2).} \textbf{F3}) gives an upper bound for the variational inequality. \\
\textbf{3).} For \textbf{F2} it is essential to have $u_0 \in M$.
\end{remark}

\begin{example}[Counterexample for the coercivity condition]
\end{example}



\begin{theorem}
	If the operator $F$ satisfies \textbf{(F1) - (F3)}, then $F$ satisfies the variational inequality
	\begin{align*}
		\langle F(u), u - v \rangle \le 0
	\end{align*} 
\end{theorem}
\begin{proof}
	By Galerkin method. 
\end{proof}


\subsection{the Monotone Operators}
\begin{definition}[The monotone operator]
	An operator $A: M \rightarrow X^*$ is a \textbf{monotone operator} if 
	\begin{enumerate}
		\item $\langle A(x) - A(y), x - y \rangle \ge 0 \quad \forall x, y \in M$
		\item \textbf{hemi-continuous:} $\forall w \in X \forall v, u \in M$:
		\begin{align*}
			t \mapsto \langle A(u + t(v - u)), w \rangle 
		\end{align*}
		is continuous on $[0, 1]$.
		\item 
	Moreover, the operator is \textbf{strictly monotone} if 
	\begin{align*}
		\langle A(u) - A(v), u - v \rangle > 0 \qquad \forall u, v \in M
	\end{align*}
	\textbf{strongly monotone} if 
	\begin{align*}
		\exists c_0 \forall u, v \in M : \langle A(u) - A(v), u - v \rangle \ge c_0 \|u - v \|^2 
	\end{align*}
	\end{enumerate}
\end{definition}
hemi-continuity can be seen as one-sided continuity. It ensure the cone property of the monotone operator. As we have seen in the case of the operator $F$, by taking the convex combination of $u$ and $v$ the non-positivity of the operator can be extended from the cone corner $u$ to the direction of $v$. The hemi-continuity has the same function to preserve the cone property w.r.t. the nonlinearity $A$. 

\begin{lemma}[Minty]
	Given a monotone operator $F$, for $u \in M$ and $u^* \in X^*$, the following conditions are equivalent
	\begin{enumerate}
		\item $ \langle F(u) - u^*, u - v \rangle \le 0 \quad \forall v \in M$
		\item $\langle F(v) - u^*, u - v \rangle \le 0 \quad \forall v \in  M$.  
	\end{enumerate}
\end{lemma}
\begin{proof}
	\quad\\
	$\Rightarrow$: 
	\begin{align*}
		\langle F(v) - u^*, u - v \rangle = 
		\langle F(u) - u^*, u - v \rangle - \langle F(u) - F(v), u - v \rangle \le 0
	\end{align*}
	$\Leftarrow$: Define $v_\epsilon:= u + \epsilon( v - u)$ for $\epsilon \in [0, 1]$,  we know by assumption
	\begin{align*}
	 	0 \ge \langle F(v_\epsilon) - u^*, u - v_\epsilon \rangle = \langle F(v_\epsilon) - u^*, \epsilon (u - v) \rangle 
	 	\Rightarrow \langle F(v_\epsilon) -  u^*, u - v\rangle \le 0
	\end{align*}
	Letting $\epsilon \rightarrow 0$ we have by the hemi-continuity:
	\begin{align*}
		\langle F(u) - u^*, u - v \rangle \le 0
	\end{align*}
\end{proof}

\begin{lemma}
	Given a monotone operator $A$ we have
	\begin{enumerate}
		\item $A$ satisfies the lower semi-continuous condition \textbf{(F3)}.  
		\item If the operator is strictly monotone, then the variational inequality
				\begin{align*}
			\langle F(u), u - v \rangle \le 0 \quad \forall v \in M
		\end{align*}
		exists at most one solution. 
		\item If the operator is strongly monotone, then the coercivity condition holds w.r.t. all $u_0 \in C$
	\end{enumerate}
\end{lemma}
\begin{proof}
\quad\\
\textbf{1)}: Observe
\begin{align*}
	&0 \le \langle F(u) - F(u_n), u - u_n \rangle = \langle F(u), u - u_n \rangle + \langle F(u_n), u_n \rangle - \langle F(u_n), u \rangle \\
	&\Rightarrow \langle F(u_n) , u \rangle \le \langle F(u_n), u_n \rangle \\
	&\Rightarrow \langle u^*, u \rangle \le \liminf_{n \rightarrow \infty} \langle F(u_n), u_n \rangle \rightharpoonup \langle u^*, u_n \rangle \rightharpoonup \langle u^*, u \rangle 
\end{align*}
therefore
\begin{align*}
	\langle u^*, u \rangle = \liminf \langle F(u_n), u_n \rangle 
\end{align*}
Furthermore, by monotonicity
\begin{align*}
	&0 \le \langle F(u_n) - F(v), u_n - v \rangle = \langle F(u_n), u_n \rangle - \langle F(u_n), v \rangle  - \langle F(v), u_n - v \rangle\\
	&\Rightarrow
	\langle F(v), u_n - v \rangle - \langle F(u_n), u_n - v\rangle \le \langle F(u_n), u_n \rangle - \langle F(u_n), u_n \rangle\\
	&\Rightarrow 
	\langle F(v) - u^*,  u - v \rangle \le \liminf \langle F(u_n), u_n \rangle - \langle u^*, u \rangle  = 0
\end{align*}
By the Minty's trick we have
\begin{align*}
	\langle F(u) - u^*, u - v \rangle \le 0 \qquad \forall v \in M
\end{align*}
\textbf{2):} Suppose $u_1, u_2$ satisfies $\langle F(u), u - v \rangle \le 0 \quad \forall v \in M$, then  we have
\begin{align*}
	\langle F(u_1) - F(u_2), u_1 - u_2 \rangle \le 0
\end{align*}
which contradicts the strict monotonicity. \\
\textbf{3)}
\begin{align*}
	\langle F(u_0), u_0 - u \rangle \slash \|u_0 - u\|
	&= \langle F(u_0) - F(u), u_0 - u \rangle  \slash \|u - u_0 \| + \langle F(u), u_0 - u \rangle \slash \|u - u_0\|\\
	&\ge c_0 \|u - u_0\| + C
\end{align*}
where $c_0$ is the strong monotone constant of $F$ and $C$ depends on the operator norm of $F(u)$. We can conclude that
\begin{align*}
	\frac{\langle F(u_0), u_0 - u \rangle }{\|u_0 - u\|} \overset{\|u_0 -u\| \rightarrow \infty}{\longrightarrow} \infty
\end{align*}
\end{proof}


\subsection{Proof of the existence theorem}
\begin{proposition}
Given $X$ a Hilbert space, if $\dim X < \infty$, $F$ satisfies \textbf{(F1)} - \textbf{(F3)}, then there exists $u \in M$ with
	\begin{align*}
		\langle F(u), u - v \rangle \le 0 \quad \forall v \in M
	\end{align*} 
\end{proposition}
\begin{proof}
\quad\\
\textbf{Step 1: In a neighbourhood $B_\epsilon M$ of $M$:}\\
\textbf{Step 2: On bounded $M$:}\\
\textbf{Step 3: Unbounded $M$:}\\
\textbf{Step 1:}\\
Since $X$ a Hilbert space and $M$ a bounded convex closed set, we define
\begin{align*}
	\tilde{F}: M \rightarrow X^*,\quad u \mapsto \langle R^{-1} F(u), \cdot \rangle 
\end{align*}
where $R$ the Riesz map. By the approximation theorem in the Hilbert space we have for $w \in X$
\begin{align*}
	\langle w - P_M w,  v - P_M w \rangle \le 0 \quad \forall v \in M
\end{align*}
Note if we have $P_M w = u \in M$ and $w = u - \tilde{F}(u) $, we would be able to conclude the variational inequality. To prove that, observe the 2 equation above defines a fixed point map:
\begin{align*}
	G: M \rightarrow M, \quad G =  P_M(I - \tilde{F}) 
\end{align*} 
If $G$ admits a fixed point, then we have
\begin{align*}
	&G(u) = u\\
	&\Leftrightarrow
	u = P_M u - P_M\tilde{F}(u)\\
	& \Rightarrow P_M w = u, \quad w - P_M w = - \tilde{F}(u)
\end{align*}
Now we prove the fixed point exists by Brouwer, we need to prove $F$ therefore $G$ is continuous: By the boundedness of $M$ $u_n$ is bounded, by \textbf{(F1)} $F(u_n)$ is bounded in $X$.\\
$\Rightarrow$: $F(u_n)$ admits a convergent subsequence $u_k$ since $\dim X < \infty$, $F(u_k) \rightarrow u^*$. $u^*$ is a limit therefore also a weak lower limit. \\
$\Rightarrow$: $\langle F(u) - u^*, u - v \rangle \le 0 \quad \forall v \in B_\epsilon M$ by \textbf{(F3)}. \\
$\Rightarrow$: Choose $v = u \pm \epsilon v$ we have $F(u) = u^*$
$\Rightarrow$: $F$ is continuous hence $G$. \\
\textbf{Step 2:} Define $M_\delta := M \backslash B_\delta(\partial M)$, since $\dim X < \infty$, by the boundedness of $M$, there exists a convergent subsequence for $u_{\delta_k}$ in $X$, $F(u_{\delta_k})$ in $X^*$:
\begin{align*}
	\texttt{as }\delta_k \downarrow 0, \quad u_{\delta_k} \rightarrow u \texttt{ in } X, \quad
	F(u_{\delta_k}) \rightarrow u^* \texttt{ in } X^* 
\end{align*}
Furthermore, by step 1 we have
\begin{align*}
	\langle F(u_{\delta_k}), u_{\delta_k} - v \rangle \le 0 \quad \forall v \in M
\end{align*}
hence
\begin{align*}
	\langle u^*, u - v \rangle \leftarrow \langle F(u_{\delta_k}), u_{\delta_k} - v \rangle \le 0
\end{align*}
By the assumption \textbf{F3} we have
\begin{align*}
	\langle F(u), u - v\rangle \le \langle u^*, u - v \rangle \le 0
\end{align*}
\textbf{Step 3: }
$\forall R \in \mathbb{N}$, define $M_R:= \{ v \in X: \|v\|_X \le R\}$, by step 2 we have
\begin{align*}
	\forall R  \exists u_R \in M_R: \langle F(u_R), u_R - v \rangle \le 0 \qquad \forall v \in M_R
\end{align*}
and by the coercivity we know for all $u_0$ there exists some $C_0$, such that
\begin{align*}
	\langle F(u), u - u_0 \rangle > 0  \quad \forall \|u \| \ge C_0
\end{align*}
choose any large $R \ge \max\{C_0, \|u_0\|\}$, in particular $u_0 \in M_R$, we know by the coercivity the $u_R$ we get by step 2 must satisfies
\begin{align*}
	\|u_R\| \le C_0
\end{align*}
$\forall v \in M$: by Minty's trick contracting $v \in M$ and make  $v_\epsilon = (1 - \epsilon) u + \epsilon v \in M \cap B_\epsilon(u_R) $, in particular $v_\epsilon \in M_R$ for $\epsilon$ small enough, we have
\begin{align*}
	0 \ge \langle F(u_R), u_R - v_\epsilon \rangle = \epsilon \langle F(u_R), u_R - v \rangle \quad \forall v \in M
\end{align*} 
\end{proof}

\begin{theorem}
	Given $X$ a Hilbert space and $F: M \rightarrow X^*$ satisfies \textbf{F1-F3}. The variational inequality admits a solution. 
\end{theorem}
\begin{proof}
	By Galerkin approximation. The setup: a nested sequence of subspace $X_n \subset X_{n + 1} \subset X$. $u_0 \in X_1$, $\bigcup_n X_n = X, M_n = M \cap X_n$. The construction is possible since $X$ is separable. \\
	Define:
	\begin{align*}
		F_n: M_n \rightarrow X_n^*, \quad \langle F_n(u), v \rangle = \langle F(u), v \rangle \qquad \forall v \in X_n
	\end{align*}
	\textbf{1):}Verify that $F_n$ satisfies the assumption \textbf{F1-3}.
	\begin{enumerate}
		\item $\|F_n\|_{X^*} \le \|F\|_{X^*}$, hence $F_n$ must also map the bounded set to bounded set. 
		\item $u_0 \in X_1 \Rightarrow u_0 \in M_n$.
		\item Suppose $u_n \rightharpoonup u, F_n(u_n) \rightharpoonup u^*$, since $\dim X < \infty$ the weak convergence implies strong convergence, we have
		\begin{align*}
			\langle F_n(u_n), u_n \rangle \rightarrow \langle u^*, u \rangle  
		\end{align*}
		For $\langle F_n(u) - u^*, u - v \rangle \forall v \in M_n$: Since $u_n \in M_n$ weakly(and strong) convergent in $X_n$  therefore  
		\begin{align*}
			\langle F_n(u_m), u_m) \rangle = \langle F(u_m), u_m) \rangle \rightarrow \langle \tilde{u}^*, u\rangle
		\end{align*}
		Since $\langle F_n(u_m), v \rangle = \langle F(u_m), v \rangle \le \langle u^*, v \rangle$
	\end{enumerate}
	\textbf{2): The weak compactness}
	By the coercivity condition \textbf{F2} we have
	\begin{align*}
		\|u_n\| \le C_0 \quad \forall n \in \mathbb{N}
	\end{align*}
	By \textbf{F1} we know $F$ maps bounded set to bounded set, therefore
	\begin{align*}
		F(u_n) \le \infty \quad \textbf{ in } X
	\end{align*}  
	by reflexivity of $X$ we have (without relabeling)
	\begin{align*}
		F(u_n) \rightharpoonup u^*
	\end{align*}
	\textbf{3): Passage to the limit}
	we have
	\begin{align*}
		\langle F(u_n) , u_n - v \rangle \le 0  \qquad \forall v \in M_n
	\end{align*}
	It follows
	\begin{align*}
		&\langle F(u_n), u_n \rangle \le \langle F(u_n), v \rangle \rightarrow \langle u^*, v \rangle \texttt{ in } M_\infty\\
		(\textbf{F3})
		& \Rightarrow \langle u^*, u \rangle \le   \liminf_n \langle F(u_n), u_n \rangle \le \langle u^*, v \rangle\\
		&\Rightarrow \langle u^*, u - v \rangle \le 0
	\end{align*}
	Choose $v = u$ we have
	\begin{align*}
		\langle u^*, u \rangle \le  \liminf \langle F(u_n), u_n \le \langle u^*, u \rangle 
	\end{align*}
	We have
	\begin{align*}
		\liminf_n \langle F(u_n), u_n \rangle = \langle u^*, u \rangle
	\end{align*}
	therefore by \textbf{F3}
	\begin{align*}
		\langle F(u) - u^*, u - v \rangle \le 0 \quad \forall v \in M
	\end{align*}
	therefore
	\begin{align*}
		\langle F(u), u - v \rangle \le \langle u^*, u - v \rangle \le 0
	\end{align*}
\end{proof}

\begin{example}[Elastic Beam]
	Observe
	\begin{align*}
		I(v) := \int_I \frac{a_0}{2} |v''| + f \cdot \nu  dx 
	\end{align*}
	with 
	\begin{align*}
		M:= \{ v \in H^2(I) \cap H^1_0(I) | \quad v \ge g \quad a.e. \quad g(0), g(l) \le 0 \}
	\end{align*}
\end{example}


\subsection{The p-Laplacian}
Consider for $\Omega \subset \mathbb{R}^d$
\begin{align*}
	\inf_{v \in W^{1, p}_0}I(v):= \int_\Omega \frac{1}{p}|\nabla v |^p + \frac{\mu v }{2} + fv dx 
\end{align*}
we arrive at the ELE:
\begin{align*}
	-\div (|\nabla u |^p \nabla u | + \mu u = f &\texttt{ in } \Omega\\
	u = 0 &\texttt{ on } \partial\Omega
\end{align*}



\subsection{Perturbed monotone operators}
In the perturbed case we have the operator consisting of 2 parts, usually one described by the non-linearity and one by linearity(or continuity). We consider the non-linear part the principle term and the linear part the perturbed term, and extend our previous result regarding the non-linear elliptic variational inequality. 
\begin{theorem}
	Given $A: M \times M \rightarrow X^*$ with the general setting. (M closed, convex, X reflexive). Define $F(u):= A(u, u)$, if $A$ satisfies
	\begin{enumerate}
		\item $A$ is monotone in the second argument.
		\item $\forall v \in M$:  $u_n \rightharpoonup u \Rightarrow A(u_n, v)$ in M $\overset{*}{\rightharpoonup} A(u, v)$ in $X^*$.
		\item $\forall v \in M$: $u_n \rightharpoonup u$ $\Rightarrow \langle A(u, v), u \rangle \le \liminf_{n \rightarrow \infty} \langle A(u_n, v), u_n  \rangle$.   
	\end{enumerate}
\end{theorem}
\begin{proof}
\end{proof}


\subsection{Quasi-linear Elliptic systems}
Consider
\begin{align*}
	- \div (a(u, \nabla u) ) + f(u) = 0 & \texttt{ in } \Omega\\
	u = 0 & \texttt{ on } \partial \Omega 
\end{align*}

We assume:
\textbf{The Growth condition}
\begin{enumerate}
	\item $\forall \delta > 0 \exists C_\delta$:
	\begin{align*}
		a(z, q) \le C_\delta (|q|^{p -1} + 1) + \delta |z|^{p -1}
	\end{align*}
	\item There exists $C > 0$ s.t.
	\begin{align*}
		f(z) \le C |z|^{p - 1}
	\end{align*}
\end{enumerate}




\section{Classical Maximum Principle}
\subsection{Weak Maximum Principle}

\subsection{Harnack's Inequality}
\begin{theorem}
	Given $u$ a solution( i.e., $\partial_t u + Lu = 0$) with $u \ge 0$ in $\Omega_T$ and $u \in C^\infty(\Omega_T)$. Assuming $a_{ij} = a_{ji} \in C^\infty(\Omega_T), b_i = 0, c_i = 0$, $\Omega' \subset \subset \Omega$ subdomain,  then there exists $\forall t_1, t_2$ with  $0 < t_1 < t_2 \le T$ a constant $C$ such that
	\begin{align*}
		\sup_{\Omega'} u(t_1, \cdot) \le C \inf_{\Omega'} u(t_2, \cdot)
	\end{align*}
\end{theorem}

\subsection{Strong Maximum Principle}
\begin{theorem}[Strong maximum principle]
The coefficients of $L$ defined as before. 
Given $u$ a sub-solution(i.e., $\partial_t u + L u \le 0$) and $u \in C^\infty(\Omega_T) \cap C(\overline{\Omega_T})$. If $u$ admits its maximum at $(t_0, x_0) \in \Omega_T$, then 
\begin{align*}
	u \equiv const \texttt{ in } \Omega_{t_0}:= (0, t_0] \times \Omega
\end{align*}
\end{theorem}
\begin{proof}
Define $\Omega' \subset\subset$ with $(t_0, x_0) \in \Omega''$.
	Define $w$: 
	\begin{align*}
		&\partial_t w + Lw = -(\partial_t u + L u) \ge 0 \quad (t, x) \in \Omega_T'\\
		& w = 0 \quad \{0\} \times \Omega'
	\end{align*}
	The parabolic regularity asserts $w$ admits the same regularity. Then $v = w + u$ satisfies
	\begin{align*}
		& \partial_t v + Lv = 0 \quad (t, x) \in \Omega_T'\\
		& v = 0 \quad \{0\} \times \Omega'
	\end{align*}
	by the weak maximum principle we have
	\begin{align*}
		u \le v \le u \le u(t_0, x_0) := M
	\end{align*}
	Further define
	\begin{align*}
		\tilde{v}: = M - v
	\end{align*}
	we have
	\begin{align*}
		&\partial_t \tilde{v} + L\tilde{v} = 0 \quad \Omega' \\
		&\tilde{v} \ge 0 \quad \Omega'
	\end{align*}
	Apply Harnack's inequality we have
	\begin{align*}
		&\sup_{t \le t_0} \tilde{v} \le \tilde{v}(t_0, \cdot) = 0\\
		&\Rightarrow \tilde{v} \equiv 0 \quad \forall t \in (0, t_0)\\
		&\Rightarrow v \equiv const
	\end{align*}
	approximate $\Omega$ by $\Omega'$ we get the result. 
\end{proof}
 
 

\section{Reaction-Diffusion System}
\subsection{Global Existence of Lipschitz nonlinearity}
\begin{theorem}
Given the PDE
\begin{align*}
	\partial_t u - \Delta u = f(u) \quad &\texttt{ in } (0, T] \times \Omega\\
	u(0, ) = u_0 \quad &\texttt{ in } \{0\} \times \Omega\\
	u \equiv 0 \quad &\texttt{ 0n } (0, T] \times \partial\Omega
\end{align*}
where $f$ Lipschitz continuous with
\begin{align*}
	|f(u) - f(v)| \le C |u - v|
\end{align*}
Then there exists a unique solution 
\begin{align*}
	u \in C([0, T], L^2(\Omega))
\end{align*}
\end{theorem}
\begin{proof}
	proof with the contracting map. \\
	\textbf{Step 1:} Transform to a linear problem: define $h = f(u)$ and observe the system:
	\begin{align*}
		\partial_t w - \Delta w &= h\\
		w(0, ) &= u_0
	\end{align*}
	by the linear theorem there exists a \textbf{unique} solution 
	\begin{align*}
		w \in W(0, T; V, H) \hookrightarrow C([0, T], L^2)
	\end{align*}
	we can therefore define the map
	\begin{align*}
		\mathcal{A}: X_t \rightarrow X_t , \mathcal{A}(u) = w
	\end{align*}
	\textbf{Step2:} Show $\mathcal{A}$ is a contraction map
	observe the weak form:
	\begin{align*}
		\int_\Omega w_t \varphi + \mathbb{A} \nabla w \cdot \nabla \varphi dx = \int_\Omega f(u) \varphi dx \qquad \forall_{a.a.} t \in [0, T]
	\end{align*}
	Given $w, \tilde{w} \in \ran(\mathcal{A})$ and test with we have
	\begin{align*}
		\langle \partial_t(w - \tilde{w}), w - \tilde{w} \rangle 
		+ \|\nabla (w - \tilde{w})\|_{L^2} 
		&= \langle f(u) - f(\tilde{u}), w - \tilde{w} \rangle\\
		\frac{d}{dt}\|w - \tilde{w}\|_{L^2}^2 + C_p \|w - \tilde{w}\|_{L^2}^2 
		&\le C \|u - \tilde{u}\|_{L^\infty} \|w - \tilde{w}\|_V\\
	\end{align*}
	where the bracket indicates the linear mapping $\langle , \rangle_{V^*, V}$. By the Young's inequality the righthand side is bounded by 
	$C_\delta( \|u - \tilde{u}\|_{L^\infty}^2 + \delta \|w - \tilde{w}\|^2_{V}$, and again by Poincare the $V$ norm could be written as $L^2$ norm up to a constant. In summary we have after integration over $t$
	\begin{align*}
		\|(w - \tilde{w})(t)\|_H^2 \le t C \|u - \tilde{u}\|_{L^\infty}^2
	\end{align*}
\end{proof}





\subsection{The Heat kernel}



\subsection{Local well-posedness for locally Lipschitz nonlinearity}


\subsection{Maximum Life Span}


\subsection{Blow-up}



\section{Hyperbolic Equation}
We observe the problem:
\begin{equation}
	\partial_t u + \sum_{i = 1}^d A_j \partial_j u = B u + f
\end{equation}
where $u:[0, T) \times \mathbb{R}^d \rightarrow \mathbb{R}^n$ for $T \in (0, \infty]$ with $A_j(t, x), B(t, x) \in C_b^\infty$. We introduce the notation:
\begin{align*}
	&L = L(t) = \sum_{j = 1}^d A_j(t) \partial_j x - B(t)\\
	&Pu = \partial_t u + Lu
\end{align*}


\subsection{fractional Sobolev spaces}
Recall the definition of the Schwartz space:
\begin{definition}[Schwartz space]
\begin{equation}
	\mathcal{S}(\mathbb{R}^d):= \{ f \in C^\infty(\mathbb{R}^d, \mathbb{C}): |x^\alpha \partial^\beta u(x) \le C_{\alpha, \beta}, \texttt{for all } x \in \mathbb{R}^d\}	
\end{equation}
\end{definition}

In the Fourier transform we know the differentiation becomes multiplication of the monoid, i.e.


\subsection{Constant coefficient system}
Observe the problem
\begin{align*}
	& \partial_t u + \sum_{ j = 1}^d A_j \partial_j u = 0 \texttt{ in } (0, T) \times \mathbb{R}^d\\
	& u(0, \cdot) = u_0 \texttt{ in } \mathbb{R}^d
\end{align*}


\subsection{Linear symmetric hyperbolic system}
Observe the system
\begin{align*}
	& \partial_t u + \sum_{j = 1}^d A_j \partial_j u = Bu + f(t, x) \texttt{ in } (0, T)\\
	& u(0, \cdot) = u_0 \texttt{ in } \mathbb{R}^d
\end{align*}
The coefficients are assumed to be bounded and smooth with bounded derivatives. We now define 
\begin{align*}
	& L(t) = \sum_{j = 1}^d A_j \partial_j u - B\\
	& P(u) = Lu + \partial_t u
\end{align*}

\begin{definition}[The hyperbolic system]
\end{definition}

\begin{definition}[The symmetric hyperbolic system]
\end{definition}



\subsection{The energy estimate}
\begin{definition}[The distributional solution]
\end{definition}
why do we need the concept of distributional solution. 

\begin{proposition}[The energy estimate($L^2$ energy estimate)]
	Given $T > 0$ and $S_0(x) \in C_b^\infty(\mathbb{R}^d, \mathbb{R}^{n \times n})$. with $0 < \lambda \le \Lambda$ s.t. 
	\begin{align*}
		\lambda \le S_0(x) \le \Lambda \quad \forall x \in \mathbb{R}^d,
	\end{align*}
	then we have the  
\end{proposition}
and the equation may be rewritten as 
\begin{equation}
	Pu = f
\end{equation}

\begin{proposition}[The general energy estimate]
\end{proposition}







\end{document}